% =====================================================================
%  Characterization and Convergence of Markov Processes
%  via Martingale Problems
%
%  A manuscript prepared as a blueprint for formalization in Lean 4.
%
%  Compile with:  pdflatex MartingaleProblem.tex  (twice)
% =====================================================================

\documentclass[11pt,a4paper]{article}

\usepackage[T1]{fontenc}
\usepackage[utf8]{inputenc}
\usepackage{amsmath, amssymb, amsthm, mathrsfs}
\usepackage{mathtools}
\usepackage[margin=3cm]{geometry}
\usepackage{enumitem}
\usepackage{booktabs}
\usepackage[colorlinks=true, linkcolor=blue, citecolor=blue, urlcolor=blue]{hyperref}

% ------------------------- theorem environments ----------------------
\theoremstyle{plain}
\newtheorem{theorem}{Theorem}[section]
\newtheorem{lemma}[theorem]{Lemma}
\newtheorem{proposition}[theorem]{Proposition}
\newtheorem{corollary}[theorem]{Corollary}

\theoremstyle{definition}
\newtheorem{definition}[theorem]{Definition}
\newtheorem{setting}[theorem]{Setting}
\newtheorem{assumption}[theorem]{Assumption}
\newtheorem{example}[theorem]{Example}

\theoremstyle{remark}
\newtheorem{remark}[theorem]{Remark}
\newtheorem{fact}[theorem]{Fact}

% ------------------------- macros ------------------------------------
\newcommand{\R}{\mathbb{R}}
\newcommand{\N}{\mathbb{N}}
\newcommand{\Q}{\mathbb{Q}}
\newcommand{\Rp}{\R_+}
\newcommand{\T}{\mathbb{T}}                     % abstract time index
\newcommand{\Bor}{\mathcal{B}}                 % Borel sigma-field
\newcommand{\Meas}{M}                          % measurable functions
\newcommand{\Bdd}{B}                           % bounded measurable functions
\newcommand{\Cb}{C_b}                          % bounded continuous functions
\newcommand{\Prob}{\mathcal{P}}                % probability measures
\newcommand{\DE}{D_E[0,\infty)}
\newcommand{\CE}{C_E[0,\infty)}
\newcommand{\dom}{\mathcal{D}}                 % domain of an operator
\newcommand{\ran}{\mathcal{R}}                 % range of an operator
\newcommand{\Filt}{\mathcal{F}}
\newcommand{\Gilt}{\mathcal{G}}
\newcommand{\XX}{\mathbb{X}}                   % family of test processes
\newcommand{\ZZ}{\mathbb{Z}}                   % determining set
\newcommand{\MP}{\mathrm{MP}}
\newcommand{\Msol}{\mathcal{M}}
\newcommand{\one}{\mathbf{1}}
\newcommand{\dif}{\,\mathrm{d}}
\newcommand{\Chat}{\hat{C}}                    % continuous functions vanishing at infinity
\newcommand{\Vsig}{V}                          % finite signed measures
\newcommand{\bplim}{\operatorname{bp-lim}}
\newcommand{\Lsp}{L}                           % a Banach space
\newcommand{\wstar}{\mathrm{w}^{*}}
\DeclareMathOperator*{\wto}{\Rightarrow}

\newcommand{\EK}{\cite{EK86}}
\newcommand{\CPS}{\cite{CPS23}}
\newcommand{\KA}{\cite{Kal21}}
\newcommand{\JS}{\cite{JS03}}
\newcommand{\BM}{\cite{BM25}}

\title{Characterization and Convergence of Markov Processes\\
       via Martingale Problems\\[2ex]
       \large A blueprint for formalization}
\author{}
\date{\today}

\begin{document}
\maketitle

\begin{abstract}
This manuscript collects, in a self-contained and formalization-oriented form,
four results on martingale problems:
(i)~existence of a c\`adl\`ag modification of a solution
(Ethier--Kurtz \EK, Theorem~4.3.6);
(ii)~the fact that uniqueness of the one-dimensional distributions already implies
uniqueness of the finite-dimensional distributions and the (strong) Markov property
(\EK, Theorem~4.4.2);
(iii)~duality as a tool for proving uniqueness
(\EK, Lemma~4.4.10, Theorem~4.4.11 and its corollaries);
(iv)~existence: four routes, of which two produce a solution out of nothing ---
an explicit construction of jump processes, and the correspondence with
stochastic differential equations --- while the fourth obtains solutions as weak
limits of solutions of approximating martingale problems (\EK, Lemma~4.5.1 and
Remark~4.5.2), together with the strictly more general convergence theorem of
Criens--Pfaffelhuber--Schmidt \CPS.
Each of (i)--(iii) is first proved in the abstract setting of \CPS, without any
operator, and \EK's statement is then recovered as a corollary; the abstract
versions are new, since \CPS{} prove no characterization results.  Throughout,
the time index is an abstract ordered set carrying a \emph{clock}, the state
space is graded from ``measurable space'' to ``Polish'', and local martingale
problems are treated wherever the statement permits it.  Every statement is
tagged with the weakest hypotheses on the index and on the state space that it
uses, and Section~\ref{ssec:bundletable} tabulates the result: the partially
ordered index survives in (ii) for the Markov property but provably not for
uniqueness, and in (iii) not at all; a merely measurable state space suffices for
(ii) and (iii) throughout.
Section~\ref{sec:setting} explains, and justifies, the choice of the underlying
setting; Section~\ref{sec:formalization} records the design decisions that are
relevant for the Lean formalization.
\end{abstract}

\tableofcontents

% =====================================================================
\section{Scope, and the choice of the setting}
\label{sec:setting}
% =====================================================================

\subsection{The two candidate settings}

There are two natural frameworks in which the results of this manuscript can be
phrased.

\paragraph{(EK) The Ethier--Kurtz setting \EK, Chapter~4.}
A martingale problem is attached to an \emph{operator}
$A \subset \Meas(E) \times \Meas(E)$ on a metric space $E$: a measurable
$E$-valued process $X$ solves the martingale problem for $A$ if
\begin{equation}
  \label{eq:EKtest}
  f(X(t)) - \int_0^t g(X(s)) \dif s , \qquad t \geq 0,
\end{equation}
is a martingale for every $(f,g) \in A$.  This is the setting in which
\emph{characterization} results are traditionally stated: uniqueness, the Markov
property, duality, and the correspondence with semigroups and generators.

\paragraph{(A) The abstract setting.}
A martingale problem is attached to a \emph{family $\XX$ of real-valued test
processes} on a filtered space carrying, as yet, no measure: a probability
measure $P$ solves the martingale problem $\MP(\XX)$ if every $Y \in \XX$ is a
$P$-martingale, and the local martingale problem if every $Y \in \XX$ is a
$P$-local martingale.  The state space is Polish and the path space is an
arbitrary Polish subset of $E^{\T}$; no operator and no Markov structure are
assumed.

It is worth being accurate about where this second setting comes from, since
this manuscript leans on it heavily.  It is not new with \CPS.  In essentially
the form used here it is \JS, Definition~III.1.3, which attaches a martingale
problem to a family of optional processes and does so with \emph{local}
martingales as the default rather than as a variant, and with the initial
condition given on an initial $\sigma$-field rather than by an initial
distribution.  \CPS{} rediscover it in the shape that suits weak convergence ---
canonical test processes, determining sets, $P$-continuity --- and it is their
formulation that Section~\ref{sec:MP} follows, but Remark~\ref{rem:absMPprov}
records what \JS{} already had and what is worth importing from them.
\KA, Chapter~32 works in the same spirit for diffusions, again with local
martingales throughout.

\paragraph{What each setting is good for.}
The division of labour is not the one the first draft of this manuscript
assumed, and the correction is the main structural finding here.  Setting (A) is
where \emph{all four} target results are naturally proved --- including the three
that \CPS{} do not state.  Setting (EK) is where they are naturally
\emph{applied}: in practice one is handed an operator, not a family of test
processes, and the hypotheses of Section~\ref{ssec:uniqmarkov} are the ones that
can be checked on $A$ directly.  Section~\ref{ssec:recommendation} sets out the
resulting three-layer architecture.

\subsection{Recommendation}
\label{ssec:recommendation}

We use a \emph{hybrid}, organised in three layers.  The reader should be warned
that the balance between them shifted while this manuscript was written, and
that the outcome is not the one the first draft anticipated.

\begin{enumerate}[label=(\arabic*)]
\item \textbf{Definitional layer: (A).}
  ``Solution of a martingale problem'' is defined as in \CPS, Definition~3.2,
  relative to a family $\XX$ of test processes, with a \emph{local} variant
  available from the outset (Definition~\ref{def:absMP}).  This costs one line
  and is what makes local martingale problems and non-c\`adl\`ag path spaces
  expressible at all.

\item \textbf{Abstract working layer: also (A).}
  Each of the three characterization results is stated and proved in this
  setting, with no operator anywhere: Theorem~\ref{thm:absreg} for the c\`adl\`ag
  modification, Theorems~\ref{thm:absuniq} and \ref{thm:absstrongmarkov} for
  uniqueness and the (strong) Markov property, and Lemma~\ref{lem:chain} with
  Proposition~\ref{prop:haar} for duality.

\item \textbf{Markovian layer: (EK).}
  The martingale problem for an operator $A$ is the special case
  $\XX = \XX_A(X)$ of \eqref{eq:EKtest} (Definition~\ref{def:markovMP}), and the
  four target results are recovered on it: Theorem~\ref{thm:cadlag},
  Theorem~\ref{thm:uniqueness}, Theorem~\ref{thm:duality} and
  Lemma~\ref{lem:EKconv}.  Each is a corollary of layer~(2), and each proof is a
  verification of hypotheses rather than an argument.
\end{enumerate}

Layer~(2) is the substance of this manuscript, and it is the point on which the
first draft was wrong.  It assumed that uniqueness, the Markov property and
duality are \emph{intrinsically} statements about an operator $A$ and its domain,
and that (A) could therefore serve only as packaging.  That is not so.  Each of
the three proofs, read carefully, uses of the operator only that it produces
test processes of a particular shape, and each hypothesis it really needs can be
stated directly:

\begin{itemize}
\item the c\`adl\`ag theorem needs a class $\Phi$ for which $f(X)$ splits into a
  martingale plus a time-regular remainder --- Definition~\ref{def:regclass};
\item uniqueness and the Markov property need the test family to be stable under
  the shift --- Definition~\ref{def:shiftstable} --- and nothing else, the whole
  proof reducing to one lemma applied four times (Lemma~\ref{lem:restart});
\item duality needs only two increment representations of a function of two time
  variables (Lemma~\ref{lem:chain}), and no probability at all.
\end{itemize}

\noindent
Remarks~\ref{rem:absreggain}, \ref{rem:absuniqgain} and \ref{rem:dualscope}
record what each abstraction buys.  Briefly: path-dependent and atomic
compensators, which cannot be written with an operator
$A \subset \Cb(E) \times \Bdd(E)$ at all; the disappearance of the auxiliary
filtration \eqref{eq:starfilt} from everything except
Proposition~\ref{prop:fddchar}; and a precise account of which time indices and
which clocks each result admits (Section~\ref{ssec:bundletable}).

Two caveats keep this from being an argument for (A) alone.

\begin{itemize}
\item \emph{\CPS{} contains none of the three.}  That paper is concerned
  exclusively with identifying weak limits; it proves no uniqueness statement, no
  Markov property and no path regularity statement.  For the c\`adl\`ag theorem
  the omission is structural rather than accidental: their Setting~\ref{set:abstract}
  \emph{presupposes} a Polish path space $F$ and right continuous test processes,
  which is exactly what a regularization theorem has to produce.  Layer~(2) is
  therefore new here, not imported.
\item \emph{Layer~(3) is not decoration.}  The abstract hypotheses are the ones a
  proof needs; the operator is the one a user has.  In applications one is handed
  $A$, and the value of Theorem~\ref{thm:uniqueness} over
  Theorem~\ref{thm:absuniq} is that its hypotheses are checkable on $A$ without
  first reconstructing $\XX_A$, $\theta_r$ and a determining set.
\end{itemize}

For the fourth target result --- existence by convergence, \EK, Lemma~4.5.1 ---
the \CPS{} formulation is a strict generalization (Theorem~\ref{thm:CPSconv}): it
drops boundedness of $f$ and $g$, replaces exact solutions of approximating
martingale problems by \emph{approximate} martingales, and replaces uniform
boundedness by uniform integrability.  Here \CPS{} is used directly, and
Theorem~\ref{thm:absconv} isolates the abstract statement behind it, which needs
neither c\`adl\`ag paths nor the Skorokhod topology and is proved here in full.

\subsection{Generality: time index, Polish state space, local martingale problems}

Four deliberate strengthenings relative to \EK{} are made.

\paragraph{Abstract time index.}
\EK{} and \CPS{} index time by $[0,\infty)$.  Here the index is an abstract set
$\T$, and every statement is tagged with the weakest of the bundles
(T0)--(T4) of Definition~\ref{def:bundles} that it needs; the compensator
becomes a separate datum, a \emph{clock} in the sense of
Definition~\ref{def:clock}.  Section~\ref{ssec:bundletable} tabulates the result.
There are two reasons for the effort.  The first is that the classical
formulation conflates a topological hypothesis on $\T$ with an algebraic one:
Section~\ref{sec:cadlag} needs the former, Sections~\ref{sec:uniqueness} and
\ref{sec:duality} the latter, and $[0,\infty)$ happens to supply both.  The
second is that Mathlib indexes \texttt{Filtration}, \texttt{Martingale} and
\texttt{IsStoppingTime} by an arbitrary \texttt{[Preorder \(\iota\)]}, so fixing
$\T = \Rp$ at the outset would mean working against the library rather than with
it.

The gain is not merely cosmetic.  Discrete time ($\T = \N_0$, $q$ counting
measure) and processes with fixed times of discontinuity ($q$ with atoms) become
special cases rather than analogues, and the partially ordered case survives in
Section~\ref{sec:MP}, in the Markov half of Section~\ref{sec:uniqueness} and in
the abstract convergence theorem.  It fails in exactly two places, and for
reasons worth knowing: the uniqueness induction needs the times to form a chain
(Remark~\ref{rem:chainonly}), and duality needs order intervals to be carried
onto order intervals by reflection (Remark~\ref{rem:dualscope}).

\paragraph{State space: measurable, not Polish, where possible.}
\EK{} and \CPS{} take $E$ Polish throughout.  Here $E$ is graded exactly as the
time index is: (E0) a measurable space, (E1) standard Borel, (E2) separable
metrizable, (E3) Polish (Definition~\ref{def:Ebundles}).  The grading is not
decoration.  Sections~\ref{sec:MP} and \ref{sec:duality} and almost all of
Section~\ref{sec:uniqueness} run under (E0); regular conditional distributions,
and hence (E1), are needed at exactly three places; the topology of $E$ is used
only in Sections~\ref{sec:cadlag} and \ref{sec:convergence}.  The payoff is
concrete: martingale problems for distribution-valued processes, where
$E = \mathcal{S}'(\R^d)$ is standard Borel but not metrizable, are covered by the
first three sections with no new proof.  See Remarks~\ref{rem:Ebundleuse} and
\ref{rem:E1why}.

Where a topology is used, $E$ is \emph{not} assumed locally compact.  On a
general Polish space there is no space $C_c(E)$ and no one-point
compactification argument, and the r\^ole of local compactness is taken over
throughout by the \emph{compact containment condition} \eqref{eq:cc} resp.\
\eqref{eq:ccn}.  This is the essential mechanism, and it is worth isolating it as
a definition (Definition~\ref{def:cc}) in the formalization.

\paragraph{Local martingale problems.}
Whether ``martingale'' may be weakened to ``local martingale'' differs from
result to result, and it is important not to overclaim.  The situation is:

\begin{itemize}
\item \emph{Theorem~\ref{thm:cadlag} (c\`adl\`ag modification):} must be stated
  for genuine martingales.  Localization requires stopping times, stopping times
  require a right-continuous filtration and a progressively measurable process,
  and path regularity is precisely what is being proved.  The argument is
  circular otherwise.  See Remark~\ref{rem:cadlaglocal}.
\item \emph{Theorem~\ref{thm:uniqueness} (uniqueness and Markov property):} the
  proof conditions on events of positive probability and uses the martingale
  property under the conditioned measures; boundedness of $f$ and $g$ enters
  through the integrability needed for \eqref{eq:fdd}.  A local version is
  available, but only after localization along stopping times that are
  \emph{common} to all solutions; see Remark~\ref{rem:uniquelocal}.
\item \emph{Theorem~\ref{thm:duality} (duality):} is already stated for
  unbounded $f,g,h \in \Meas(E_1 \times E_2)$, controlled by the integrability
  hypotheses \eqref{eq:dual1}--\eqref{eq:dual2}.  Nothing has to be changed.
\item \emph{Section~\ref{ssec:jump} (jump processes):} the local version is not
  an afterthought but the general case --- with unbounded rates the process
  explodes, and the jump times are a localizing sequence by construction; see
  Remark~\ref{rem:jumpexplosion}.
\item \emph{Theorem~\ref{thm:CPSconv} (convergence):} concludes that the limit
  solves the \emph{global} martingale problem.  A local martingale problem is
  obtained in the limit by applying the theorem to a localized family of test
  processes; see Remark~\ref{rem:convlocal}.
\end{itemize}

\noindent
Section~\ref{ssec:localmp} carries out the local theory in the abstract setting,
following \KA, Theorems~32.10 and 32.11.  Two hypotheses are needed and are
isolated there: a localizing sequence that is a path functional common to all
$P$, and covariance of that sequence under the shift.  The point of isolating
them is that the local martingale property is \emph{not} linear in $P$, so that
the two workhorse lemmas of Section~\ref{ssec:absuniq} --- mixture stability and
the restart lemma --- do not survive localization for free.

% =====================================================================
\section{Preliminaries}
\label{sec:notation}
% =====================================================================

Following the project plan, all material of \EK, Chapters~1--3, may be cited.
This section makes that interface explicit: it collects, in the form of
numbered \emph{Facts}, precisely those definitions and results of
\EK, Chapters~1--3 (together with a few items from the Appendixes of \EK) that
are used later on.  Nothing here is proved; the r\^ole of the section is to fix
statements and to serve as the list of prerequisites for the formalization.
Section~\ref{ssec:usage} records where each item is used.

\subsection{Notation}
\label{ssec:notation}

Let $(E,\mathcal{E})$ be a measurable space; the topological assumptions on $E$
are made bundle by bundle in Definition~\ref{def:Ebundles}, and $\Bor(E)$ is
written for $\mathcal{E}$ whenever $E$ carries a topology.  We write
\begin{itemize}[nosep]
\item $\Meas(E)$ for the real-valued Borel measurable functions on $E$;
\item $\Bdd(E) \subset \Meas(E)$ for the bounded ones, with
  $\lVert f\rVert = \sup_{x \in E} |f(x)|$;
\item $C(E)$ for the continuous functions, $\Cb(E) = C(E) \cap \Bdd(E)$, and,
  for $E$ locally compact, $\Chat(E)$ for the continuous functions vanishing at
  infinity;
\item $\Prob(E)$ for the Borel probability measures on $E$ and $\Vsig(E)$ for the
  finite signed Borel measures on $E$;
\item $\DE$ for the space of c\`adl\`ag functions $[0,\infty) \to E$ and $\CE$
  for the continuous ones.
\end{itemize}
$\Lsp$ always denotes a real Banach space with norm $\lVert \cdot \rVert$.

An \emph{operator} is a subset $A \subset \Meas(E) \times \Meas(E)$; it need not
be single-valued nor linear.  Its domain is
$\dom(A) = \{f : (f,g) \in A \text{ for some } g\}$, its range is
$\ran(A) = \{g : (f,g) \in A \text{ for some } f\}$, and $A_S$ denotes its linear
span inside $\Meas(E) \times \Meas(E)$.  For $\lambda \in \R$ we write
$\lambda - A = \{ (f, \lambda f - g) : (f,g) \in A \}$.

Throughout, references of the form ``\EK, Theorem~3.9.1'' mean Chapter~3,
Section~9, item~1 of \EK.

% ---------------------------------------------------------------------
\subsection{The time index and the clock}
\label{ssec:timeindex}
% ---------------------------------------------------------------------

\EK{} and \CPS{} both index time by $\Rp$.  Almost nothing in this manuscript
needs that, and the parts that do need it need it for visibly different reasons:
Section~\ref{sec:cadlag} needs a \emph{topological} property of the index,
Sections~\ref{sec:uniqueness} and \ref{sec:duality} an \emph{algebraic} one.  We
therefore fix an abstract index $\T$ and record, for each statement, the weakest
bundle of assumptions it uses.  Two motives: it separates the topological from
the algebraic hypotheses, which the classical formulation conflates; and it
matches Mathlib, which indexes \texttt{Filtration}, \texttt{Martingale} and
\texttt{IsStoppingTime} by an arbitrary \texttt{[Preorder \(\iota\)]}, so that
committing to $\Rp$ at the outset would mean working against the library.

\subsubsection*{Structure on the index}

\begin{definition}
  \label{def:bundles}
  We use the following bundles of assumptions on $\T$.
  \begin{itemize}[nosep,leftmargin=3.2em]
  \item[\textbf{(T0)}] $\T$ is a preordered set with a least element $0$.
  \item[\textbf{(T1)}] (T0), and $\T$ is a directed lattice: $s \wedge t$ and
    $s \vee t$ exist.
  \item[\textbf{(T1$'$)}] (T1) in the concrete form of \EK, Section~2.8: a metric
    lattice whose order intervals are separable from above.
  \item[\textbf{(T2a)}] (T0) and $\T$ is linearly ordered.
  \item[\textbf{(T2b)}] (T2a), $\T$ carries the order topology, admits a
    countable dense subset $D$, and has no isolated points from the right.
  \item[\textbf{(T3)}] $\T = [0,\infty)$ or $\T = [0,T]$.
  \item[\textbf{(T4)}] (T0) and $\T$ is a cancellative ordered commutative
    monoid: $r + \cdot$ is order preserving, and $t - s$ exists for $s \leq t$.
  \end{itemize}
\end{definition}

Three remarks on the design of Definition~\ref{def:bundles}.

\emph{(T4) is independent of (T2a)--(T3).}  The monoid structure is what the
shift $X(r+\cdot)$ needs; it presupposes no linear order.  Keeping it separate
makes visible that Sections~\ref{sec:uniqueness} and \ref{sec:duality} rest on an
algebraic property of $\T$ whereas Sections~\ref{sec:cadlag} and
\ref{sec:convergence} rest on a topological one.

\emph{Why (T2a) and (T2b) are separate.}  The uniqueness half of
Theorem~\ref{thm:absuniq} and the whole of Section~\ref{sec:duality} need the
order to be linear but use no topology at all; only Section~\ref{sec:cadlag} and
the concrete convergence results need the topology.  Bundling the two, as one is
tempted to do, hides that.

\emph{(T1$'$) is not used in this manuscript.}  It is listed because it is the
only place where \EK{} themselves treat a general index set, and because it is
easy to misremember what it is for; Remark~\ref{rem:ekindex} sets that out.  But
every statement here that involves a stopping time --- the strong Markov property
of Theorem~\ref{thm:absstrongmarkov}, and the localized statements of
Remarks~\ref{rem:uniquelocal}, \ref{rem:cadlaglocal} and \ref{rem:convlocal} ---
already assumes (T2b), hence a linear order, for which the lattice structure is
automatic.  (T1$'$) would become relevant only for a strong Markov property on a
genuinely directed index, and Remark~\ref{rem:chainonly} shows that such a
statement is blocked earlier, on the uniqueness side.

\subsubsection*{The clock}

The index set alone does not carry a martingale problem.  The test processes of
Definition~\ref{def:markovMP} have the form $Y_t = f(X_t) - C^{f,g}_t$, and what
distinguishes the compensator of a \emph{generator} from an arbitrary adapted
process is that it is additive in time and generated pointwise from $g$.  That
is what carries Proposition~\ref{prop:fddchar}, and without a
finite-dimensional characterization the phrase ``martingale problem for the
operator $A$'' has no content, since every process solves the martingale problem
for \emph{some} family of test processes, namely the family of its own
martingales.

\begin{definition}[Clock]
  \label{def:clock}
  A \emph{clock} on $\T$ is a $\sigma$-field $\mathcal{T}$ on $\T$ for which the
  down-sets
  \[
    \T_{\leq t} = \{u \in \T : u \leq t\} ,
    \qquad
    \T_{<t} = \T_{\leq t} \setminus \{ u \in \T : t \leq u \}
  \]
  are measurable for every $t$, together with a measure $q$ on
  $(\T,\mathcal{T})$ such that $q(\T_{\leq t}) < \infty$ for every $t$.  For
  $s \leq t$ we write
  \begin{equation}
    \label{eq:clockinterval}
    (s,t] \;\coloneqq\; \T_{\leq t} \setminus \T_{\leq s} ,
    \qquad
    [s,t) \;\coloneqq\; \T_{<t} \setminus \T_{<s} ,
  \end{equation}
  and call these the \emph{optional} and the \emph{predictable} interval.  Both
  $t \mapsto \T_{\leq t}$ and $t \mapsto \T_{<t}$ are monotone, so both intervals
  are additive:
  \begin{equation}
    \label{eq:clockadd}
    (s,u] = (s,t] \uplus (t,u] , \qquad
    [s,u) = [s,t) \uplus [t,u) , \qquad (s \leq t \leq u) .
  \end{equation}
  We write $\iota \in \{\mathrm{o},\mathrm{p}\}$ for the choice of convention and
  $\langle s,t \rangle_\iota$ for $(s,t]$ resp.\ $[s,t)$.

  The clock is \emph{shift invariant} if $\T$ satisfies (T4) and
  $q(r + B) = q(B)$ for all $r \in \T$ and $B \in \mathcal{T}$; it is
  \emph{reflective} if in addition $\T$ satisfies (T2a); and it is
  \emph{atomless} if $q(\{u : t \leq u \leq t\}) = 0$ for every $t$, which is
  exactly the condition under which the two conventions coincide.
\end{definition}

The additivity in \eqref{eq:clockinterval} is not an assumption: it follows from
transitivity of $\leq$, which is why a preorder suffices.  This is the point at
which a measure is genuinely needed rather than an additive interval function
$q_{s,u} = q_{s,t} + q_{t,u}$: the proof of
Proposition~\ref{prop:fddchar} uses Fubini in the form
\[
  E\Bigl[ \int_{(s,t]} g(X_u)\, q(\dif u) \cdot Z \Bigr]
    \;=\; \int_{(s,t]} E\bigl[ g(X_u) Z \bigr]\, q(\dif u) ,
\]
and additivity along chains does not give that.  On $\T = \Rp^d$, for instance,
chain additivity is strictly weaker than $d$-monotonicity, and only the latter
produces a measure on boxes.

\begin{example}
  \label{ex:clocks}
  \begin{enumerate}[label=(\roman*),nosep]
  \item $\T = \Rp$, $q$ Lebesgue measure: $\T_{\leq t} = [0,t]$ and
    $C_t = \int_0^t g(X_u)\dif u$.  This is \EK{} and \CPS{} throughout.
  \item $\T = \N_0$, $q$ counting measure: $C_n = \sum_{0 < k \leq n} g(X_k)$,
    discrete time.
  \item $\T = \Rp$, $q$ with atoms: processes with fixed times of
    discontinuity, the subject of \CPS, Section~5.3.
  \item $\T = \Rp^d$, $q$ Lebesgue measure: $\T_{\leq t}$ is the box $[0,t]$ and
    $C_t = \int_{[0,t]} g(X_u)\dif u$, the multiparameter martingale problem.
  \end{enumerate}
  All four are clocks in the sense of Definition~\ref{def:clock}; (i), (ii) and
  (iv) are shift invariant, and only (i) and (ii) are reflective.
\end{example}

\begin{remark}[Optional or predictable: both]
  \label{rem:convention}
  The compensator of Definition~\ref{def:markovMP} is
  $C_t = \int_{\langle 0,t\rangle_\iota} g(X_u)\,q(\dif u)$, and the two choices of
  $\iota$ genuinely differ.  For $\iota = \mathrm{o}$ the compensator is right
  continuous; for $\iota = \mathrm{p}$ it is predictable and left continuous.
  They agree exactly when the clock is atomless, in particular in
  Example~\ref{ex:clocks}(i) --- which is why the choice is invisible in \EK{}
  and \CPS{}, and why neither source has to make it.

  We carry both.  The reason is that neither serves the whole manuscript:
  Theorem~\ref{thm:absreg} holds only for $\iota = \mathrm{o}$, because an atom of
  $q$ at $t$ must lie inside the compensator up to $t$ for (R3) to hold, and it is
  that limit which makes the c\`adl\`ag process a \emph{modification};
  Proposition~\ref{prop:haar} holds only for $\iota = \mathrm{p}$, because
  otherwise the chain identity telescopes to $\gamma(t,0)-\gamma(0,t) \neq 0$ and
  discrete-time duality fails.  Everything else --- Sections~\ref{sec:MP} and
  \ref{sec:uniqueness}, and Lemma~\ref{lem:chain} itself --- is
  convention-agnostic, using only the additivity \eqref{eq:clockadd}, which holds
  for both.  Remark~\ref{rem:conventionclash} states the trade-off precisely, and
  Section~\ref{ssec:bundletable} records which $\iota$ each result needs.

  Two special cases are worth keeping in view.  For $\T = \N_0$ with counting
  measure, $\iota = \mathrm{p}$ gives $C_n = \sum_{k<n} g(X_k)$, the Doob
  decomposition, with $A = P - I$; $\iota = \mathrm{o}$ gives
  $C_n = \sum_{k \leq n} g(X_k)$ and the condition $Pf - f = Pg$.  For
  $\T = \Rp$ with Lebesgue measure the distinction is empty.
\end{remark}

\begin{remark}[Comparison with \EK, Section 2.8]
  \label{rem:ekindex}
  \EK, Section~2.8 is the only part of \EK{} that works with a general index set,
  and it is worth recording exactly what it does, because three things about it
  are easy to get wrong.

  \emph{(i) What \EK{} assume.}  Their index $\mathscr{I}$ is a
  \emph{partially} ordered set --- their (8.1)--(8.3) include antisymmetry ---
  which is moreover a \emph{metric lattice}: a metric space in which $u \wedge v$
  and $u \vee v$ exist, are unique, and depend \emph{continuously} on $(u,v)$.
  Separability from above of the intervals is a further hypothesis, imposed where
  it is needed (their Propositions 8.4, 8.5(c) and Theorem~8.7), not globally.
  Our (T0) is weaker on one point: we do not assume antisymmetry, because nothing
  below uses it and Mathlib's \texttt{Preorder} does not have it.

  \emph{(ii) The lattice is needed for $\vee$, not for $\wedge$.}  It is tempting
  to say that stopping times require $\tau_1 \wedge \tau_2$.  The opposite is
  true.  \EK, Proposition 8.1(a) shows that $\max_{k \leq n} \tau_k$ is a stopping
  time --- $\{\max_k \tau_k \leq u\} = \bigcap_k \{\tau_k \leq u\}$ --- whereas
  their Remark~8.3 warns that $\tau \wedge a$ \emph{need not be} one.  The reason
  is that in a lattice $x \wedge y \leq u$ is strictly weaker than
  ``$x \leq u$ or $y \leq u$'': in $\Rp^2$ one has
  $(1,0) \wedge (0,1) = (0,0) \leq (0,0)$ with neither factor below $(0,0)$.  So
  $\{\tau \wedge a \leq u\}$ is genuinely larger than
  $\{\tau \leq u\} \cup \{a \leq u\}$ and there is no reason for it to lie in
  $\Filt_u$.  \EK{} therefore replace $\tau \wedge a$ by the truncation
  $\tau^a = \tau$ on $\{\tau \leq a\}$ and $\tau^a = a$ otherwise, their (8.10).
  Nor does $\Filt_\tau$ require infima: their (8.6) is the usual
  $\{A : A \cap \{\tau \leq u\} \in \Filt_u\}$, verbatim.

  \emph{(iii) What it is for.}  \EK{} build Section~2.8 for their Chapter~6, and
  it is worth being precise about how it is used there.  In their
  Theorem~6.3.4 the multiparameter index carries the \emph{filtration}
  $\{\Filt_u : u \in \Rp^k\}$ and a nondecreasing family of $\Rp^k$-valued
  stopping times $\tau(t)$; the processes themselves live in
  $\prod_k D_{E_k}[0,\infty)$, and so does the time-changed process
  $Z(t) = Y(\tau(t))$.  That is, \EK{} never index a \emph{process} by $\Rp^d$:
  the multiparameter index is a device inside a one-parameter theory, and what
  Section~2.8 delivers for it is optional sampling and nothing else.  This is
  consistent with --- indeed is indirect evidence for ---
  Remark~\ref{rem:chainonly}: there is no multiparameter uniqueness or Markov
  theorem in \EK{} to measure ours against, because there is none to be had.

  \emph{(iv) A conflation we avoid.}  \EK{} declare ``we assume throughout this
  section that $\mathscr{I}$ is a metric lattice'' before defining filtration,
  stopping time, $\Filt_\tau$, adaptedness and, in their (8.15), the martingale
  property itself --- none of which uses the lattice, and all of which are
  correct on a bare preorder.  The lattice is first used in their
  Proposition~8.1.  This is exactly the conflation that
  Definition~\ref{def:bundles} is meant to remove, and it is reassuring rather
  than surprising: the entire martingale layer of this manuscript rests on
  \EK{} (8.15) alone, which is stated for a partially ordered index.
\end{remark}

\subsubsection*{Notation for processes}

For an $E$-valued measurable process $X$ on $(\Omega,\Filt,P)$, indexed by $\T$,
we write $\Filt^X_t = \sigma(X(s) : s \leq t)$ and, following \EK, (4.3.2),
\begin{equation}
  \label{eq:starfilt}
  {}^{*}\Filt^X_t
  \;=\; \Filt^X_t \vee
        \sigma\Bigl( \int_{\langle 0,s \rangle_\iota} h(X(u)) \, q(\dif u)
          \;:\; s \leq t,\; h \in \Bdd(E) \Bigr).
\end{equation}
If $X$ is progressively measurable --- in particular if $\T$ satisfies (T2b) and
$X$ is right continuous --- then ${}^{*}\Filt^X_t = \Filt^X_t$.  Only (T0) and a
clock are needed for \eqref{eq:starfilt} to make sense.

% ---------------------------------------------------------------------
\subsection{The state space}
\label{ssec:statespace}
% ---------------------------------------------------------------------

\EK{} and \CPS{} take the state space $E$ to be Polish throughout.  As with the
time index, almost nothing needs that, and the parts that do need it need it for
visibly different reasons.  We therefore separate them.

\begin{definition}
  \label{def:Ebundles}
  We use the following bundles of assumptions on the state space.
  \begin{itemize}[nosep,leftmargin=3.2em]
  \item[\textbf{(E0)}] $(E,\mathcal{E})$ is a measurable space.
  \item[\textbf{(E1)}] (E0), and $(E,\mathcal{E})$ is \emph{standard Borel}:
    $\mathcal{E}$ is the Borel $\sigma$-field of some Polish topology on $E$,
    equivalently $(E,\mathcal{E})$ is Borel isomorphic to a Borel subset of a
    Polish space.
  \item[\textbf{(E2)}] (E1), and $E$ carries a separable metrizable topology with
    $\mathcal{E} = \Bor(E)$.
  \item[\textbf{(E3)}] $E$ is Polish.
  \end{itemize}
  Thus (E3) $\Rightarrow$ (E2) $\Rightarrow$ (E1) $\Rightarrow$ (E0), and both
  gaps are inhabited: a Borel non-$G_\delta$ subset of $\R$ satisfies (E2) but
  not (E3), and the space $\mathcal{S}'(\R^d)$ of tempered distributions
  satisfies (E1) but not (E2), being a Lusin space and hence standard Borel while
  not metrizable.
\end{definition}

\begin{remark}[What each bundle is for]
  \label{rem:Ebundleuse}
  The four bundles correspond to four different uses, and
  Section~\ref{ssec:bundletable} records which result makes which.
  \begin{enumerate}[label=(\roman*)]
  \item \emph{(E0): nothing but measurability.}  Sections~\ref{sec:MP} and
    \ref{sec:duality}, and the whole of Section~\ref{ssec:absuniq} apart from
    Lemma~\ref{lem:disint}, mention $E$ only through $\Bdd(E)$ and $\mathcal{E}$.
    \CPS{} already state their duality theorem for measurable $E_1$, $E_2$; the
    observation here is that uniqueness and the Markov property are of the same
    kind.
  \item \emph{(E1): regular conditional distributions.}  Exactly three items need
    them --- Lemma~\ref{lem:disint}, Fact~\ref{fact:sepcond}, and the
    identification of a law on $F$ by its finite-dimensional distributions.  Each
    is a statement about the Borel structure, not about the topology, and
    standard Borel is what makes disintegration available.
  \item \emph{(E2): separating classes and compact sets.}
    Section~\ref{sec:cadlag} needs a topology: the compact containment condition
    \eqref{eq:cc} refers to compact subsets of $E$, and the two hypotheses on
    $\dom(A)$ refer to $\Cb(E)$ separating points and separating measures.
  \item \emph{(E3): weak convergence.}  Only Section~\ref{sec:convergence}, and
    there only because $\DE$ has to be Polish for Prohorov's theorem and the
    Skorokhod representation.
  \end{enumerate}
\end{remark}

\begin{remark}[The point of (E1)]
  \label{rem:E1why}
  The gap between (E1) and (E2) is not academic.  Martingale problems for
  distribution-valued processes --- the martingale problem formulation of an
  SPDE, with $E = \mathcal{S}'(\R^d)$ or $E = \mathcal{D}'(\R^d)$ --- live on a
  state space that is standard Borel but not metrizable.  Under (E0) and (E1),
  Sections~\ref{sec:MP}, \ref{sec:uniqueness} and \ref{sec:duality} apply to them
  verbatim, with no new proof.  Sections~\ref{sec:cadlag} and
  \ref{sec:convergence} do not, and the repair is a genuine piece of work rather
  than a change of hypotheses: one needs a substitute for $\DE$ and for
  tightness in it, which for nuclear duals is Mitoma's theorem.  That is out of
  scope here, and is flagged rather than attempted.
\end{remark}

% ---------------------------------------------------------------------
\subsection{From Chapter 1: operator semigroups}
\label{ssec:ch1}
% ---------------------------------------------------------------------

\begin{definition}[\EK, Section 1.1]
  \label{def:semigroup}
  A one-parameter family $\{T(t) : t \geq 0\}$ of bounded linear operators on
  $\Lsp$ is a \emph{semigroup} if $T(0) = I$ and $T(s+t) = T(s)T(t)$ for all
  $s,t \geq 0$.  It is \emph{strongly continuous} if
  $\lim_{t \to 0} T(t) f = f$ for every $f \in \Lsp$, and a \emph{contraction
  semigroup} if $\lVert T(t) \rVert \leq 1$ for all $t \geq 0$.

  A (possibly unbounded) \emph{linear operator} $A$ on $\Lsp$ is a linear map
  whose domain $\dom(A)$ is a subspace of $\Lsp$ and whose range $\ran(A)$ lies
  in $\Lsp$.  The \emph{(infinitesimal) generator} of a semigroup $\{T(t)\}$ on
  $\Lsp$ is the linear operator $A$ defined by
  \begin{equation}
    \label{eq:generator}
    Af \;=\; \lim_{t \to 0} \tfrac{1}{t} \bigl( T(t) f - f \bigr) ,
  \end{equation}
  with $\dom(A)$ the set of $f \in \Lsp$ for which the limit exists.
\end{definition}

\begin{definition}[\EK, Section 1.2]
  \label{def:dissipative}
  A linear operator $A$ on $\Lsp$ is \emph{dissipative} if
  $\lVert \lambda f - A f \rVert \geq \lambda \lVert f \rVert$ for every
  $f \in \dom(A)$ and $\lambda > 0$.  For $A \subset \Lsp \times \Lsp$ linear,
  $A$ is dissipative if $\lVert \lambda f - g \rVert \geq \lambda \lVert f\rVert$
  for all $(f,g) \in A$ and $\lambda > 0$.
\end{definition}

\begin{definition}[\EK, Section 1.5]
  \label{def:fullgenerator}
  Let $\Lsp$ be a function space of the type considered in \EK, Section~1.5 (for
  instance $\Lsp = \Bdd(E)$ with the sup norm).  A semigroup $\{T(t)\}$ on
  $\Lsp$ is \emph{measurable} if $t \mapsto T(t) f$ is measurable for each
  $f \in \Lsp$.  The \emph{full generator} of a measurable contraction semigroup
  $\{T(t)\}$ on $\Lsp$ is the multivalued operator
  \begin{equation}
    \label{eq:fullgenerator}
    \hat{A} \;=\; \Bigl\{ (f,g) \in \Lsp \times \Lsp \;:\;
      T(t) f - f = \int_0^t T(s) g \dif s \text{ for all } t \geq 0 \Bigr\} .
  \end{equation}
\end{definition}

\begin{fact}[\EK, Proposition 1.5.1]
  \label{fact:fullgenerator}
  The full generator $\hat{A}$ of a measurable contraction semigroup
  $\{T(t)\}$ on $\Lsp$ is linear and dissipative, and
  $(\lambda - \hat{A})^{-1} h = \int_0^\infty e^{-\lambda t} T(t) h \dif t$ for
  all $h \in \ran(\lambda - \hat{A})$ and $\lambda > 0$.
\end{fact}

\begin{remark}
  \label{rem:noch1}
  Nothing further from \EK, Chapter~1 is used below.  In particular the
  Hille--Yosida theorem, cores, multivalued dissipative operators and the
  exponential formula (\EK, Theorems~1.2.6 and~1.4.3, Propositions~1.3.1
  and~1.3.3, Corollary~1.6.8) are \emph{not} needed.  They enter only through the
  semigroup-based uniqueness criteria \EK, Theorem~4.4.1 and Corollary~4.4.4,
  which are deliberately not part of this manuscript: uniqueness is obtained here
  from Theorem~\ref{thm:uniqueness} together with duality
  (Corollary~\ref{cor:uniqviadual}), which requires no semigroup theory at all.
  Should those criteria be wanted later, this subsection is the place to extend.
\end{remark}

% ---------------------------------------------------------------------
\subsection{From Chapter 2: stochastic processes and martingales}
\label{ssec:ch2}
% ---------------------------------------------------------------------

\begin{definition}[\EK, Section 2.1; (T0), except where noted]
  \label{def:process}
  Let $\T$ carry a clock (Definition~\ref{def:clock}), with $\sigma$-field
  $\mathcal{T}$.  An $E$-valued process $X$ on $(\Omega,\Filt,P)$ is
  \emph{measurable} if $X : \T \times \Omega \to E$ is
  $\mathcal{T} \otimes \Filt$-measurable.  A \emph{filtration} is an increasing
  family $(\Filt_t)_{t \in \T}$ of sub-$\sigma$-fields of $\Filt$.  $X$ is
  \emph{$(\Filt_t)$-adapted} if $X(t)$ is $\Filt_t$-measurable for each $t$, and
  \emph{$(\Filt_t)$-progressive} if for each $t \in \T$ the restriction of $X$ to
  $\T_{\leq t} \times \Omega$ is
  $(\mathcal{T} \cap \T_{\leq t}) \otimes \Filt_t$-measurable.

  $Y$ is a \emph{version} of $X$ if the two have the same finite-dimensional
  distributions; a \emph{modification} of $X$ if $P\{X(t) = Y(t)\} = 1$ for each
  $t \in \T$; and \emph{indistinguishable} from $X$ if
  $P\{X(t) = Y(t) \text{ for all } t \in \T\} = 1$.

  A $\T \cup \{\infty\}$-valued random variable $\tau$ is an
  \emph{$(\Filt_t)$-stopping time} if $\{\tau \leq t\} \in \Filt_t$ for every
  $t \in \T$, and then
  $\Filt_\tau = \{ F \in \Filt : F \cap \{\tau \leq t\} \in \Filt_t
  \ \forall t \in \T \}$.

  Under (T2b) one says in addition that $(\Filt_t)$ is \emph{right continuous} if
  $\Filt_t = \Filt_{t+} = \bigcap_{s>t}\Filt_s$ for all $t$, and that $X$ is
  \emph{right continuous} if its paths are; every right continuous adapted
  process is then progressive.  All of the above except this last paragraph
  needs only (T0), and is exactly \EK{} (8.1)--(8.6) of their Section~2.8; see
  Remark~\ref{rem:ekindex}.

\begin{fact}[\EK, Propositions 2.1.2 and 2.1.4; (T2b)]
  \label{fact:stoppingtimes}
  Assume (T2b).  If $\tau, \tau_1, \tau_2, \dots$ are $(\Filt_t)$-stopping
  times and $c \in \T$, then $\tau_1 + c$, $\tau_1 \wedge c$, $\sup_n \tau_n$ and
  $\min_{k \leq n} \tau_k$ are $(\Filt_t)$-stopping times; if $(\Filt_t)$ is
  right continuous, so are $\inf_n \tau_n$, $\varliminf_n \tau_n$ and
  $\varlimsup_n \tau_n$.  If $X$ is $(\Filt_t)$-progressive and $\tau$ is an
  $(\Filt_t)$-stopping time, then $X(\tau)$ is $\Filt_\tau$-measurable on
  $\{\tau < \infty\}$ and $X(\cdot \wedge \tau)$ is $(\Filt_t)$-progressive.
  The linear order is used throughout: on a merely directed index
  $\tau_1 \wedge \tau_2$ need not be a stopping time at all
  (Remark~\ref{rem:ekindex}(ii)).
\end{fact}

\begin{definition}[\EK, Section 2.2 and (8.15); (T0)]
  \label{def:martingale}
  A real-valued $(\Filt_t)$-adapted process $X$ with $E[|X(t)|] < \infty$ for all
  $t$ is an \emph{$(\Filt_t)$-martingale} if $E[X(t) \mid \Filt_s] = X(s)$ a.s.\
  for all $s \leq t$ in $\T$, and an \emph{$(\Filt_t)$-submartingale} if
  $E[X(t) \mid \Filt_s] \geq X(s)$ a.s.  This is \EK{} (8.15) verbatim, and it is
  the whole of what the martingale layer of this manuscript uses; it is also
  \texttt{MeasureTheory.Martingale} of Mathlib, which is stated for
  \texttt{[Preorder \(\iota\)]}.
\end{definition}

\begin{fact}[Regularization along a dense set; \EK, Lemma 2.2.8; (T2b)]
  \label{fact:cadlagext}
  Let $(E,r)$ be a metric space, $F \subset \T$ dense, and
  $x : F \to E$.  If
  $y(t) \coloneqq \lim_{s \to t+,\, s \in F} x(s)$ exists for every $t \in \T$,
  then $y$ is right continuous; if
  $y^-(t) \coloneqq \lim_{s \to t-,\, s \in F} x(s)$ exists for every $t > 0$,
  then $y^-$ is left continuous on $\T \setminus \{0\}$; and if both exist for
  every $t>0$, then $y^-(t) = y(t-)$ for all $t > 0$.  In particular $y$ is
  c\`adl\`ag, and $y \in \DE$ when $\T$ satisfies (T3).
\end{fact}

\begin{fact}[Submartingale regularization; \EK, Proposition 2.2.9; (T2b)]
  \label{fact:submgreg}
  Let $X$ be a real-valued $(\Filt_t)$-submartingale and let
  $F \subset \T$ be countable and dense.  Then there is $\Omega_0$ with
  $P(\Omega_0) = 1$ such that for every $\omega \in \Omega_0$ the limits
  \begin{equation}
    \label{eq:regY}
    Y(t,\omega) = \lim_{s \to t+,\, s \in F} X(s,\omega) \ \ (t \in \T),
    \qquad
    Y^-(t,\omega) = \lim_{s \to t-,\, s \in F} X(s,\omega) \ \ (t > 0)
  \end{equation}
  exist.  Moreover $\Gamma = \{t : P\{Y(t) \neq Y(t-)\} > 0\}$ is countable,
  $P\{X(t) = Y(t)\} = 1$ for $t \notin \Gamma$, and setting $\tilde{X}(t) = Y(t)$
  for $t \notin \Gamma$ and $\tilde{X}(t) = X(t)$ for $t \in \Gamma$ defines a
  modification of $X$ almost all of whose sample paths have right and left limits
  at all $t \in \T$ and are right continuous at all $t \notin \Gamma$.
\end{fact}

\begin{fact}[Optional sampling; \EK, Theorem 2.2.13, Remark 2.2.14, and Chapter 2, Problem 11; (T2b)]
  \label{fact:optsampl}
  Let $X$ be a right continuous $(\Filt_t)$-submartingale and let $\tau_1$,
  $\tau_2$ be $(\Filt_t)$-stopping times.  Then for each $T > 0$,
  \begin{equation}
    \label{eq:optsampl}
    E\bigl[ X(\tau_2 \wedge T) \bigm| \Filt_{\tau_1} \bigr]
      \;\geq\; X(\tau_1 \wedge \tau_2 \wedge T) .
  \end{equation}
  If in addition $\tau_2$ is a.s.\ finite, $E[|X(\tau_2)|] < \infty$ and
  $\lim_{T \to \infty} E[|X(T)| \one_{\{\tau_2 > T\}}] = 0$, then
  $E[X(\tau_2) \mid \Filt_{\tau_1}] \geq X(\tau_1 \wedge \tau_2)$.  Equality holds
  throughout if $X$ is a martingale.  In particular, a right continuous
  martingale with bounded increments satisfies the optional sampling identity at
  arbitrary a.s.\ finite stopping times.
\end{fact}

\begin{fact}[Doob's inequalities; \EK, Corollary 2.2.17; (T2b)]
  \label{fact:doob}
  Let $X$ be a right continuous martingale.  Then for $x > 0$, $T > 0$ and
  $\alpha > 1$,
  \begin{equation}
    \label{eq:doob}
    P\Bigl\{ \sup_{t \leq T} |X(t)| \geq x \Bigr\} \leq x^{-1} E[|X(T)|],
    \qquad
    E\Bigl[ \sup_{t \leq T} |X(t)|^\alpha \Bigr]
      \leq \Bigl( \tfrac{\alpha}{\alpha-1} \Bigr)^{\!\alpha} E[|X(T)|^\alpha] .
  \end{equation}
\end{fact}

\begin{definition}[\EK, Section 2.3]
  \label{def:localmartingale}
  A real-valued process $X$ is an \emph{$(\Filt_t)$-local martingale} if there
  exist $(\Filt_t)$-stopping times $\tau_1 \leq \tau_2 \leq \cdots$ with
  $\tau_n \to \infty$ a.s.\ such that $X^{\tau_n} = X(\cdot \wedge \tau_n)$ is an
  $(\Filt_t)$-martingale for each $n$.  The sequence $(\tau_n)$ is a
  \emph{localizing sequence}.  Local submartingales are defined analogously.
\end{definition}

\begin{fact}[\EK, Proposition 2.3.1]
  \label{fact:stoppedlocalmg}
  If $X$ is a right continuous $(\Filt_t)$-local martingale and $\tau$ is an
  $(\Filt_t)$-stopping time, then $X^\tau = X(\cdot \wedge \tau)$ is an
  $(\Filt_t)$-local martingale.
\end{fact}

% ---------------------------------------------------------------------
\subsection{From Chapter 3: weak convergence and the Skorokhod space}
\label{ssec:ch3}
% ---------------------------------------------------------------------

Everything in this subsection assumes (T3), $\T = [0,\infty)$, and is stated for
it without further comment: the Skorokhod topology is defined through time
changes $\lambda : [0,\infty) \to [0,\infty)$ (Definition~\ref{def:skorokhod})
and has no counterpart on a general index.  This is the one block of
prerequisites that does not abstract, and
Section~\ref{ssec:bundletable} records which results inherit the restriction.

Throughout this subsection $(S,d)$ is a metric space.

\begin{definition}[\EK, Section 3.3]
  \label{def:weakconv}
  $P_n \wto P$ in $\Prob(S)$ (\emph{weak convergence}) if
  $\int f \dif P_n \to \int f \dif P$ for every $f \in \Cb(S)$.  $S$-valued
  random variables $X_n$ \emph{converge in distribution} to $X$, written
  $X_n \wto X$, if $P X_n^{-1} \wto P X^{-1}$.
\end{definition}

\begin{fact}[Portmanteau; \EK, Theorem 3.3.1]
  \label{fact:portmanteau}
  For $\{P_n\} \subset \Prob(S)$ and $P \in \Prob(S)$, conditions (b)--(f) below
  are equivalent and are implied by~(a); if $S$ is separable, all six are
  equivalent.
  \begin{enumerate}[label=(\alph*), nosep]
  \item $\rho(P_n, P) \to 0$, where $\rho$ is the Prohorov metric;
  \item $P_n \wto P$;
  \item $\int f \dif P_n \to \int f \dif P$ for all uniformly continuous
    $f \in \Cb(S)$;
  \item $\varlimsup_n P_n(F) \leq P(F)$ for all closed $F \subset S$;
  \item $\varliminf_n P_n(G) \geq P(G)$ for all open $G \subset S$;
  \item $P_n(A) \to P(A)$ for all $P$-continuity sets $A \subset S$, i.e.\ all
    $A \in \Bor(S)$ with $P(\partial A) = 0$.
  \end{enumerate}
\end{fact}

\begin{fact}[Continuous mapping theorem; \EK, Corollary 3.1.9 and Corollary 3.3.3]
  \label{fact:cmt}
  Let $(S,d)$ and $(S',d')$ be separable metric spaces and $h : S \to S'$ Borel
  measurable.  Let $C_h$ be the set of points at which $h$ is continuous.  If
  $P_n \wto P$ and $P(C_h) = 1$, then $P_n h^{-1} \wto P h^{-1}$.  Moreover, if
  $X_n \wto X$ and $d(X_n, Y_n) \to 0$ in probability, then $Y_n \wto X$.
\end{fact}

\begin{fact}[\EK, Theorems 3.1.7 and 3.1.8]
  \label{fact:PSpolish}
  If $S$ is separable, then $\Prob(S)$ is separable; if $(S,d)$ is complete and
  separable, then so is $\Prob(S)$ with the Prohorov metric.  Moreover
  (\emph{Skorokhod representation}), if $S$ is separable and $P_n \wto P$, there
  are $S$-valued random variables $X_n$, $X$ on a common probability space with
  laws $P_n$, $P$ and $X_n \to X$ a.s.
\end{fact}

\begin{fact}[Prohorov; \EK, Lemma 3.2.1 and Theorem 3.2.2]
  \label{fact:prohorov}
  Call $\mathcal{M} \subset \Prob(S)$ \emph{tight} if for every
  $\varepsilon > 0$ there is a compact $K \subset S$ with
  $\inf_{P \in \mathcal{M}} P(K) \geq 1 - \varepsilon$.  If $(S,d)$ is complete
  and separable, then every $P \in \Prob(S)$ is tight, and
  $\mathcal{M} \subset \Prob(S)$ is tight if and only if it is relatively
  compact.
\end{fact}

\begin{definition}[\EK, Section 3.4 and Appendix 3]
  \label{def:bp}
  A sequence $\{f_n\} \subset \Bdd(S)$ \emph{converges boundedly and pointwise}
  to $f \in \Bdd(S)$, written $\bplim_n f_n = f$, if $\sup_n \lVert f_n \rVert
  < \infty$ and $f_n(x) \to f(x)$ for every $x \in S$.  A set
  $M \subset \Bdd(S)$ is \emph{bp-closed} if it is closed under bp-limits of
  sequences, and the \emph{bp-closure} of $M$ is the smallest bp-closed set
  containing it.  The same terminology is used for subsets of
  $\Bdd(S) \times \Bdd(S)$, applying $\bplim$ in each coordinate.
\end{definition}

\begin{fact}[\EK, Lemma 3.4.1, Proposition 3.4.2, and Appendix 3, Proposition 3.1]
  \label{fact:bp}
  The bp-closure of a subspace of $\Bdd(S)$ is a subspace.  $\Cb(S)$ is bp-dense
  in $\Bdd(S)$, and if $S$ is separable there is a sequence $\{f_n\}$ of
  nonnegative functions in $\Cb(S)$ whose linear span is bp-dense in $\Bdd(S)$.
  Furthermore, for sequences in $\Bdd(S)$ bounded pointwise convergence coincides
  with convergence in the weak$^*$ topology induced by $\Vsig(S)$.
\end{fact}

\begin{definition}[\EK, Section 3.4]
  \label{def:separating}
  $M \subset \Cb(S)$ is \emph{separating} if, for $P,Q \in \Prob(S)$,
  $\int f \dif P = \int f \dif Q$ for all $f \in M$ implies $P = Q$; it is
  \emph{convergence determining} if, for $\{P_n\} \subset \Prob(S)$ and
  $P \in \Prob(S)$, $\int f \dif P_n \to \int f \dif P$ for all $f \in M$
  implies $P_n \wto P$.  $M$ \emph{separates points} if for $x \neq y$ there is
  $f \in M$ with $f(x) \neq f(y)$, and \emph{strongly separates points} if for
  every $x \in S$ and $\delta > 0$ there are $h_1,\dots,h_k \in M$ with
  $\inf_{y : d(y,x) \geq \delta} \max_{i \leq k} |h_i(y) - h_i(x)| > 0$.
  Every convergence determining set is separating.
\end{definition}

\begin{fact}[\EK, Proposition 3.4.4]
  \label{fact:convdet}
  If $(S,d)$ is separable, then the set of uniformly continuous $f \in \Cb(S)$
  with bounded support is convergence determining.  If $S$ is in addition
  locally compact, then $C_c(S)$ is convergence determining.
\end{fact}

\begin{fact}[Stone--Weierstrass for separating classes; \EK, Theorem 3.4.5]
  \label{fact:stoneweierstrass}
  Let $(S,d)$ be complete and separable and let $M \subset \Cb(S)$ be an algebra.
  If $M$ separates points, then $M$ is separating.  If $M$ strongly separates
  points, then $M$ is convergence determining.
\end{fact}

\begin{fact}[\EK, Proposition 3.4.6 and Proposition 3.7.1]
  \label{fact:fdd}
  Let $(S_k, d_k)$, $k \geq 1$, be separable and $S = \prod_{k \geq 1} S_k$.  If
  $M_k \subset \Cb(S_k)$ are separating, then
  \begin{equation}
    \label{eq:prodsep}
    M = \Bigl\{ x \mapsto \textstyle\prod_{k=1}^n f_k(x_k) \;:\;
      n \geq 1,\ f_k \in M_k \cup \{1\} \Bigr\}
  \end{equation}
  is separating on $S$; if the $(S_k,d_k)$ are complete and separable and the
  $M_k$ are convergence determining, so is $M$.  In particular the
  finite-dimensional distributions of a process determine its law.  Moreover, for
  $\pi_t : \DE \to E$, $\pi_t(x) = x(t)$, one has
  $\Bor(\DE) = \sigma(\pi_t : t \geq 0) = \sigma(\pi_t : t \in D)$ for any dense
  $D \subset [0,\infty)$, provided $E$ is separable; hence the
  finite-dimensional distributions of a process with sample paths in $\DE$
  determine its law on $\DE$.
\end{fact}

\begin{definition}[\EK, Section 3.5]
  \label{def:skorokhod}
  $\DE$ is the space of right continuous $x : [0,\infty) \to E$ having left
  limits $x(t-)$ for all $t > 0$ (with $x(0-) = x(0)$).  Let $q = r \wedge 1$,
  let $\Lambda'$ be the set of strictly increasing continuous
  $\lambda : [0,\infty) \to [0,\infty)$ onto with $\lambda(0)=0$, and let
  $\Lambda \subset \Lambda'$ be those $\lambda$ which are Lipschitz with
  \begin{equation}
    \label{eq:gamma}
    \gamma(\lambda) \;=\; \operatorname*{ess\,sup}_{t \geq 0} |\log \lambda'(t)|
      \;=\; \sup_{s > t \geq 0}
        \Bigl| \log \tfrac{\lambda(s)-\lambda(t)}{s-t} \Bigr| \;<\; \infty .
  \end{equation}
  The \emph{Skorokhod ($J_1$) metric} on $\DE$ is
  \begin{equation}
    \label{eq:skorokhodmetric}
    d(x,y) \;=\; \inf_{\lambda \in \Lambda}
      \Bigl[ \gamma(\lambda) \vee
        \int_0^\infty e^{-u} d(x,y,\lambda,u) \dif u \Bigr],
  \end{equation}
  where
  \begin{equation}
    \label{eq:skorokhodmetricaux}
    d(x,y,\lambda,u)
      = \sup_{t \geq 0} q\bigl( x(t \wedge u), y(\lambda(t) \wedge u) \bigr).
  \end{equation}
\end{definition}

\begin{fact}[\EK, Lemma 3.5.1 and Theorem 3.5.6]
  \label{fact:DEpolish}
  Every $x \in \DE$ has at most countably many discontinuities.  If $E$ is
  separable, then $\DE$ is separable; if $(E,r)$ is complete, then $(\DE, d)$ is
  complete.  In particular, $\DE$ is Polish whenever $E$ is Polish.
\end{fact}

\begin{fact}[Compact sets of $\DE$; \EK, Lemma 3.6.2 and Theorem 3.6.3]
  \label{fact:DEcompact}
  For $x \in \DE$, $\delta > 0$ and $T > 0$ define the modulus
  \begin{equation}
    \label{eq:modulus}
    w'(x,\delta,T) \;=\; \inf_{\{t_i\}} \max_i
      \sup_{s,t \in [t_{i-1}, t_i)} r\bigl( x(s), x(t) \bigr) ,
  \end{equation}
  the infimum being over all partitions $0 = t_0 < t_1 < \dots < t_n$ with
  $T \leq t_n$ and $\min_{i} (t_i - t_{i-1}) > \delta$.  Then
  $w'(x,\delta,T) \to 0$ as $\delta \to 0$, $w'(\cdot,\delta,T)$ is Borel
  measurable, and, for $(E,r)$ complete, the closure of $A \subset \DE$ is
  compact if and only if
  \begin{enumerate}[label=(\alph*), nosep]
  \item for every rational $t \geq 0$ there is a compact $\Gamma_t \subset E$
    with $x(t) \in \Gamma_t$ for all $x \in A$, and
  \item $\lim_{\delta \to 0} \sup_{x \in A} w'(x,\delta,T) = 0$ for each $T>0$.
  \end{enumerate}
\end{fact}

\begin{fact}[\EK, Lemma 3.7.7]
  \label{fact:Dcountable}
  If $X$ has sample paths in $\DE$, then
  $D(X) \coloneqq \{t \geq 0 : P\{X(t) = X(t-)\} = 1\}$ has at most countable
  complement in $[0,\infty)$.  In particular $D(X)$ is dense.
\end{fact}

\begin{fact}[\EK, Theorem 3.7.8]
  \label{fact:fddconv}
  Let $E$ be separable and let $X_n$, $X$ have sample paths in $\DE$.
  \begin{enumerate}[label=(\alph*), nosep]
  \item If $X_n \wto X$ in $\DE$, then
    $(X_n(t_1),\dots,X_n(t_k)) \wto (X(t_1),\dots,X(t_k))$ for every finite
    $\{t_1,\dots,t_k\} \subset D(X)$.
  \item If $\{X_n\}$ is relatively compact and there is a dense
    $D \subset [0,\infty)$ such that the convergence in (a) holds for every
    finite subset of $D$, then $X_n \wto X$ in $\DE$.
  \end{enumerate}
\end{fact}

\begin{fact}[Relative compactness, I; \EK, Theorem 3.9.1]
  \label{fact:relcompact}
  Let $(E,r)$ be complete and separable and let $\{X_n\}$ be a family of
  processes with sample paths in $\DE$ satisfying the compact containment
  condition: for every $\eta, T > 0$ there is a compact
  $\Gamma_{\eta,T} \subset E$ with
  \begin{equation}
    \label{eq:ccn}
    \inf_n P\{ X_n(t) \in \Gamma_{\eta,T} \text{ for } 0 \leq t \leq T \}
      \;\geq\; 1 - \eta .
  \end{equation}
  Let $H \subset \Cb(E)$ be dense in the topology of uniform convergence on
  compact sets.  Then $\{X_n\}$ is relatively compact if and only if
  $\{f \circ X_n\}$ is relatively compact in $D_{\R}[0,\infty)$ for each
  $f \in H$.
\end{fact}

\begin{fact}[Relative compactness, II; \EK, Theorem 3.9.4]
  \label{fact:relcompact2}
  Let $(E,r)$ be arbitrary and let $X_n$ have sample paths in $\DE$, adapted to
  filtrations $(\Filt^n_t)$.  Let $\mathcal{L}_n$ be the Banach space of
  real-valued $(\Filt^n_t)$-progressive processes with
  $\lVert Y \rVert = \sup_{t \geq 0} E[|Y(t)|] < \infty$ and put
  \begin{equation}
    \label{eq:An}
    \mathcal{A}_n \;=\; \Bigl\{ (Y,Z) \in \mathcal{L}_n \times \mathcal{L}_n :
      Y(t) - \int_0^t Z(s) \dif s \text{ is an } (\Filt^n_t)\text{-martingale} \Bigr\}.
  \end{equation}
  Let $C_a \subset \Cb(E)$ be a subalgebra and let $D$ be the collection of
  $f \in \Cb(E)$ such that for every $\varepsilon, T > 0$ there is
  $(Y_n, Z_n) \in \mathcal{A}_n$ with
  \begin{equation}
    \label{eq:relcompact2}
    \sup_n E\Bigl[ \sup_{t \in [0,T] \cap \Q} |Y_n(t) - f(X_n(t))| \Bigr]
      < \varepsilon,
    \qquad
    \sup_n E\bigl[ \lVert Z_n \rVert_{p,T} \bigr] < \infty
    \text{ for some } p \in (1,\infty].
  \end{equation}
  If $C_a$ is contained in the sup-norm closure of $D$, then
  $\{f \circ X_n\}$ is relatively compact for each $f \in C_a$; more generally
  $\{(f_1,\dots,f_k) \circ X_n\}$ is relatively compact in
  $D_{\R^k}[0,\infty)$ for all $f_1,\dots,f_k \in C_a$.
\end{fact}

% ---------------------------------------------------------------------
\subsection{From the Appendixes of \EK}
\label{ssec:appendix}
% ---------------------------------------------------------------------

\begin{fact}[Uniform integrability; \EK, Appendix 2]
  \label{fact:ui}
  A family $\{X_i\}$ of real-valued random variables is \emph{uniformly
  integrable} if $\sup_i E[|X_i|] < \infty$ and for every $\varepsilon > 0$ there
  is $\delta > 0$ such that $P(A) < \delta$ implies
  $|E[X_i \one_A]| < \varepsilon$ for all $i$.  Equivalently,
  $\lim_{N \to \infty} \sup_i E[|X_i| - |X_i| \wedge N] = 0$; equivalently, there
  is an increasing convex $\varphi$ on $[0,\infty)$ with
  $\varphi(x)/x \to \infty$ and $\sup_i E[\varphi(|X_i|)] < \infty$.  If
  $X_n \wto X$ and $\{X_n\}$ is uniformly integrable, then
  $E[X_n] \to E[X]$.
\end{fact}

\begin{fact}[Monotone class theorem; \EK, Appendix 4]
  \label{fact:monotoneclass}
  Let $H$ be a linear space of bounded real-valued functions on $\Omega$
  containing the constants and closed under bounded monotone limits, and let
  $K \subset H$ be closed under multiplication.  Then $H$ contains every bounded
  $\sigma(K)$-measurable function.
\end{fact}

\begin{fact}[Conditional determination by separating sets; \EK, Chapter 3, Problem 7]
  \label{fact:sepcond}
  Assume (E1) and (E2), let $\Gilt \subset \Filt$ be a sub-$\sigma$-field, let $M \subset \Cb(E)$ be
  separating, and let $U,V$ be $E$-valued random variables with $V$
  $\Gilt$-measurable.  If $E[f(U) \mid \Gilt] = f(V)$ a.s.\ for every $f \in M$,
  then $P\{U = V\} = 1$.
\end{fact}

\begin{remark}[\EK, Chapter 3, Problem 7 --- a proof]
  \label{rem:sepcondproof}
  This is stated as a \emph{problem} in \EK{} and is used in the last step of
  Theorem~\ref{thm:absreg}, so we prove it.  The point worth making is that the
  countability the argument needs comes from the \emph{state space}, not from
  $M$: no countable subfamily of $M$ is assumed, and none is available in general.

  \emph{Step 1: $M$ separating gives conditional equality in law.}
  Let $G \in \Gilt$ with $P(G) > 0$.  For $f \in M$, the hypothesis gives
  $E[\one_G f(U)] = E[\one_G E[f(U)\mid\Gilt]] = E[\one_G f(V)]$.  The two
  probability measures
  \[
    \nu_1(B) = \frac{P(\{U \in B\} \cap G)}{P(G)} , \qquad
    \nu_2(B) = \frac{P(\{V \in B\} \cap G)}{P(G)} ,
    \qquad B \in \mathcal{E},
  \]
  therefore satisfy $\int f \dif\nu_1 = \int f \dif\nu_2$ for every $f \in M$, and
  since $M$ is separating, $\nu_1 = \nu_2$.  Hence
  \begin{equation}
    \label{eq:sepcondstep}
    P\bigl( \{U \in B\} \cap G \bigr) = P\bigl( \{V \in B\} \cap G \bigr)
    \qquad \text{for all } B \in \mathcal{E},\ G \in \Gilt ,
  \end{equation}
  i.e.\ $P(U \in B \mid \Gilt) = \one_B(V)$ a.s.\ for each fixed $B$.

  \emph{Step 2: one null set, using (E1).}
  By (E1) the $\sigma$-field $\mathcal{E}$ is countably generated; let
  $\mathcal{E}_0$ be a countable algebra generating it.  Again by (E1) there is a
  regular conditional distribution $\mu(\omega,\cdot)$ of $U$ given $\Gilt$.
  Applying \eqref{eq:sepcondstep} to the countably many $B \in \mathcal{E}_0$
  yields a single set of full measure off which
  $\mu(\omega,B) = \one_B(V(\omega))$ for all $B \in \mathcal{E}_0$
  simultaneously.  Two probability measures agreeing on a generating algebra
  agree, so $\mu(\omega,\cdot) = \delta_{V(\omega)}$ a.s.

  \emph{Step 3.}  The diagonal $\{(x,\omega) : x = V(\omega)\}$ is
  $\mathcal{E}\otimes\Gilt$-measurable, again because $\mathcal{E}$ is countably
  generated and separates points, so
  $P\{U = V\} = E\bigl[ \mu(\cdot,\{V(\cdot)\}) \bigr]
   = E\bigl[ \delta_{V}(\{V\}) \bigr] = 1$.

  Only (E2) is used, in Step~1, to make sense of ``$M \subset \Cb(E)$
  separating''; all the measure theory is (E1).
\end{remark}

% ---------------------------------------------------------------------
\subsection{Where the prerequisites are used}
\label{ssec:usage}
% ---------------------------------------------------------------------

\noindent
\begin{tabular}{@{}ll@{}}
\toprule
Prerequisite & Used in \\
\midrule
Definition~\ref{def:dissipative}
  & Proposition~\ref{prop:dissipative}, Fact~\ref{fact:fullgenerator} \\
Fact~\ref{fact:fullgenerator}
  & Remark~\ref{rem:fullgenerator} \\
Fact~\ref{fact:stoppingtimes}
  & Theorem~\ref{thm:uniqueness}\ref{it:uniq_b}, Corollary~\ref{cor:dualstopped} \\
Fact~\ref{fact:cadlagext}, \ref{fact:submgreg}
  & Theorem~\ref{thm:absreg} \\
Fact~\ref{fact:optsampl}
  & Theorem~\ref{thm:absstrongmarkov} \\
Fact~\ref{fact:doob}
  & Remark~\ref{rem:EKrelcompact} (via Fact~\ref{fact:relcompact2}) \\
Fact~\ref{fact:stoppedlocalmg}
  & Remarks~\ref{rem:uniquelocal}, \ref{rem:convlocal} \\
Fact~\ref{fact:portmanteau}, \ref{fact:cmt}
  & Lemma~\ref{lem:EKconv}, Theorem~\ref{thm:CPSconv} \\
Fact~\ref{fact:PSpolish}, \ref{fact:prohorov}
  & Remark~\ref{rem:EKrelcompact} \\
Fact~\ref{fact:bp}
  & Remark~\ref{rem:fddconsequences} \\
Fact~\ref{fact:convdet}
  & Remark~\ref{rem:ccverify} \\
Fact~\ref{fact:stoneweierstrass}
  & Remark~\ref{rem:EKrelcompact} \\
Fact~\ref{fact:fdd}
  & Theorem~\ref{thm:absuniq}, Corollary~\ref{cor:DEuniqueness},
    Example~\ref{ex:determining} \\
Fact~\ref{fact:DEpolish}
  & Setting~\ref{set:abstract}, Theorem~\ref{thm:CPSconv} \\
Fact~\ref{fact:DEcompact}
  & Fact~\ref{fact:relcompact} \\
Fact~\ref{fact:Dcountable}
  & Lemma~\ref{lem:EKconv}, Theorem~\ref{thm:CPSconv} \\
Fact~\ref{fact:fddconv}
  & Remark~\ref{rem:EKrelcompact} \\
Fact~\ref{fact:relcompact}, \ref{fact:relcompact2}
  & Remark~\ref{rem:EKrelcompact} \\
Fact~\ref{fact:ui}
  & Theorem~\ref{thm:CPSconv} \\
Fact~\ref{fact:monotoneclass}
  & Proposition~\ref{prop:fddchar}, Example~\ref{ex:determining} \\
Fact~\ref{fact:sepcond}
  & Theorem~\ref{thm:absreg} \\
\bottomrule
\end{tabular}

% ---------------------------------------------------------------------
\subsection{Which result needs which bundle}
\label{ssec:bundletable}
% ---------------------------------------------------------------------

\noindent
The counterpart of Section~\ref{ssec:usage} for
Definitions~\ref{def:bundles}, \ref{def:clock} and \ref{def:Ebundles}.
``Clock'' means Definition~\ref{def:clock} with no further property;
``reflective'' is as defined there, and $\iota$ is the convention of
Definition~\ref{def:clock}, shown only where it matters.  Bold entries mark the
places where the weakest bundle in its column does \emph{not} suffice, for a
mathematical reason rather than for convenience; the two $\iota$-entries are the
only places in the manuscript where the two conventions part company.

\medskip
\noindent
\begin{tabular}{@{}p{0.38\textwidth}p{0.22\textwidth}p{0.13\textwidth}p{0.19\textwidth}@{}}
\toprule
\textbf{Result} & \textbf{Index} & \textbf{State} & \textbf{Clock} \\
\midrule
Definition~\ref{def:process}, Definition~\ref{def:martingale}
  & (T0) & (E0) & clock \\
Facts~\ref{fact:stoppingtimes}, \ref{fact:optsampl}, \ref{fact:doob}
  & (T2b) & (E0) & --- \\
Facts~\ref{fact:cadlagext}, \ref{fact:submgreg} (regularization)
  & (T2b) & (E2) & --- \\
Section~\ref{ssec:ch3} entire (Skorokhod, weak convergence)
  & (T3) & (E3) & --- \\
\midrule
Setting~\ref{set:abstract}, Definitions~\ref{def:absMP}, \ref{def:canonical}
  & (T0) & (E0) & --- \\
Definition~\ref{def:markovMP}, Proposition~\ref{prop:fddchar},
  Definition~\ref{def:wellposed}, \eqref{eq:starfilt}
  & (T0) & (E0) & clock \\
Lemma~\ref{lem:resolvent}, Proposition~\ref{prop:dissipative}
  & (T3) & (E0) & Lebesgue \\
Definition~\ref{def:cc}, Example~\ref{ex:determining}
  & (T2b) & (E2) & --- \\
\midrule
Definition~\ref{def:regclass}, Theorem~\ref{thm:absreg} (c\`adl\`ag)
  & (T2b) & (E2) & $\boldsymbol{\iota = \mathrm{o}}$ \\
Theorem~\ref{thm:cadlag} (\EK{} 4.3.6)
  & (T2b) & (E2) & clock, $\boldsymbol{\iota = \mathrm{o}}$ \\
\midrule
Lemma~\ref{lem:mixture} (mixtures)
  & (T0) & (E0) & --- \\
Lemma~\ref{lem:disint} (disintegration)
  & (T0) & \textbf{(E1)} & --- \\
Definition~\ref{def:shiftstable}, Lemma~\ref{lem:restart} (restart)
  & (T0) + (T4) & (E0) & clock \\
Theorem~\ref{thm:absuniq}\ref{it:absuniq_a} (Markov)
  & (T0) + (T4) & (E0) & clock \\
Theorem~\ref{thm:absuniq}\ref{it:absuniq_b} (uniqueness)
  & \textbf{(T2a)} + (T4) & (E0) & clock \\
Theorem~\ref{thm:absstrongmarkov} (strong Markov)
  & (T2b) + (T4) & (E2) & clock \\
Definition~\ref{def:localizing}, Lemmas~\ref{lem:localmix},
  \ref{lem:L1auto}, \ref{lem:localrestart},
  Theorem~\ref{thm:localuniq} (local theory)
  & (T2b) + (T4) & (E0) & clock \\
Definition~\ref{def:pasting}, Lemma~\ref{lem:pasting},
  Theorem~\ref{thm:localuniqueness} (local uniqueness)
  & (T2b) + (T4) & (E0) & clock \\
\midrule
Lemma~\ref{lem:chain} (chain identity)
  & (T0) + (T4) & (E0) & clock \\
Proposition~\ref{prop:haar}, Corollary~\ref{cor:dualdiscrete}
  & \textbf{(T2a)} + (T4) & (E0) & \textbf{reflective},
    $\boldsymbol{\iota = \mathrm{p}}$ \\
Lemma~\ref{lem:calculus}, Theorem~\ref{thm:duality}
  & (T3) & (E0) & Lebesgue \\
Corollary~\ref{cor:uniqviadual} (uniqueness via duality)
  & (T3) & (E2) & Lebesgue \\
\midrule
Setting~\ref{set:jumpdata}, Theorem~\ref{thm:jumpMP} (jump processes)
  & (T3) & (E0) & Lebesgue \\
Proposition~\ref{prop:jumpwellposed}, Remark~\ref{rem:jumpexplosion}
  & (T3) & (E2) & Lebesgue \\
Corollary~\ref{cor:sdewellposed} (from SDEs)
  & (T3) & $E = \R^d$ & Lebesgue \\
Theorem~\ref{thm:absconv} (abstract convergence)
  & (T0), or (T2b) for $D \subsetneq \T$ & (E3) & --- \\
Lemma~\ref{lem:EKconv}, Theorem~\ref{thm:CPSconv}
  & (T3) & (E3) & Lebesgue \\
\bottomrule
\end{tabular}

\medskip
\noindent
Two readings of this table are worth recording.  First, the partially ordered
case --- and with it the multiparameter index $\T = \Rp^d$ --- survives in
Section~\ref{sec:MP}, in the Markov half of Section~\ref{sec:uniqueness} and in
the abstract convergence theorem, and dies everywhere else; the two bold entries
are where it dies for a \emph{mathematical} reason rather than for want of a
Skorokhod space, and Remarks~\ref{rem:chainonly} and \ref{rem:dualscope} say
which.  Second, (T1$'$) --- the metric lattice of \EK, Section~2.8, which the
first draft of this manuscript treated as the base assumption --- does not occur
at all: every statement involving a stopping time already assumes (T2b), under
which the lattice structure is automatic.  See Remark~\ref{rem:ekindex}.

% =====================================================================
\section{Martingale problems}
\label{sec:MP}
% =====================================================================

\subsection{The abstract (local) martingale problem}

\begin{setting}[after \CPS, Section 3.1; (T0) + (E0)]
  \label{set:abstract}
  Let $(\Omega, \Filt, \mathbb{F} = (\Filt_t)_{t \in \T})$ be a filtered space
  carrying a family $\XX$ of real-valued adapted processes (the \emph{test
  processes}), right continuous if $\T$ satisfies (T2b).  Let $(E,\mathcal{E})$
  be a measurable space and let $F \subset E^{\T}$ be a \emph{path space}: a set
  of maps $\T \to E$ equipped with a $\sigma$-field $\mathcal{S}$ for which the
  coordinates $\pi_t : F \to E$ are measurable and which they generate,
  \begin{equation}
    \label{eq:pathsigma}
    \mathcal{S} \;=\; \sigma\bigl( \pi_t \;:\; t \in \T \bigr) .
  \end{equation}
  Let $X = (X_t)_{t \in \T}$ be an $E$-valued process with $X(\omega) \in F$ for
  every $\omega$ and $\omega \mapsto X(\omega)$ measurable
  $\Filt / \mathcal{S}$.

  Under (E3) we take $F$ to carry a Polish topology and
  $\mathcal{S} = \Bor(F)$; \eqref{eq:pathsigma} is then a theorem rather than an
  assumption for $F = \DE$ (Fact~\ref{fact:fdd}), and it is what makes the
  finite-dimensional distributions determine the law.
\end{setting}

\begin{definition}[\CPS, Definition 3.2; (T0) + (E0)]
  \label{def:absMP}
  A probability measure $P$ on $(\Omega, \Filt)$ is a \emph{solution of the
  martingale problem $\MP(\XX)$} if every $Y \in \XX$ is an
  $(\mathbb{F},P)$-martingale, and a \emph{solution of the local martingale
  problem} if every $Y \in \XX$ is an $(\mathbb{F},P)$-local martingale.  The
  corresponding sets of solutions are denoted $\Msol(\XX)$ and
  $\Msol_{\mathrm{loc}}(\XX)$.
\end{definition}

\begin{definition}[\CPS, Definitions 3.3 and 3.5; (T0) + (E0)]
  \label{def:canonical}
  $\XX$ is \emph{canonical} for $X$ if for every $t \in \T$ there is a set
  $\XX^\circ_t$ of $\Bor(F)/\Bor(\R)$-measurable maps $Y^\circ_t : F \to \R$
  such that every $Y \in \XX$ satisfies $Y_t = Y^\circ_t(X)$ for some
  $Y^\circ_t \in \XX^\circ_t$.  A family
  $\ZZ^\circ = \{\ZZ^\circ_t : t \in \T\}$ of sets of bounded measurable maps
  $F \to \R$ is \emph{determining} for $\XX$ if for every probability measure
  $P$, every $Y \in \XX$ and all $s \leq t$ with $Y_s, Y_t \in L^1(P)$,
  \[
    E^P[Y_t Z^\circ_s(X)] = E^P[Y_s Z^\circ_s(X)]
      \quad \text{for all } Z^\circ_s \in \ZZ^\circ_s
    \qquad \Longrightarrow \qquad
    E^P[Y_t \mid \Filt_s] = Y_s \ \ P\text{-a.s.}
  \]
\end{definition}

\begin{example}[\CPS, Example 3.6(i); (T2b) + (E2)]
  \label{ex:determining}
  If $\Filt_t = \sigma(X_s : s \leq t)$, $F = \DE$ or $F = \CE$, and
  $D \subset \T$ is dense, then
  \[
    \ZZ^\circ_t
    = \Bigl\{ \textstyle\prod_{i=1}^n h_i(X_{t_i}) \;:\;
        n \in \N,\ t_1,\dots,t_n \in D \cap [0,t],\ h_i \in \Cb(E) \Bigr\}
  \]
  is determining.  This uses $\Bor(F) = \sigma(X_t : t \in \T)$
  (Fact~\ref{fact:fdd}) and the monotone class theorem.
\end{example}

\subsection{The Markovian martingale problem}

\begin{definition}[(T0) + (E0) + clock]
  \label{def:markovMP}
  Fix a clock $q$ on $\T$ and a convention $\iota \in \{\mathrm{o},\mathrm{p}\}$
  (Definition~\ref{def:clock}).  Let $A \subset \Meas(E) \times \Meas(E)$ and let
  $X$ be a jointly measurable $E$-valued process on $(\Omega,\Filt,P)$.  Put
  \begin{equation}
    \label{eq:XXA}
    \XX^\iota_A(X) \;=\;
      \Bigl\{ \bigl( f(X(t))
        - \int_{\langle 0,t \rangle_\iota} g(X(s)) \, q(\dif s) \bigr)_{t \in \T}
              \;:\; (f,g) \in A \Bigr\} ,
  \end{equation}
  and write $\XX_A(X)$ when the convention is fixed by the context or immaterial.
  $X$ is a \emph{solution of the martingale problem for $A$} if every element of
  $\XX_A(X)$ is a martingale with respect to $({}^{*}\Filt^X_t)_{t \in \T}$ of
  \eqref{eq:starfilt}.  If $(\Gilt_t)$ is a filtration with
  $\Gilt_t \supset {}^{*}\Filt^X_t$ for all $t$ and every element of $\XX_A(X)$
  is a $(\Gilt_t)$-martingale, $X$ is a \emph{solution of the martingale problem
  for $A$ with respect to $(\Gilt_t)$}.  Given $\mu \in \Prob(E)$, $X$ solves the
  martingale problem for $(A,\mu)$ if in addition $P X(0)^{-1} = \mu$.
  Replacing ``martingale'' by ``local martingale'' throughout defines the
  \emph{local martingale problem} for $A$.

  A probability measure $P \in \Prob(\DE)$ is called a solution of the
  \emph{$\DE$ martingale problem} for $A$, respectively for $(A,\mu)$, if the
  coordinate process $X(t,\omega) = \omega(t)$ on the space
  $(\DE, \Bor(\DE), P)$ is one.
\end{definition}

Definition~\ref{def:markovMP} is the special case $\XX = \XX_A(X)$ of
Definition~\ref{def:absMP}.

\begin{remark}[Provenance of Definition~\ref{def:absMP}]
  \label{rem:absMPprov}
  The abstract formulation is older than \CPS.  \JS, Definition~III.1.3 already
  attaches a martingale problem to a family $\mathcal{X}$ of optional processes
  on a filtered space carrying no measure, and does so with \emph{local}
  martingales as the default rather than as a variant.  Two features of their
  formulation are worth importing.  First, the initial condition is a measure
  $P_{\mathcal{H}}$ on an \emph{initial $\sigma$-field} $\mathcal{H} \subset \Filt_0$,
  not merely an initial distribution $\mu = P\pi_0^{-1}$; everything below goes
  through with $\sigma(\pi_0)$ replaced by $\mathcal{H}$, and
  Lemma~\ref{lem:disint} in particular.  Second, \JS, Remark~III.1.6 notes
  explicitly that the elements of $\mathcal{X}$ need be neither c\`adl\`ag nor
  even adapted --- which is the observation that Setting~\ref{set:absreg} needs
  and that \CPS's standing assumptions exclude.
\end{remark}

\begin{proposition}[\EK, (3.4); (T0) + (E0) + clock]
  \label{prop:fddchar}
  Let $X$ be jointly measurable, fix $\iota \in \{\mathrm{o},\mathrm{p}\}$, and let
  $(f,g) \in A$ be such that
  $Y^{f,g}_t = f(X(t)) - \int_{\langle 0,t \rangle_\iota} g(X(s))\,q(\dif s)$ is
  integrable for every $t \in \T$.  Then $Y^{f,g}$ is a $({}^{*}\Filt^X_t)$-martingale if and only if
  \begin{equation}
    \label{eq:fdd}
    E\Bigl[ \Bigl( f(X(t)) - f(X(s))
        - \int_{\langle s,t \rangle_\iota} g(X(u)) \, q(\dif u) \Bigr)
        \prod_{k=1}^n h_k(X(t_k)) \Bigr] \;=\; 0
  \end{equation}
  for all $s \leq t$ in $\T$, all $n \geq 0$, all $t_1,\dots,t_n \in \T_{\leq s}$
  and all $h_1,\dots,h_n \in \Bdd(E)$ (equivalently, $h_k \in \Cb(E)$).
  Consequently $X$ solves the martingale problem for $A$ if and only if
  \eqref{eq:fdd} holds for every $(f,g) \in A$.
\end{proposition}

\begin{remark}[The times need not form a chain]
  \label{rem:fddnochain}
  Earlier versions of this manuscript quantified \eqref{eq:fdd} over
  \emph{chains} $t_1 \leq \dots \leq t_n \leq s$, following \EK, (3.4), where
  $\T = \Rp$ is linearly ordered and the two formulations coincide.  On a
  preordered index they do not, and the chain version is too weak: the proof
  below tests against the generators of $\Filt^X_s = \sigma(X_u : u \in \T_{\leq s})$,
  and these are products over \emph{arbitrary} finite subsets of $\T_{\leq s}$.
  On $\T = \Rp^2$ the times $(1,0)$ and $(0,1)$ lie in $\T_{\leq (1,1)}$ without
  being comparable, and no chain reaches the product $h_1(X_{(1,0)})h_2(X_{(0,1)})$.
  This is the same phenomenon as in Remark~\ref{rem:chainonly}, but with the
  opposite moral: here it is the \emph{hypothesis} that has to be strengthened,
  and once it is, nothing is lost.
\end{remark}

\begin{proof}
  Throughout, $s \leq t$ are fixed, $Y = Y^{f,g}$, and we abbreviate
  \[
    I^h_u \;=\; \int_{\langle 0,u \rangle_\iota} h(X(v))\, q(\dif v) ,
    \qquad u \in \T,\ h \in \Bdd(E) ,
  \]
  so that ${}^{*}\Filt^X_s = \sigma\bigl( X(u),\, I^h_u \;:\; u \in \T_{\leq s},\,
  h \in \Bdd(E) \bigr)$ by \eqref{eq:starfilt}.  Note $|I^h_u| \leq \lVert h
  \rVert\, q(\T_{\leq u}) < \infty$, so each $I^h_u$ is \emph{bounded}, by
  Definition~\ref{def:clock}.  By \eqref{eq:clockadd},
  \begin{equation}
    \label{eq:fddincrement}
    Y_t - Y_s \;=\; f(X(t)) - f(X(s))
      - \int_{\langle s,t \rangle_\iota} g(X(u))\, q(\dif u) ,
  \end{equation}
  which is the bracket in \eqref{eq:fdd}; it lies in $L^1$ by hypothesis.

  \emph{Necessity.}  Each $\prod_k h_k(X(t_k))$ with $t_k \leq s$ is bounded and
  $\Filt^X_s \subset {}^{*}\Filt^X_s$-measurable, so the martingale property gives
  $E[(Y_t - Y_s)\prod_k h_k(X(t_k))] = 0$, which by \eqref{eq:fddincrement} is
  \eqref{eq:fdd}.

  \emph{Sufficiency.}  Assume \eqref{eq:fdd} and put
  \[
    \mathcal{H} \;=\; \bigl\{ W : \Omega \to \R \ \text{bounded},\
      {}^{*}\Filt^X_s\text{-measurable},\ E[(Y_t - Y_s) W] = 0 \bigr\} .
  \]
  We must show that $\mathcal{H}$ contains \emph{all} such $W$, for then
  $E[Y_t \mid {}^{*}\Filt^X_s] = Y_s$.

  \emph{Step 1: $\mathcal{H}$ is a monotone vector space.}  It is a vector space
  because $W \mapsto E[(Y_t-Y_s)W]$ is linear, and it contains the constants by
  \eqref{eq:fdd} with $n = 0$.  If $W_m \in \mathcal{H}$, $0 \leq W_m \uparrow W$
  with $W$ bounded, then $(Y_t-Y_s)W_m \to (Y_t-Y_s)W$ pointwise and
  $|(Y_t-Y_s)W_m| \leq \lVert W \rVert\, |Y_t - Y_s| \in L^1$, so dominated
  convergence gives $W \in \mathcal{H}$.

  \emph{Step 2: a multiplicative generating class.}  Let $\mathcal{K}$ be the set
  of finite products
  \begin{equation}
    \label{eq:multclass}
    Z \;=\; \prod_{k=1}^n h_k(X(t_k)) \cdot \prod_{j=1}^m I^{g_j}_{u_j} ,
    \qquad n,m \geq 0,\ t_k, u_j \in \T_{\leq s},\ h_k, g_j \in \Bdd(E) .
  \end{equation}
  Every such $Z$ is bounded, $\mathcal{K}$ is closed under multiplication, and
  $\sigma(\mathcal{K}) = {}^{*}\Filt^X_s$ by \eqref{eq:starfilt}.

  \emph{Step 3: $\mathcal{K} \subset \mathcal{H}$.}  This is the only step with
  content, and it is Fubini rather than an approximation.  Fix $Z$ as in
  \eqref{eq:multclass} and write the product of the $m$ integrals as one integral
  over the product:
  \[
    \prod_{j=1}^m I^{g_j}_{u_j}
      \;=\; \int_{(0,u_1] \times \dots \times (0,u_m]}
              \prod_{j=1}^m g_j(X(v_j)) \; q^{\otimes m}(\dif v) ,
    \qquad v = (v_1,\dots,v_m) .
  \]
  The map $(v,\omega) \mapsto (Y_t-Y_s)(\omega) \prod_k h_k(X(t_k,\omega))
  \prod_j g_j(X(v_j,\omega))$ is $\mathcal{T}^{\otimes m} \otimes \Filt$-measurable,
  because $X$ is jointly measurable, and it is dominated by
  $\lVert h \rVert \lVert g \rVert \, |Y_t - Y_s|$, which is integrable for
  $P \otimes q^{\otimes m}$ since $Y_t - Y_s \in L^1(P)$ and
  $q^{\otimes m}\bigl( (0,u_1] \times \dots \times (0,u_m] \bigr)
   \leq \prod_j q(\T_{\leq u_j}) < \infty$.  Fubini therefore applies and gives
  \[
    E\bigl[ (Y_t - Y_s) Z \bigr]
      \;=\; \int_{(0,u_1] \times \dots \times (0,u_m]}
        E\Bigl[ (Y_t - Y_s) \prod_{k=1}^n h_k(X(t_k))
          \prod_{j=1}^m g_j(X(v_j)) \Bigr] \, q^{\otimes m}(\dif v) .
  \]
  For each $v$ in the domain of integration the times
  $t_1,\dots,t_n,v_1,\dots,v_m$ all lie in $\T_{\leq s}$, and the functions
  $h_1,\dots,h_n,g_1,\dots,g_m$ are bounded; so the inner expectation vanishes by
  \eqref{eq:fdd}, applied with $n+m$ in place of $n$.  Hence
  $E[(Y_t-Y_s)Z] = 0$.  \emph{This is the step for which the times in
  \eqref{eq:fdd} must be allowed to be an arbitrary finite subset of $\T_{\leq s}$:
  the $v_j$ are produced by the integration and cannot be ordered into a chain
  with the $t_k$.}

  \emph{Step 4: conclusion.}  By Steps 1--3 and the functional monotone class
  theorem (Fact~\ref{fact:monotoneclass}), $\mathcal{H}$ contains every bounded
  $\sigma(\mathcal{K})$-measurable function, i.e.\ every bounded
  ${}^{*}\Filt^X_s$-measurable function.  Since $s \leq t$ were arbitrary,
  $Y$ is a $({}^{*}\Filt^X_t)$-martingale.

  \emph{The variant with $h_k \in \Cb(E)$.}  On a metric space $C_b(E)$ generates
  $\Bor(E)$, so the class \eqref{eq:multclass} built from $\Cb(E)$ has the same
  generated $\sigma$-field, and Steps 1--4 are unchanged.  Passing from
  $h_k \in \Cb(E)$ back to $h_k \in \Bdd(E)$ in \eqref{eq:fdd} is again Step 1
  applied one coordinate at a time.

  \emph{What was used.}  Only the preorder, through
  $\T_{\leq u} \subset \T_{\leq s}$ for $u \leq s$ and through the additivity
  \eqref{eq:clockadd}; the finiteness $q(\T_{\leq u}) < \infty$, twice, for the
  boundedness of $I^h_u$ and for Fubini; and the joint measurability of $X$.  In
  particular no topology, no linear order, no property of $q$ beyond
  Definition~\ref{def:clock}, \emph{and no choice of $\iota$}: the argument uses
  only that $\langle\cdot,\cdot\rangle_\iota$ is additive, which
  \eqref{eq:clockadd} gives for both conventions.
\end{proof}

\begin{remark}
  \label{rem:fddconsequences}
  Proposition~\ref{prop:fddchar} says that being a solution of the martingale
  problem for $A$ is a statement about the \emph{finite-dimensional
  distributions} of $X$ only.  Three consequences are used repeatedly:
  \begin{enumerate}[label=(\alph*)]
  \item any measurable modification of a solution is a solution;
  \item a solution for $A$ is a solution for the linear span $A_S$, and for the
    bounded pointwise closure of $A_S$;
  \item if $A^{(1)} \subset A^{(2)}$, every solution for $A^{(2)}$ is one for
    $A^{(1)}$.  In particular, if the bp-closures of $A^{(1)}_S$ and $A^{(2)}_S$
    agree, the two martingale problems have the same solutions
    (\EK, Proposition 4.3.1).
  \end{enumerate}
\end{remark}

\begin{definition}[\EK, Section 4.4; (T0) + (E0) + clock]
  \label{def:wellposed}
  \emph{Uniqueness} holds for the martingale problem for $(A,\mu)$ if any two
  solutions have the same finite-dimensional distributions.  The martingale
  problem for $(A,\mu)$ is \emph{well-posed} if a solution exists and uniqueness
  holds; the martingale problem for $A$ is well-posed if it is well-posed for
  every $\mu \in \Prob(E)$.  It is \emph{well-posed in $\DE$} (resp.\ $\CE$) if
  there is a unique $P \in \Prob(\DE)$ (resp.\ $\Prob(\CE)$) solving the $\DE$
  (resp.\ $\CE$) martingale problem for $(A,\mu)$.  Well-posedness in $\DE$ does
  not imply well-posedness; Theorem~\ref{thm:cadlag} shows that this gap is rare.
\end{definition}

The following two results connect Definition~\ref{def:markovMP} with the
operator-theoretic notions of Section~\ref{ssec:ch1}.

\begin{lemma}[\EK, Lemma 4.3.2; (T3)]
  \label{lem:resolvent}
  Let $X$ be a measurable process, $\Gilt_t \supset {}^{*}\Filt^X_t$, and
  $f,g \in \Bdd(E)$.  Then for fixed $\lambda \in \R$ the process
  $f(X(t)) - \int_0^t g(X(s)) \dif s$ is a $(\Gilt_t)$-martingale if and only if
  \begin{equation}
    \label{eq:resolventmg}
    e^{-\lambda t} f(X(t))
      + \int_0^t e^{-\lambda u} \bigl( \lambda f(X(u)) - g(X(u)) \bigr) \dif u
  \end{equation}
  is a $(\Gilt_t)$-martingale.  Consequently, if $(f,g) \in A$, $\lambda > 0$ and
  $X$ solves the martingale problem for $A$ with respect to $(\Gilt_t)$, then
  \begin{equation}
    \label{eq:resolventrep}
    f(X(t)) \;=\; E\Bigl[ \int_0^\infty e^{-\lambda s}
      \bigl( \lambda f(X(t+s)) - g(X(t+s)) \bigr) \dif s \Bigm| \Gilt_t \Bigr] .
  \end{equation}
\end{lemma}

\begin{proof}
  Write $Y_t = f(X(t)) - \int_0^t g(X(s))\dif s$ and let $Z_t$ denote
  \eqref{eq:resolventmg}.

  \emph{Step 1: an identity.}  We claim
  \begin{equation}
    \label{eq:ZfromY}
    Z_t \;=\; e^{-\lambda t} Y_t + \lambda \int_0^t e^{-\lambda u} Y_u \dif u .
  \end{equation}
  Indeed, substituting $Y_u = f(X(u)) - \int_0^u g(X(v))\dif v$ into the
  right-hand side and using Fubini,
  \[
    \lambda \int_0^t e^{-\lambda u} \! \int_0^u \! g(X(v)) \dif v \dif u
      = \lambda \int_0^t \! g(X(v)) \! \int_v^t \! e^{-\lambda u} \dif u \dif v
      = \int_0^t g(X(v)) \bigl( e^{-\lambda v} - e^{-\lambda t} \bigr) \dif v ,
  \]
  so the two terms involving $g$ combine to
  $- e^{-\lambda t}\int_0^t g \dif v - \int_0^t e^{-\lambda v} g \dif v
    + e^{-\lambda t}\int_0^t g\dif v = - \int_0^t e^{-\lambda v} g(X(v))\dif v$,
  and \eqref{eq:ZfromY} is \eqref{eq:resolventmg}.  Solving \eqref{eq:ZfromY} for
  $Y$ --- put $W_t = \int_0^t e^{-\lambda u}Y_u\dif u$, so that $Z = W' + \lambda W$
  and hence $(e^{\lambda t}W)' = e^{\lambda t}Z_t$ --- gives the inverse relation
  \begin{equation}
    \label{eq:YfromZ}
    Y_t \;=\; e^{\lambda t} Z_t - \lambda \int_0^t e^{\lambda u} Z_u \dif u .
  \end{equation}

  \emph{Step 2: either relation transports the martingale property.}
  Suppose $Y$ is a $(\Gilt_t)$-martingale.  For $s \leq t$ we use
  \eqref{eq:ZfromY} together with Fubini for conditional expectations, which is
  legitimate because $\sup_{u\leq t}E|Y_u| < \infty$:
  \begin{align*}
    E[Z_t \mid \Gilt_s]
      &= e^{-\lambda t} Y_s + \lambda \int_0^s e^{-\lambda u} Y_u \dif u
         + \lambda \int_s^t e^{-\lambda u} Y_s \dif u \\
      &= e^{-\lambda t} Y_s + \lambda \int_0^s e^{-\lambda u} Y_u \dif u
         + Y_s \bigl( e^{-\lambda s} - e^{-\lambda t} \bigr)
       = Z_s .
  \end{align*}
  The converse is the same computation applied to \eqref{eq:YfromZ}, with
  $e^{\lambda u}$ in place of $e^{-\lambda u}$ and $-\lambda$ in place of
  $\lambda$.  This proves the equivalence.

  \emph{Step 3: the resolvent representation.}
  Let $(f,g) \in A$ with $f,g$ bounded and $\lambda > 0$, and let $X$ solve the
  martingale problem for $A$ with respect to $(\Gilt_t)$, so that $Z$ is a
  martingale by Step~2.  For $r > 0$,
  \[
    Z_t \;=\; E[Z_{t+r} \mid \Gilt_t]
      \;=\; E\Bigl[ e^{-\lambda(t+r)} f(X(t+r))
        + \int_0^{t+r} e^{-\lambda u}(\lambda f - g)(X(u))\dif u \Bigm| \Gilt_t \Bigr] .
  \]
  As $r \to \infty$ the first term is bounded by $e^{-\lambda(t+r)}\lVert f\rVert
  \to 0$ and the second increases to an integral dominated by
  $\lambda^{-1}\lVert \lambda f - g\rVert$, so dominated convergence for
  conditional expectations gives
  \[
    Z_t \;=\; E\Bigl[ \int_0^\infty e^{-\lambda u}(\lambda f - g)(X(u))\dif u
      \Bigm| \Gilt_t \Bigr] .
  \]
  Subtracting $\int_0^t e^{-\lambda u}(\lambda f-g)(X(u))\dif u$, which is
  $\Gilt_t$-measurable, from both sides leaves
  $e^{-\lambda t}f(X(t)) = E[\int_t^\infty e^{-\lambda u}(\lambda f-g)(X(u))\dif u
  \mid \Gilt_t]$, and the substitution $u = t+s$ gives
  \eqref{eq:resolventrep}.
\end{proof}

\begin{remark}
  \label{rem:resolventscope}
  Step~1 is the only place where $\T = \Rp$ and $q = $ Lebesgue measure are used,
  and they are used twice: for the exponential $e^{-\lambda t}$, which is the
  solution of $\phi' = -\lambda\phi$, and for the Fubini rearrangement.  On
  $\T = \N_0$ with $q$ counting measure the analogue holds with $e^{-\lambda}$
  replaced by $(1+\lambda)^{-1}$ and the integrals by sums; we do not need it,
  and record only that Lemma~\ref{lem:resolvent} is (T3) for a reason of calculus,
  not of probability.
\end{remark}

\begin{proposition}[\EK, Proposition 4.3.5; (T3)]
  \label{prop:dissipative}
  Let $A$ be a linear subset of $\Bdd(E) \times \Bdd(E)$.  If for each
  $x \in E$ there exists a solution $X_x$ of the martingale problem for
  $(A,\delta_x)$, then $A$ is dissipative in the sense of
  Definition~\ref{def:dissipative}.
\end{proposition}

\begin{proof}
  Given $(f,g) \in A$ and $\lambda > 0$, \eqref{eq:resolventrep} with
  $t = 0$ gives
  $f(x) = E[\int_0^\infty e^{-\lambda s}(\lambda f(X_x(s)) - g(X_x(s)))\dif s]$,
  whence
  $|f(x)| \leq \int_0^\infty e^{-\lambda s} \lVert \lambda f - g \rVert \dif s
   = \lambda^{-1} \lVert \lambda f - g \rVert$ for every $x \in E$.
\end{proof}

\begin{remark}
  \label{rem:fullgenerator}
  Only dissipative operators arise as generators of Markov processes.  Indeed, if
  $X$ is a Markov process with measurable transition semigroup $\{T(t)\}$ on
  $\Bdd(E)$ and full generator $\hat{A}$ as in
  Definition~\ref{def:fullgenerator}, then $X$ solves the martingale problem for
  $\hat{A}$ (\EK, Proposition 4.1.7), and $\hat{A}$ is dissipative by
  Fact~\ref{fact:fullgenerator}.  Proposition~\ref{prop:dissipative} is the
  converse direction.
\end{remark}

\begin{definition}[(T2b) + (E2)]
  \label{def:cc}
  Let $D \subset \T$ be the countable dense set of (T2b).  A process $X$
  satisfies the \emph{compact containment condition} if for all
  $\varepsilon > 0$ and $T \in \T$ there is a compact
  $K_{\varepsilon,T} \subset E$ with
  \begin{equation}
    \label{eq:cc}
    P\bigl\{ X(t) \in K_{\varepsilon,T} \text{ for all } t \in \T_{\leq T} \cap D \bigr\}
      \;>\; 1 - \varepsilon .
  \end{equation}
  A family $\{X_n\}$ satisfies the compact containment condition if
  \eqref{eq:ccn} holds.
\end{definition}

% =====================================================================
\section{Sample path regularity: existence of a c\`adl\`ag modification}
\label{sec:cadlag}
% =====================================================================

\subsection{An abstract regularization theorem}
\label{ssec:absreg}

\EK{} state Theorem~\ref{thm:cadlag} below for the Markovian martingale problem,
that is, for the family $\XX_A(X)$ of \eqref{eq:XXA}.  Their proof, however, never
uses the operator $A$, nor the special filtration \eqref{eq:starfilt}, nor the
particular form $\int_0^t g(X_s)\dif s$ of the compensator.  What it uses is only
that $f(X)$ splits, for every $f$ in a sufficiently rich class, into a martingale
plus a process that is regular in time.  Isolating that hypothesis gives a
statement in the abstract setting of Definition~\ref{def:absMP}, from which
\EK, Theorem~4.3.6 follows in one line.

\CPS{} contain no counterpart of this result, and cannot: their
Setting~\ref{set:abstract} \emph{presupposes} a Polish path space $F$ and right
continuous test processes, which is exactly what a regularization theorem has to
produce.  Theorem~\ref{thm:absreg} is therefore the statement that carries a
solution \emph{into} the \CPS{} setting, with $F = \DE$; it is the missing first
step of their programme rather than a consequence of it.

\begin{setting}[(T2b) + (E2) + clock, $\iota = \mathrm{o}$]
  \label{set:absreg}
  As in Setting~\ref{set:abstract}, but with no path space and no path
  regularity assumed.  Let $\T$ satisfy (T2b) with countable dense $D \subset \T$,
  let $(\Omega,\Filt,\mathbb{F},P)$ be a filtered probability space, let $E$
  satisfy (E2), let $\XX$ be a family of real-valued $\mathbb{F}$-adapted processes
  with $P \in \Msol(\XX)$, and let $X$ be an $E$-valued, $\mathbb{F}$-adapted,
  jointly measurable process.  Neither the elements of $\XX$ nor $X$ are assumed
  to be right continuous.
\end{setting}

\begin{definition}[(T2b) + (E2)]
  \label{def:regclass}
  A set $\Phi \subset \Cb(E)$ is a \emph{regularizing class} for $(X,\XX)$ if for
  every $f \in \Phi$ there exist $Y^f \in \XX$ and an adapted real-valued process
  $C^f$ such that
  \begin{enumerate}[label=(R\arabic*)]
  \item \label{it:R1}
    $f(X(t)) = Y^f(t) + C^f(t)$ $P$-a.s., for every $t \in \T$;
  \item \label{it:R2}
    $P$-a.s., $t \mapsto C^f(t)$ has a limit along $D$ from the right at every
    $t \in \T$ and from the left at every $t > 0$;
  \item \label{it:R3}
    $C^f$ is right continuous in $L^1$, i.e.\
    $E\bigl[ |C^f(s) - C^f(t)| \bigr] \to 0$ as $s \downarrow t$, for every
    $t \in \T$.
  \end{enumerate}
\end{definition}

\begin{theorem}[Abstract c\`adl\`ag modification; (T2b) + (E2)]
  \label{thm:absreg}
  In Setting~\ref{set:absreg}, let $\Phi$ be a regularizing class for
  $(X,\XX)$ such that
  \begin{enumerate}[label=(\alph*)]
  \item $\Phi$ contains a countable subset separating the points of $E$, and
  \item $\Phi$ is separating (Definition~\ref{def:separating}),
  \end{enumerate}
  and suppose that $X$ satisfies the compact containment condition \eqref{eq:cc}.
  Then $X$ has a modification with sample paths in $\DE$.
\end{theorem}

\begin{proof}
  Let $\{f_i\} \subset \Phi$ be countable and separating the points of $E$, and
  write $Y^i = Y^{f_i}$, $C^i = C^{f_i}$ for the associated data of
  Definition~\ref{def:regclass}.

  \emph{Step 1: regularity along the rationals.}
  Each $Y^i$ is an $(\mathbb{F},P)$-martingale, so
  Fact~\ref{fact:submgreg}, applied with the countable dense set $D$, provides a set of
  full measure on which the limits of $Y^i$ through the rationals from above
  (and, for $t>0$, from below) exist at every $t$, for every $i$.  Intersecting
  it with the sets of full measure furnished by \ref{it:R2} and \ref{it:R1}
  yields $\Omega' \subset \Omega$ with $P(\Omega')=1$ such that for
  $\omega \in \Omega'$, every $i$ and every $t \in \T$ the one-sided limits of
  \begin{equation}
    \label{eq:absreg1}
    f_i(X(t)) \;=\; Y^i(t) + C^i(t)
  \end{equation}
  along $D$ exist --- as a sum of two summands each of which has them.  This is
  the only place where \ref{it:R2} is used, and it is weaker than continuity of
  $C^i$: the two summands are allowed to jump, as long as each does so in a
  c\`adl\`ag manner.

  \emph{Step 2: compactness.}
  By \eqref{eq:cc} there is $\Omega'' \subset \Omega'$ with $P(\Omega'') = 1$
  such that $\{X(t,\omega) : t \in \T_{\leq T} \cap D\}$ has compact closure for
  every $T$ in a countable cofinal subset of $\T$ and every $\omega \in \Omega''$.
  (Take the intersection over $\varepsilon = 1/m$ and over those $T$ of the events
  in \eqref{eq:cc}.)

  \emph{Step 3: existence of one-sided limits.}
  Fix $\omega \in \Omega''$ and $t \in \T$, and let $s_n \in D$ with
  $s_n \downarrow t$, which exist because $\T$ has no isolated points from the
  right.  By Step~2 the sequence $(X(s_n,\omega))_n$ lies in a
  compact set, so it has convergent subsequences.  If $x$ and $x'$ are limits
  along two subsequences, then by Step~1 and the continuity of $f_i$,
  \[
    f_i(x) \;=\; \lim_{s \to t+,\, s \in D} f_i(X(s,\omega)) \;=\; f_i(x')
    \qquad \text{for all } i ,
  \]
  and since $\{f_i\}$ separates points, $x = x'$.  Hence
  \begin{equation}
    \label{eq:absreg2}
    Y(t,\omega) \;\coloneqq\; \lim_{s \to t+,\, s \in D} X(s,\omega)
  \end{equation}
  exists for all $t \in \T$, and similarly
  $Y^-(t,\omega) = \lim_{s \to t-,\, s \in D} X(s,\omega)$ exists for $t > 0$.
  By Fact~\ref{fact:cadlagext}, $Y(\cdot,\omega) \in \DE$.  Set
  $Y(\cdot,\omega) \equiv x_0$ for $\omega \notin \Omega''$ and some fixed
  $x_0 \in E$.

  \emph{Step 4: $Y$ is a modification of $X$.}
  Let $f \in \Phi$ and $t \in \T$.  For $s > t$, \ref{it:R1} and the martingale
  property of $Y^f$ give
  \[
    E\bigl[ f(X(s)) \bigm| \Filt_t \bigr]
      \;=\; E\bigl[ Y^f(s) \bigm| \Filt_t \bigr] + E\bigl[ C^f(s) \bigm| \Filt_t \bigr]
      \;=\; Y^f(t) + E\bigl[ C^f(s) \bigm| \Filt_t \bigr] ,
  \]
  and by \ref{it:R3} and conditional Jensen,
  $E[C^f(s) \mid \Filt_t] \to C^f(t)$ in $L^1$ as $s \downarrow t$.  Hence, using
  \ref{it:R1} once more,
  $E[f(X(s)) \mid \Filt_t] \to Y^f(t) + C^f(t) = f(X(t))$ in $L^1$.  On the other
  hand $X(s) \to Y(t)$ a.s.\ along $D$ by \eqref{eq:absreg2}, so
  $f(X(s)) \to f(Y(t))$ a.s.\ and boundedly, whence
  $E[f(X(s)) \mid \Filt_t] \to E[f(Y(t)) \mid \Filt_t]$ in $L^1$.  Therefore
  \begin{equation}
    \label{eq:absreg3}
    E\bigl[ f(Y(t)) \bigm| \Filt_t \bigr] \;=\; f(X(t))
    \qquad \text{for all } f \in \Phi,\ t \in \T .
  \end{equation}
  Since $X(t)$ is $\Filt_t$-measurable and $\Phi$ is separating,
  Fact~\ref{fact:sepcond} yields $P\{Y(t) = X(t)\} = 1$ for every $t \in \T$.
\end{proof}

\begin{remark}[What the abstract form buys]
  \label{rem:absreggain}
  Three hypotheses of \EK, Theorem~4.3.6 disappear, and each of them is one that
  \CPS{} would want gone.
  \begin{enumerate}[label=(\roman*)]
  \item \emph{No operator.}  The compensator $C^f$ is an arbitrary adapted
    process subject to \ref{it:R2} and \ref{it:R3}.  In particular it may be
    \emph{path dependent}: the \CPS{} prototype
    $C^f(t) = \int_{(0,t]} g(s,L,X)\, q(\mathrm{d}s)$, with a control variable $L$
    and a locally finite measure $q$, satisfies \ref{it:R2} because it is of
    finite variation, and \ref{it:R3} because
    $E|C^f(s) - C^f(t)| \leq \int_{(t,s]} E|g(u,L,X)| \, q(\mathrm{d}u) \to 0$
    under the integrability condition that \CPS{} impose anyway.  A statement of
    this generality cannot be made in the language of an operator
    $A \subset \Cb(E)\times\Bdd(E)$.
  \item \emph{Atoms are harmless.}  Nothing above requires $q$ to be atomless.
    \ref{it:R2} asks for one-sided limits, not for continuity; the classical
    statement ``$t \mapsto \int_0^t g(X_s)\dif s$ is Lipschitz'' is the special
    case $q = $ Lebesgue measure.  In \ref{it:R3} an atom at $t$ contributes to
    $(0,t]$ but not to $(t,s]$, and so does not obstruct the limit --- which is
    precisely why the half-open convention $(0,t]$ is the right one here.
  \item \emph{No special filtration.}  \eqref{eq:starfilt} plays no r\^ole;
    $\mathbb{F}$ is any filtration making $X$ adapted and the $Y^f$ martingales.
    The filtration ${}^{*}\Filt^X$ is needed for
    Proposition~\ref{prop:fddchar}, not here.
  \end{enumerate}
  What is \emph{not} gained is a local version; see
  Remark~\ref{rem:cadlaglocal}, whose argument is unaffected.
\end{remark}

\subsection{The Markovian case}
\label{ssec:cadlagmarkov}

\begin{theorem}[\EK, Theorem 4.3.6; (T2b) + (E2) + clock, $\iota = \mathrm{o}$]
  \label{thm:cadlag}
  Assume (E2).  Let $A \subset \Cb(E) \times \Bdd(E)$ and suppose that
  $\dom(A)$ is separating and contains a countable subset that separates points.
  Let $X$ be a solution of the martingale problem for $A$ satisfying the compact
  containment condition \eqref{eq:cc}.  Then $X$ has a modification with sample
  paths in $\DE$.
\end{theorem}

\begin{proof}
  Apply Theorem~\ref{thm:absreg} with $\XX = \XX_A(X)$,
  $\mathbb{F} = ({}^{*}\Filt^X_t)$ and $\Phi = \dom(A)$: for $f \in \dom(A)$ pick
  $g$ with $(f,g) \in A$ and put
  $Y^f = f(X(\cdot)) - \int_{(0,\cdot]} g(X(s))\,q(\dif s) \in \XX^{\mathrm{o}}_A(X)$
  and $C^f(t) = \int_{(0,t]} g(X(s))\,q(\dif s)$ --- here, and only here, the
  convention must be $\iota = \mathrm{o}$.  Then \ref{it:R1} holds by
  definition.  Since $g$ is bounded, $C^f$ is of finite variation with
  $|C^f(t) - C^f(s)| \leq \lVert g \rVert \, q((s,t])$, so it has one-sided
  limits, which is \ref{it:R2}; and $q((t,s]) \downarrow q(\emptyset) = 0$ as
  $s \downarrow t$ by continuity from above, which is \ref{it:R3}.  For
  $q = $ Lebesgue measure this is the familiar statement that $C^f$ is Lipschitz,
  but no atomlessness is needed: an atom of $q$ at $t$ lies in $(0,t]$ and not in
  $(t,s]$, which is exactly why the convention of Remark~\ref{rem:convention} is
  the right one here.  Hypotheses (a) and (b) are the assumptions on $\dom(A)$.
\end{proof}

\begin{remark}
  \label{rem:cadlagseparating}
  The two hypotheses on $\dom(A)$ --- hypotheses (a) and (b) of
  Theorem~\ref{thm:absreg} --- play different r\^oles and should be kept apart in
  the formalization: ``contains a countable subset separating points'' is used in
  Step~3 of the proof of Theorem~\ref{thm:absreg} to identify limit points,
  ``separating'' is used in Step~4 to identify $Y(t)$ with $X(t)$.  They are
  independent: neither implies the other.
\end{remark}

\begin{remark}
  \label{rem:cadlaglocal}
  Theorem~\ref{thm:cadlag} is stated for genuine martingales and this is not an
  artefact of the proof.  For a local martingale problem the localizing sequence
  consists of stopping times, whose usefulness (via optional stopping) presupposes
  right continuity of the filtration, and typically of the paths --- which is
  exactly the conclusion being sought.  If one \emph{postulates} a localizing
  sequence of \emph{deterministic} times, or --- equivalently for the present
  purpose --- if $A$ is replaced by
  $A^{K} = \{(f,g) \in A : f,g \text{ bounded on } K\}$ for the compact sets $K$
  of \eqref{eq:cc}, then the proof goes through verbatim.  This is the honest
  local version of Theorem~\ref{thm:cadlag}.
\end{remark}

\begin{remark}[Verifying compact containment]
  \label{rem:ccverify}
  Condition \eqref{eq:cc} is the price paid for allowing non-locally-compact
  Polish $E$.  If $E$ is locally compact and separable and $\dom(A)$ is dense in
  $\hat{C}(E)$ (continuous functions vanishing at infinity) in the norm topology,
  then \eqref{eq:cc} may be dispensed with (see also Fact~\ref{fact:convdet}):
  pass to the one-point
  compactification $E^\Delta$, apply Theorem~\ref{thm:cadlag} there (where
  compact containment is trivial), and conclude that the modification has paths in
  $D_{E^\Delta}[0,\infty)$; see \EK, Corollary~4.3.7.
\end{remark}

% =====================================================================
\section{Uniqueness, the Markov property, and the strong Markov property}
\label{sec:uniqueness}
% =====================================================================

The theorem of this section is the workhorse of the whole theory: it reduces
uniqueness for a martingale problem, which is a statement about all
finite-dimensional distributions, to a statement about \emph{one-dimensional}
distributions only, and it produces the Markov property for free.  No semigroup
is assumed to exist, and $\dom(A)$ is not assumed to be separating.

As with Theorem~\ref{thm:cadlag}, the operator plays no r\^ole in the proof.
Section~\ref{ssec:absuniq} extracts what does, which turns out to be a single
lemma applied four times; Section~\ref{ssec:uniqmarkov} then recovers
\EK, Theorem~4.4.2.  \CPS{} have no counterpart of this result either --- they
prove no uniqueness theorem at all --- but, unlike the situation in
Section~\ref{sec:cadlag}, the obstruction is not structural: the statement fits
their setting once one adds the one ingredient a uniqueness theory needs and a
convergence theory does not, namely a way of \emph{restarting} the martingale
problem.

\subsection{The abstract uniqueness and Markov theorem}
\label{ssec:absuniq}

Throughout this subsection we work on the path space, as in
Setting~\ref{set:abstract}; $\T$ satisfies (T0) and $E$ satisfies (E0) unless
stated otherwise.  Let $F \subset E^{\T}$ be a path space with $\sigma$-field
$\mathcal{S} = \sigma(\pi_t : t \in \T)$,
coordinate process $\pi = (\pi_t)_{t \in \T}$ and canonical filtration
$\Filt^\circ_t = \sigma(\pi_u : u \leq t)$, and let $\XX^\circ$ be a family of
$(\Filt^\circ_t)$-adapted real-valued functionals on $F$.  Write
\[
  \Msol(\XX^\circ) \;=\; \bigl\{ P \in \Prob(F) \;:\;
    \text{every } Y^\circ \in \XX^\circ \text{ is a }
    (\Filt^\circ_t, P)\text{-martingale} \bigr\} ,
\]
and $\Msol(\XX^\circ,\mu) = \{P \in \Msol(\XX^\circ) : P \pi_0^{-1} = \mu\}$.

\begin{remark}
  \label{rem:filtirrelevant}
  Nothing is lost by using the canonical filtration.  If $X$ is a process with
  paths in $F$ solving the martingale problem with respect to \emph{some}
  filtration $(\Gilt_t)$ to which it is adapted --- for instance
  \eqref{eq:starfilt} --- then, since each $Y^\circ_t(X)$ is
  $\sigma(X_u : u \leq t)$-measurable and
  $\sigma(X_u : u \leq t) \subset \Gilt_t$, the tower property makes $Y^\circ(X)$
  a martingale for the smaller filtration as well.  Hence
  $P X^{-1} \in \Msol(\XX^\circ)$, and $\Msol(\XX^\circ)$ is the set of laws of
  all solutions, for all filtrations at once.
\end{remark}

The martingale property is an identity that is \emph{linear} in the measure.
This trivial observation is what carries the mixing steps of the classical proof,
and it deserves to be recorded once.

\begin{lemma}[Mixture stability; (T0) + (E0)]
  \label{lem:mixture}
  Let $(\Theta,\mathcal{A},\nu)$ be a probability space and
  $\vartheta \mapsto P_\vartheta \in \Msol(\XX^\circ)$ a measurable family, and
  put $P = \int P_\vartheta \, \nu(\dif\vartheta)$.  If
  $\int E^{P_\vartheta}|Y^\circ_t| \, \nu(\dif\vartheta) < \infty$ for every
  $Y^\circ \in \XX^\circ$ and $t \in \T$, then $P \in \Msol(\XX^\circ)$.  In
  particular $\Msol(\XX^\circ)$ is convex.
\end{lemma}

\begin{proof}
  For $G \in \Filt^\circ_s$ the identity
  $E^{P_\vartheta}[Y^\circ_t \one_G] = E^{P_\vartheta}[Y^\circ_s \one_G]$ holds
  for every $\vartheta$; integrate it $\nu(\dif\vartheta)$ and use Fubini, which
  is licensed by the integrability hypothesis.
\end{proof}

Its converse holds as well, and is what actually extends well-posedness from
degenerate to general initial distributions.  It is worth separating from
Theorem~\ref{thm:absuniq}, because --- unlike the argument used in
Corollary~\ref{cor:uniqviadual} --- it does not go through the Markov property.

\begin{lemma}[Disintegration; (T0) + (E1)]
  \label{lem:disint}
  Suppose $\XX^\circ$ admits a determining set $\ZZ^\circ$ and that membership in
  $\Msol(\XX^\circ)$ can be tested along countably many data, i.e.\ there are a
  countable $\XX^\circ_0 \subset \XX^\circ$, a countable $D \subset \T$ and
  countable sets $\ZZ^\circ_s$ such that $P \in \Msol(\XX^\circ)$ if and only if
  \begin{equation}
    \label{eq:countabletest}
    E^P\bigl[ (Y^\circ_t - Y^\circ_s) Z^\circ_s \bigr] = 0
    \qquad \text{for all } Y^\circ \in \XX^\circ_0,\ s \leq t \text{ in } D,\
    Z^\circ_s \in \ZZ^\circ_s .
  \end{equation}
  Let $P \in \Msol(\XX^\circ)$ with $P\pi_0^{-1} = \mu$.  Then
  $P(\,\cdot \mid \pi_0 = x) \in \Msol(\XX^\circ, \delta_x)$ for $\mu$-a.e.\ $x$.
  In particular, if $\Msol(\XX^\circ,\delta_x) = \{P_x\}$ for $\mu$-a.e.\ $x$,
  then $P = \int P_x \, \mu(\dif x)$, and $\Msol(\XX^\circ,\mu)$ is a singleton.
\end{lemma}

\begin{proof}
  Fix $Y^\circ \in \XX^\circ_0$, $s \leq t$ in $D$ and $Z^\circ_s \in \ZZ^\circ_s$.
  For bounded measurable $h : E \to \R$ the variable $Z^\circ_s\, h(\pi_0)$ is
  bounded and $\Filt^\circ_s$-measurable, so the martingale property of $Y^\circ$
  under $P$ gives $E^P[(Y^\circ_t - Y^\circ_s) Z^\circ_s h(\pi_0)] = 0$.  As $h$
  is arbitrary,
  $E^P[(Y^\circ_t - Y^\circ_s) Z^\circ_s \mid \pi_0] = 0$ $P$-a.s.  The data being
  countable, one null set serves all of them at once; off it,
  $P(\,\cdot \mid \pi_0 = x)$ satisfies \eqref{eq:countabletest} and has initial
  law $\delta_x$.
\end{proof}

\begin{remark}
  \label{rem:disintkall}
  Lemma~\ref{lem:disint} together with Lemma~\ref{lem:mixture} is
  \KA, Theorem~32.10, in abstract form; \KA{} attributes it to Stroock and
  Varadhan and uses it exactly this way, to pass from $P_x$ to
  $P_\mu = \int P_x \mu(\dif x)$ for the local martingale problem of a diffusion.
  The countability hypothesis is his too: he replaces $\hat{C}^\infty$ by a
  countable class of truncated coordinate functions and their products.  Note
  that Corollary~\ref{cor:uniqviadual} currently justifies the same passage via
  Theorem~\ref{thm:absuniq}\ref{it:absuniq_a}; Lemma~\ref{lem:disint} is the
  better route, since it needs no Markov property.
\end{remark}

The second ingredient is the one that a convergence theory does not need and a
uniqueness theory cannot do without.

\begin{definition}[Shift system; (T4) + (E0)]
  \label{def:shiftstable}
  A family of measurable maps $\theta_r : F \to F$, $r \in \T$, is a
  \emph{shift} on $F$ if $\pi_t \circ \theta_r = \pi_{r+t}$ for all $r,t \in \T$.

  A \emph{shift system} for $\XX^\circ$ is a family $(\XX^\circ_r)_{r \in \T}$ of
  families of $(\Filt^\circ_t)$-adapted real-valued processes on $F$, with
  $\XX^\circ_0 = \XX^\circ$, such that for every $r \in \T$ and every
  $\hat{Y}^\circ \in \XX^\circ_r$ there are $Y^\circ \in \XX^\circ$ and an
  $\Filt^\circ_r$-measurable $\kappa$ with
  \begin{equation}
    \label{eq:shiftstable}
    \hat{Y}^\circ_t \circ \theta_r \;=\; Y^\circ_{r+t} - Y^\circ_r + \kappa ,
    \qquad t \in \T .
  \end{equation}
  We call $\XX^\circ_r$ the \emph{problem shifted to $r$}, and say that
  $\XX^\circ$ is \emph{shift stable} in the narrow sense if
  $\XX^\circ_r = \XX^\circ$ for every $r$ is a shift system.

  The system is \emph{measurable} if there is a set $I$ and, for each $i \in I$
  and $r \in \T$, an element $\hat{Y}^{\circ,r,i} \in \XX^\circ_r$ exhausting
  $\XX^\circ_r$, such that $(\omega,r) \mapsto \hat{Y}^{\circ,r,i}_t(\omega)$ is
  $\Bor(F) \otimes \mathcal{T}$-measurable for every $i$ and $t$.  For an
  a.s.\ finite stopping time $\tau$ one then writes
  $\XX^\circ_\tau = \{ \hat{Y}^{\circ,\tau(\cdot),i} : i \in I \}$.
\end{definition}

\begin{remark}
  \label{rem:shiftsystemwhy}
  Requiring $\XX^\circ_r = \XX^\circ$ --- which is what the first version of this
  manuscript did --- is a real restriction and not a normalization: it forces the
  martingale problem to look the same from every instant, hence forces the clock
  to be shift invariant and the resulting Markov process to be
  \emph{homogeneous}.  Definition~\ref{def:shiftstable} follows
  \JS, III.2.39 instead, where a whole family of shifted triplets
  $(\rho_r B, \rho_r C, \rho_r \nu)$ is carried along and
  $(\rho_r B)_t(\theta_r\omega) = B_{r+t}(\omega) - B_r(\omega)$; the measurability
  clause is their III.2.39(i).  The cost is nil: the proofs below use only that
  the \emph{shifted} process is a $P$-martingale, never that it belongs to the
  original family.  The gain is the time-inhomogeneous case, and with it the
  removal of shift invariance of the clock from the hypotheses; see
  Remark~\ref{rem:absuniqgain}(ii) and Example~\ref{ex:shiftXA}.
\end{remark}

\begin{example}[The shifted problem for $\XX_A$; (T0) + (T4)]
  \label{ex:shiftXA}
  Let $q$ be a clock on $\T$, let $A \subset \Meas(\T \times E) \times \Meas(E)$
  allow a time-dependent coefficient, and put
  \[
    \XX^\circ \;=\; \Bigl\{ \bigl( f(\pi_t)
      - \int_{\langle 0,t \rangle_\iota} g(s,\pi_s)\, q(\dif s) \bigr)_{t \in \T}
      \;:\; (f,g) \in A \Bigr\} .
  \]
  Define the \emph{pulled-back clock} $q_r$ by $q_r(A) = q(r + A)$, i.e.\ the
  image of $q$ restricted to $\{v : v \geq r\}$ under $v \mapsto v - r$, and set
  \[
    \XX^\circ_r \;=\; \Bigl\{ \bigl( f(\pi_t)
      - \int_{\langle 0,t \rangle_\iota} g(r+u,\pi_u)\, q_r(\dif u) \bigr)_{t \in \T}
      \;:\; (f,g) \in A \Bigr\} .
  \]
  Then $(\XX^\circ_r)$ is a shift system, with $\kappa = f(\pi_r)$: substituting
  $v = r+u$,
  \begin{align*}
    \hat{Y}^\circ_t(\theta_r\omega)
      &= f(\omega(r+t)) - \int_{\langle 0,t \rangle_\iota} g(r+u,\omega(r+u))\, q_r(\dif u) \\
      &= f(\omega(r+t)) - \int_{\langle r,\,r+t \rangle_\iota} g(v,\omega(v))\, q(\dif v) ,
  \end{align*}
  which is $Y^\circ_{r+t}(\omega) - Y^\circ_r(\omega) + f(\omega(r))$.  Note
  $q_0 = q$ and $\XX^\circ_0 = \XX^\circ$.

  The shifted problem always exists.  It coincides with the original one, i.e.\
  $\XX^\circ_r = \XX^\circ$ for all $r$, precisely when $q_r = q$ and $g$ does not
  depend on time --- that is, when the clock is shift invariant in the sense of
  Definition~\ref{def:clock}.  The system is measurable as soon as
  $(r,u,x) \mapsto g(r+u,x)$ is, which for $q = $ Lebesgue measure and
  measurable $g$ is automatic.
\end{example}

\begin{lemma}[Restart; (T0) + (T4) + (E0)]
  \label{lem:restart}
  Let $(\XX^\circ_r)$ be a shift system for $\XX^\circ$ and let $\ZZ^\circ$ be a
  determining set for every $\XX^\circ_r$ (Definition~\ref{def:canonical}).  Let
  $P \in \Msol(\XX^\circ)$, let $r \in \T$ and let $Z \geq 0$ be
  $\Filt^\circ_r$-measurable with $E^P[Z] = 1$.  Then, provided all the integrals
  below exist,
  \begin{equation}
    \label{eq:restart}
    R \;\coloneqq\; (Z \cdot P) \circ \theta_r^{-1} \;\in\; \Msol(\XX^\circ_r) .
  \end{equation}
\end{lemma}

\begin{proof}
  Let $\hat{Y}^\circ \in \XX^\circ_r$, let $s < t$ and let
  $Z^\circ_s \in \ZZ^\circ_s$.  Choose $Y^\circ \in \XX^\circ$ and $\kappa$ as in
  \eqref{eq:shiftstable}.  Since $Z^\circ_s$ is $\Filt^\circ_s$-measurable and
  $\pi_u \circ \theta_r = \pi_{r+u}$, the random variable
  $Z^\circ_s \circ \theta_r$ is $\Filt^\circ_{r+s}$-measurable; so is $Z$,
  because $\Filt^\circ_r \subset \Filt^\circ_{r+s}$, and so is $\kappa$.  Hence
  $W \coloneqq Z \cdot (Z^\circ_s \circ \theta_r)$ is bounded and
  $\Filt^\circ_{r+s}$-measurable, and the martingale property of $Y^\circ$ under
  $P$ gives $E^P[(Y^\circ_{r+t} - Y^\circ_{r+s}) W] = 0$.  Therefore
  \begin{align*}
    E^R\bigl[ \hat{Y}^\circ_t Z^\circ_s \bigr]
      &= E^P\bigl[ Z \, (\hat{Y}^\circ_t \circ \theta_r)(Z^\circ_s \circ \theta_r) \bigr]
       = E^P\bigl[ (Y^\circ_{r+t} - Y^\circ_r + \kappa) W \bigr] \\
      &= E^P\bigl[ (Y^\circ_{r+s} - Y^\circ_r + \kappa) W \bigr]
       = E^R\bigl[ \hat{Y}^\circ_s Z^\circ_s \bigr] .
  \end{align*}
  As $Z^\circ_s \in \ZZ^\circ_s$ was arbitrary, part~(ii) of
  Definition~\ref{def:canonical} yields
  $E^R[\hat{Y}^\circ_t \mid \Filt^\circ_s] = \hat{Y}^\circ_s$.
\end{proof}

Note where the two hypotheses act.  Shift stability is used only to rewrite
$\hat{Y}^\circ_t \circ \theta_r$; the determining set is used only in the last
step, and it is exactly the r\^ole that Proposition~\ref{prop:fddchar} plays in
\EK.  Everything else is the tower property.

\begin{theorem}[Abstract uniqueness and Markov property; (T0) + (T4), and (T2a) for \ref{it:absuniq_b}]
  \label{thm:absuniq}
  Let $(\XX^\circ_r)$ be a shift system for $\XX^\circ$, let $\ZZ^\circ$ be
  determining for every $\XX^\circ_r$, and suppose that for every $r \in \T$
  \begin{equation}
    \label{eq:absonedim}
    P, Q \in \Msol(\XX^\circ_r), \ P \pi_0^{-1} = Q \pi_0^{-1}
    \quad \Longrightarrow \quad
    P \pi_t^{-1} = Q \pi_t^{-1} \ \text{ for all } t \in \T .
  \end{equation}
  Then, subject to the integrability provisos of Lemmas~\ref{lem:mixture} and
  \ref{lem:restart}:
  \begin{enumerate}[label=(\alph*)]
  \item \label{it:absuniq_a}
    every solution is Markov --- in general \emph{time inhomogeneously} so: if
    $X$ has paths in $F$ and solves the martingale problem with respect to
    $(\Gilt_t)$, then for $f \in \Bdd(E)$, $r,t \in \T$,
    \[
      E\bigl[ f(X(r+t)) \bigm| \Gilt_r \bigr]
        \;=\; E\bigl[ f(X(r+t)) \bigm| X(r) \bigr] ;
    \]
  \item \label{it:absuniq_b}
    if $\T$ satisfies (T2a), then $\Msol(\XX^\circ,\mu)$ contains at most one
    element, for every $\mu \in \Prob(E)$.
  \end{enumerate}
\end{theorem}

\begin{proof}
  \emph{\ref{it:absuniq_a}.}
  Fix $r \in \T$ and $F_0 \in \Gilt_r$ with $P(F_0)>0$, and put
  \[
    Z_1 = \frac{\one_{F_0}}{P(F_0)} , \qquad
    Z_2 = \frac{E[\one_{F_0} \mid X(r)]}{P(F_0)} .
  \]
  Both are nonnegative with mean $1$, and both are measurable with respect to
  $\sigma(X(u):u\leq r) \subset \Gilt_r$ --- for $Z_2$ because it is even
  $\sigma(X(r))$-measurable.  Transporting to $F$ by $X$ and applying
  Lemma~\ref{lem:restart} at $r$, the two laws
  $R_i = (Z_i \cdot P) \circ (X(r+\cdot))^{-1}$ lie in $\Msol(\XX^\circ_r)$ ---
  the \emph{same} shifted problem for $i=1$ and $i=2$, which is all that
  \eqref{eq:absonedim} needs.
  A direct computation gives, for $B \in \Bor(E)$,
  \[
    R_1 \pi_0^{-1}(B) = \frac{E[\one_{F_0} \one_B(X(r))]}{P(F_0)}
      = R_2 \pi_0^{-1}(B) ,
  \]
  the middle expression for $i=2$ because
  $E[\one_{F_0} \one_B(X(r))] = E\bigl[E[\one_{F_0}\mid X(r)] \one_B(X(r))\bigr]$.
  So \eqref{eq:absonedim} applies and yields
  $E^{R_1}[f(\pi_t)] = E^{R_2}[f(\pi_t)]$, i.e.
  \[
    E\bigl[ \one_{F_0} \, E[f(X(r+t)) \mid \Gilt_r] \bigr]
      \;=\; E\bigl[ \one_{F_0} \, E[f(X(r+t)) \mid X(r)] \bigr] ,
  \]
  where on the right we used the $\sigma(X(r))$-measurability of $Z_2$ once more.
  Since $F_0 \in \Gilt_r$ was arbitrary, the Markov property follows.

  \emph{\ref{it:absuniq_b}.}
  Let $P,Q \in \Msol(\XX^\circ,\mu)$.  We show by induction on $m$ that
  $E^P[\prod_{k \leq m} f_k(\pi_{t_k})] = E^Q[\prod_{k\leq m} f_k(\pi_{t_k})]$
  for $t_1 \leq \dots \leq t_m$ in $\T$ and $f_k \in \Bdd(E)$, $f_k > 0$; by
  \eqref{eq:pathsigma} this identifies $P$ and $Q$ --- no topology and no
  regularity are involved, only that the coordinates generate $\mathcal{S}$.  For
  $m=1$ this is
  \eqref{eq:absonedim}.  Assume it for $m = n$ and put
  \[
    Z \;=\; \frac{\prod_{k=1}^n f_k(\pi_{t_k})}
                 {E^P\bigl[\prod_{k=1}^n f_k(\pi_{t_k})\bigr]} ,
  \]
  which is nonnegative, $\Filt^\circ_{t_n}$-measurable and of $P$-mean $1$; let
  $Z'$ be its analogue under $Q$.  By Lemma~\ref{lem:restart} at $r = t_n$,
  $R = (Z\cdot P)\circ\theta_{t_n}^{-1}$ and $R' = (Z'\cdot Q)\circ\theta_{t_n}^{-1}$
  lie in $\Msol(\XX^\circ_{t_n})$ --- again the same shifted problem for both ---
  and the induction hypothesis applied to the $n$
  functions $f_1,\dots,f_{n-1}, f_n f$ gives $R \pi_0^{-1} = R' \pi_0^{-1}$.
  Hence \eqref{eq:absonedim} yields $E^R[f(\pi_t)] = E^{R'}[f(\pi_t)]$ for all
  $t$, which after clearing denominators is the assertion for $m = n+1$ with
  $t_{n+1} = t_n + t$.
\end{proof}

\begin{remark}[Why \ref{it:absuniq_b} needs a linear order]
  \label{rem:chainonly}
  Part \ref{it:absuniq_a} uses nothing about $\T$ beyond (T0) and (T4): the
  proof needs $\theta_r$, the inclusion $\Filt^\circ_r \subset \Filt^\circ_{r+s}$
  and the martingale property at $r+s \leq r+t$, all of which hold in a
  preordered monoid.  There is no topology and no linear order in it.

  Part \ref{it:absuniq_b} is different.  The induction restarts at $t_n$, and for
  that $\prod_{k \leq n} f_k(\pi_{t_k})$ must be $\Filt^\circ_{t_n}$-measurable,
  i.e.\ $t_k \leq t_n$ for every $k \leq n$: the times must form a \emph{chain}.
  On a linearly ordered $\T$ every finite set of times is a chain and nothing is
  lost.  On a merely partially ordered $\T$ the induction proves equality of the
  finite-dimensional distributions along chains only, and these do not determine
  the law: on $\T = \Rp^2$ the times $(1,0)$ and $(0,1)$ are incomparable, and no
  chain sees their joint distribution.  Replacing $t_n$ by $t_1 \vee \dots \vee
  t_n$, which (T1) would allow, does not repair this, since the next time
  $t_{n+1}$ need not dominate that join.

  So the slogan ``uniqueness of the one-dimensional distributions implies
  uniqueness'' is \emph{false} for a multiparameter index, and the obstruction is
  the induction, not the martingale problem.  This is the sharpest instance of
  the general pattern of Section~\ref{ssec:timeindex}: the Markov property is
  algebraic and survives on a poset, uniqueness is order-theoretic and does not.
\end{remark}

For the strong Markov property one restarts at a stopping time.  The only extra
input is optional sampling, which is where path regularity finally enters.

\begin{theorem}[Abstract strong Markov property; (T2b) + (T4)]
  \label{thm:absstrongmarkov}
  In the situation of Theorem~\ref{thm:absuniq}, assume in addition that
  $F \subset \DE$, that $(\theta_r)$ extends to a measurable map
  $(\omega,r) \mapsto \theta_r\omega$, that the shift system is measurable in the
  sense of Definition~\ref{def:shiftstable}, and that every
  $Y^\circ \in \XX^\circ$ has right continuous paths and uniformly integrable
  increments on bounded intervals.  Let $X$ have paths in $F$, solve the
  martingale problem with respect to $(\Gilt_t)$, and let $\tau$ be an a.s.\
  finite $(\Gilt_t)$-stopping time.  Then
  \[
    E\bigl[ f(X(\tau+t)) \bigm| \Gilt_\tau \bigr]
      \;=\; E\bigl[ f(X(\tau+t)) \bigm| X(\tau) \bigr] ,
    \qquad f \in \Bdd(E),\ t \in \T .
  \]
  If moreover there is a measurable family $(P_{x,r})_{x \in E,\, r \in \T}$ with
  $P_{x,r} \in \Msol(\XX^\circ_r,\delta_x)$, then
  \[
    E\bigl[ f(X(\tau+t)) \bigm| \Gilt_\tau \bigr] \;=\;
      \bigl( T_{\tau,\tau+t} f \bigr)(X(\tau)) ,
    \qquad
    \bigl( T_{r,r+t} f \bigr)(x) \;=\; \int f(\pi_t) \dif P_{x,r} ,
  \]
  a two-parameter family of kernels satisfying the Chapman--Kolmogorov relation
  $T_{r,s} T_{s,t} = T_{r,t}$ for $r \leq s \leq t$.
\end{theorem}

\begin{proof}
  The proof of Lemma~\ref{lem:restart} rests on the single identity
  $E^P[(Y^\circ_{r+t} - Y^\circ_{r+s})W] = 0$ for bounded
  $\Filt^\circ_{r+s}$-measurable $W$.  Optional sampling
  (Fact~\ref{fact:optsampl}) upgrades it to its version at $r = \tau$, with $W$
  bounded and $\Gilt_{\tau+s}$-measurable, and with $\XX^\circ_\tau$ in place of
  $\XX^\circ_r$ --- which is meaningful because the shift system is measurable.
  The remainder of that proof, and the proof of
  Theorem~\ref{thm:absuniq}\ref{it:absuniq_a}, is unchanged.  For the second
  assertion write $\mu_{F_0} = P\{X(\tau) \in \cdot \mid F_0\}$ and put
  \[
    R_1 = (Z_1 \cdot P) \circ (X(\tau+\cdot))^{-1}
      \ \text{ with } Z_1 = \one_{F_0}/P(F_0) ,
    \qquad
    R_2 = \int P_{x,\tau} \, \mu_{F_0}(\dif x) .
  \]
  Then $R_2 \in \Msol(\XX^\circ_\tau)$ by Lemma~\ref{lem:mixture}, and
  $R_1 \pi_0^{-1} = \mu_{F_0} = R_2 \pi_0^{-1}$, so \eqref{eq:absonedim}
  applies.  The Chapman--Kolmogorov relation follows by applying the first
  assertion twice.
\end{proof}

\begin{remark}[What the abstract form buys]
  \label{rem:absuniqgain}
  \begin{enumerate}[label=(\roman*)]
  \item \emph{One lemma, four applications.}  In \EK{} the conditioning
    construction in the proof of the Markov property, its variant in the
    induction step of the uniqueness proof, and the two constructions in parts
    \ref{it:uniq_b} and \ref{it:uniq_c} look like four separate arguments.  They
    are one: Lemma~\ref{lem:restart} with, in turn,
    \begin{align*}
      Z &= \one_{F_0} , &
      Z &= E[\one_{F_0}\mid X(r)] , \\
      Z &= \textstyle\prod_k f_k(\pi_{t_k}) , &
      Z &= \one_{F_0} \ \text{at a stopping time} ,
    \end{align*}
    all normalized to mean $1$.  The only property of $Z$ that is ever used is
    that it is nonnegative and $\Filt^\circ_r$-measurable.
  \item \emph{The shift system is the real hypothesis, and it costs nothing.}
    Definition~\ref{def:shiftstable} is where the Markov property comes from, and
    it is checkable.  For $\XX_A$ it never fails: whatever the clock $q$ and
    however $g$ depends on time, the shifted problem exists and is given
    explicitly by Example~\ref{ex:shiftXA}, with the pulled-back clock
    $q_r = q(r + \cdot)$ and the shifted coefficient $g(r+\cdot,\cdot)$.  What
    \emph{does} depend on $q$ and $g$ is whether the shifted problem is the
    \emph{same} problem: $\XX^\circ_r = \XX^\circ$ holds if and only if $q$ is
    shift invariant and $g$ is time independent, and only then is the Markov
    process homogeneous.  The first version of this manuscript required
    $\XX^\circ_r = \XX^\circ$ and therefore had to assume shift invariance of the
    clock; carrying the shifted family instead --- as \JS, III.2.39 do --- removes
    that hypothesis and yields the inhomogeneous case for free.  \EK{} never have
    to say any of this because their clock is Lebesgue measure and their $A$ does
    not depend on time.
  \item \emph{No operator, no boundedness, no path regularity.}  Parts
    \ref{it:absuniq_a} and \ref{it:absuniq_b} use neither $\Cb(E)$ nor $\DE$;
    $F$ may be any Polish path space, for instance $L^p_{\mathrm{loc}}$.  Path
    regularity enters only in Theorem~\ref{thm:absstrongmarkov}, and there only
    through optional sampling.
  \item \emph{Mixtures.}  Lemma~\ref{lem:mixture} says $\Msol(\XX^\circ)$ is
    convex and stable under measurable mixtures.  This is what makes the passage
    from $\delta_x$ to a general $\mu$ legitimate --- in \EK{} it is invoked
    silently in part~\ref{it:uniq_c} and in Corollary~\ref{cor:uniqviadual}.
  \end{enumerate}
\end{remark}

\subsection{The local martingale problem}
\label{ssec:localmp}

Everything in Section~\ref{ssec:absuniq} is stated for $\Msol(\XX^\circ)$.  The
local analogue is not a corollary, and the reason is worth stating plainly:
\emph{``$Y$ is a local martingale'' is not a condition linear in $P$}, because
the localizing sequence may depend on $P$.  Lemma~\ref{lem:mixture} and
Lemma~\ref{lem:disint} rest on that linearity, and Lemma~\ref{lem:restart} rests
on a martingale identity that a local martingale does not supply.  All three
therefore fail for $\Msol_{\mathrm{loc}}(\XX^\circ)$ without a further hypothesis.

The hypothesis that repairs them is due to Stroock and Varadhan and is used in
this form by \KA, Theorems~32.10 and 32.11: insist that the localization be
performed by stopping times that are \emph{functionals of the path}, hence the
same for every $P$, and that the family of such times be stable under the shift.

\begin{definition}[Localizing system; (T2b) + (T4)]
  \label{def:localizing}
  A \emph{localizing system} for $\XX^\circ$ is a family $\Sigma$ of
  $(\Filt^\circ_t)$-stopping times on $F$ --- functionals of the path, not of any
  measure --- such that the following hold.
  \begin{enumerate}[label=(L\arabic*)]
  \item \label{it:L1}
    \emph{Uniform localization.}  $\Sigma$ contains a sequence
    $\tau_1 \leq \tau_2 \leq \dots$ with $\tau_n \uparrow \infty$ pointwise, and
    \[
      P \in \Msol_{\mathrm{loc}}(\XX^\circ)
      \iff
      Y^{\circ,\tau_n} \coloneqq Y^\circ_{\cdot \wedge \tau_n}
      \text{ is a $P$-martingale for all } Y^\circ \in \XX^\circ,\ n \in \N .
    \]
  \item \label{it:L2}
    \emph{Shift covariance.}  $\sigma \in \Sigma$ and $r \in \T$ imply
    $r + \sigma \circ \theta_r \in \Sigma$.
  \item \label{it:L3}
    \emph{Integrable increments after a restart.}  For every
    $P \in \Msol_{\mathrm{loc}}(\XX^\circ)$, $Y^\circ \in \XX^\circ$,
    $r \in \T$ and $\sigma \in \Sigma$ with $\sigma \geq r$, the process
    $t \mapsto Y^\circ_{(r+t) \wedge \sigma} - Y^\circ_r$ is a martingale for the
    filtration $(\Filt^\circ_{r+t})_{t \in \T}$ under $P$.
  \end{enumerate}
  $\Sigma$ is \emph{strongly} shift covariant if \ref{it:L2} and \ref{it:L3} hold
  with $r$ replaced by any a.s.\ finite $(\Filt^\circ_t)$-stopping time $\tau$,
  the time $\tau + \sigma \circ \theta_\tau$ being required to be a stopping time.
\end{definition}

\begin{remark}
  \label{rem:localhyp}
  Three comments on Definition~\ref{def:localizing}.

  \emph{\ref{it:L1} is the whole point.}  It turns membership in
  $\Msol_{\mathrm{loc}}(\XX^\circ)$ into a family of identities
  $E^P[(Y^{\circ,\tau_n}_t - Y^{\circ,\tau_n}_s)\one_G] = 0$, each of which is
  linear in $P$.  The implication ``$\Leftarrow$'' is automatic because
  $\tau_n \uparrow \infty$; ``$\Rightarrow$'' is the assumption, and it says that
  \emph{one} sequence localizes every test process under every solution.

  \emph{\ref{it:L3} is not implied by \ref{it:L1}.}  Stopping a $P$-local
  martingale at $\sigma \in \Sigma$ leaves a local martingale, not a martingale:
  $Y^{\circ,\sigma}$ is unstopped before $r$, where nothing controls it.  What
  \ref{it:L3} asks is only that the \emph{increments after $r$} be integrable,
  which is the same kind of hypothesis as in
  Theorem~\ref{thm:absstrongmarkov} and is usually visible at once.

  \emph{The diffusion case.}  For $\XX^\circ = \{f(\pi) - f(\pi_0) - \int A f(\pi)\}$
  on $F = C_{\Rp,\R^d}$ take $\Sigma$ generated by the exit times
  $\tau_n = \inf\{t : |\pi_t| \geq n\}$ and their shifts.  Then \ref{it:L1} is
  \KA, Theorem~32.10, \ref{it:L2} holds because
  $r + \tau_n \circ \theta_r = \inf\{u \geq r : |\pi_u| \geq n\}$ is again an exit
  time, and \ref{it:L3} holds because between $r$ and $\sigma$ the path stays in
  $\overline{B}_n$, so that the increments are bounded by
  $2\lVert f \rVert_{\overline{B}_n} + q(\T_{\leq T})\lVert Af \rVert_{\overline{B}_n}$.
  \KA{} uses $\tau^f_n = \inf\{t : |M^f_t| \geq n\}$ instead, which serves equally
  well and has the advantage of not referring to the state space.
\end{remark}

With Definition~\ref{def:localizing} in hand the three lemmas of
Section~\ref{ssec:absuniq} all go through, and their proofs are the earlier ones
with $Y^\circ$ replaced by $Y^{\circ,\tau_n}$.  We give them in full, since the
point of the exercise is that nothing is being waved through.

\begin{lemma}[Local mixtures and disintegration; (T0) + (E1)]
  \label{lem:localmix}
  Parts \ref{it:lm_conv} and \ref{it:lm_a} need only (E0); (E1) is used in
  \ref{it:lm_b}.
  \begin{enumerate}[label=(\alph*)]
  \item \label{it:lm_conv}
    $\Msol_{\mathrm{loc}}(\XX^\circ)$ is convex.  \emph{No hypothesis is needed
    for this.}
  \item \label{it:lm_a}
    If $\Sigma$ satisfies \ref{it:L1}, then $\Msol_{\mathrm{loc}}(\XX^\circ)$ is
    stable under measurable mixtures $P = \int P_\vartheta \nu(\dif\vartheta)$
    subject to $\int E^{P_\vartheta}|Y^{\circ,\tau_n}_t| \, \nu(\dif\vartheta)
    < \infty$.
  \item \label{it:lm_b}
    If in addition membership can be tested along countably many data as in
    \eqref{eq:countabletest}, then for $P \in \Msol_{\mathrm{loc}}(\XX^\circ)$ one
    has $P(\,\cdot \mid \pi_0 = x) \in \Msol_{\mathrm{loc}}(\XX^\circ,\delta_x)$
    for $P\pi_0^{-1}$-a.e.\ $x$.
  \end{enumerate}
\end{lemma}

\begin{proof}
  \ref{it:lm_conv}  Let $P, P' \in \Msol_{\mathrm{loc}}(\XX^\circ)$,
  $Q = \alpha P + (1-\alpha)P'$, and let $Y^\circ \in \XX^\circ$ be localized by
  $(\tau_n)$ under $P$ and by $(\tau'_n)$ under $P'$.  Put
  $\sigma_n = \tau_n \wedge \tau'_n$.  Then $Y^{\circ,\sigma_n}$ is a martingale
  under $P$ and under $P'$ --- stopping a martingale at a further stopping time
  leaves a martingale --- hence under $Q$, and $\sigma_n \to \infty$ $Q$-a.s.
  because $\tau_n \to \infty$ $P$-a.s.\ and $\tau'_n \to \infty$ $P'$-a.s.  Here
  the two localizing sequences are allowed to differ; it is only for
  \ref{it:lm_a} that this ceases to be possible.

  \ref{it:lm_a}, \ref{it:lm_b}  By \ref{it:L1},
  $\Msol_{\mathrm{loc}}(\XX^\circ) = \Msol(\XX^\circ_\bullet)$ with
  $\XX^\circ_\bullet = \{Y^{\circ,\tau_n} : Y^\circ \in \XX^\circ, n \in \N\}$, a
  family of $(\Filt^\circ_t)$-adapted processes.  Now apply
  Lemma~\ref{lem:mixture} respectively Lemma~\ref{lem:disint} to
  $\XX^\circ_\bullet$; neither proof uses anything about the family beyond
  adaptedness, and the countability hypothesis for \ref{it:lm_b} is inherited
  because $n$ ranges over $\N$.
\end{proof}

\begin{remark}
  \label{rem:convexfree}
  The contrast between \ref{it:lm_conv} and \ref{it:lm_a} is the whole point of
  \ref{it:L1}, and it is easy to get wrong.  A \emph{finite} convex combination
  costs nothing, because two localizing sequences can be combined by
  $\tau_n \wedge \tau'_n$.  A mixture over a continuum cannot: one would need
  $\inf_\vartheta \tau^\vartheta_n$, which need not tend to infinity and need not
  even be measurable.  That is exactly where a localizing sequence common to all
  $P$ becomes indispensable, and it is why \KA{} states his Theorem~32.10 for a
  \emph{kernel} $x \mapsto P_x$ and takes the trouble to exhibit one.
\end{remark}

The hypothesis \ref{it:L1} is less exotic than it looks: in the two places in the
literature where it is used it is not assumed but \emph{constructed}, and always
by the same device --- localizing along the test process itself rather than along
the state space.

\begin{lemma}[Uniform localization for bounded jumps; (T2b)]
  \label{lem:L1auto}
  Suppose every $Y^\circ \in \XX^\circ$ is c\`adl\`ag with $Y^\circ_0 = 0$ and has
  jumps bounded by a constant $c_{Y^\circ}$.  Put
  \[
    \tau^{Y^\circ}_n \;=\; \inf\{ t \in \T : |Y^\circ_t| > n \} ,
    \qquad n \in \N .
  \]
  Then each $\tau^{Y^\circ}_n$ is an $(\Filt^\circ_t)$-stopping time, hence a
  functional of the path, and
  $\Sigma_0 = \{\tau^{Y^\circ}_n : Y^\circ \in \XX^\circ,\ n \in \N\}$ satisfies
  \ref{it:L1}.
\end{lemma}

\begin{proof}
  $Y^{\circ,\tau^{Y^\circ}_n}$ is bounded by $n + c_{Y^\circ}$: before the time
  the process is at most $n$ in absolute value, and it can overshoot only by a
  single jump.  A bounded local martingale is a uniformly integrable martingale,
  so for every $P \in \Msol_{\mathrm{loc}}(\XX^\circ)$ the stopped process is a
  $P$-martingale.  Conversely $\tau^{Y^\circ}_n \uparrow \infty$ because a
  c\`adl\`ag path is locally bounded.  Both directions of \ref{it:L1} follow.
\end{proof}

\begin{remark}
  \label{rem:L1sources}
  Lemma~\ref{lem:L1auto} is the argument of \JS, III.2.8, where it proves that
  the set of solutions of a semimartingale martingale problem is convex: they
  take $T_n = \inf\{t : |Y_t| > n\}$ and observe that $Y_0 = 0$ together with
  bounded $|\Delta Y|$ makes $Y^{T_n}$ bounded.  \KA{} uses the same device,
  $\tau^f_n = \inf\{t : |M^f_t| \geq n\}$, in his Theorem~32.10.  Both localize
  along the \emph{test process}, not along the state space; the exit times from
  balls in Remark~\ref{rem:localhyp} are the less canonical choice, since they
  refer to $E$ and presuppose that $f$ and $g$ are bounded on compacts.
\end{remark}

\begin{lemma}[Local restart; (T2b) + (T4)]
  \label{lem:localrestart}
  Let $(\XX^\circ_r)$ be a shift system for $\XX^\circ$ with $\ZZ^\circ$
  determining for each $\XX^\circ_r$, and let $\Sigma$ be a localizing system.
  Let $P \in \Msol_{\mathrm{loc}}(\XX^\circ)$, let $r \in \T$ and let $Z \geq 0$
  be $\Filt^\circ_r$-measurable with $E^P[Z] = 1$.  Then
  \[
    R \;\coloneqq\; (Z \cdot P) \circ \theta_r^{-1}
      \;\in\; \Msol_{\mathrm{loc}}(\XX^\circ_r) .
  \]
  The same holds with $r$ replaced by an a.s.\ finite stopping time $\tau$ and
  $\Filt^\circ_r$ by $\Filt^\circ_\tau$, provided $\Sigma$ is strongly shift
  covariant.
\end{lemma}

\begin{proof}
  By \ref{it:L1} it suffices to show that $\hat{Y}^{\circ,\tau_n}$ is an
  $R$-martingale for every $\hat{Y}^\circ \in \XX^\circ_r$ and $n \in \N$, and by
  the determining property it suffices to show, for $s \leq t$ and
  $Z^\circ_s \in \ZZ^\circ_s$,
  \begin{equation}
    \label{eq:localrestart}
    E^R\bigl[ (\hat{Y}^{\circ,\tau_n}_t - \hat{Y}^{\circ,\tau_n}_s) Z^\circ_s \bigr]
      \;=\; 0 .
  \end{equation}

  \emph{Step 1: the shifted stopped process.}
  Let $Y^\circ \in \XX^\circ$ and the $\Filt^\circ_r$-measurable $\kappa$ be as in
  \eqref{eq:shiftstable}, and put $\sigma \coloneqq r + \tau_n \circ \theta_r$,
  which lies in $\Sigma$ by \ref{it:L2} and satisfies $\sigma \geq r$.  For
  $\omega \in F$ write $u = t \wedge \tau_n(\theta_r\omega)$.  Then
  \[
    \hat{Y}^{\circ,\tau_n}_t(\theta_r\omega)
      = \hat{Y}^\circ_u(\theta_r\omega)
      = Y^\circ_{r+u}(\omega) - Y^\circ_r(\omega) + \kappa(\omega) ,
  \]
  and, $r + \cdot$ being order preserving and cancellative,
  \[
    r + u = r + \bigl( t \wedge \tau_n(\theta_r\omega) \bigr)
          = (r+t) \wedge \bigl( r + \tau_n(\theta_r\omega) \bigr)
          = (r+t) \wedge \sigma(\omega) .
  \]
  Hence
  \begin{equation}
    \label{eq:shiftedstopped}
    \hat{Y}^{\circ,\tau_n}_t \circ \theta_r
      \;=\; Y^\circ_{(r+t) \wedge \sigma} - Y^\circ_r + \kappa ,
    \qquad t \in \T .
  \end{equation}
  This is \KA's display~(17) in the proof of his Theorem~32.11, and it is the
  only computation in the local theory that is not already in
  Section~\ref{ssec:absuniq}.

  \emph{Step 2: conclusion.}
  Put $W \coloneqq Z \cdot (Z^\circ_s \circ \theta_r)$.  As in the proof of
  Lemma~\ref{lem:restart}, $W$ is bounded and $\Filt^\circ_{r+s}$-measurable:
  $Z^\circ_s$ is $\Filt^\circ_s$-measurable and $\pi_u \circ \theta_r = \pi_{r+u}$,
  while $Z$ is $\Filt^\circ_r$-measurable and $\Filt^\circ_r \subset
  \Filt^\circ_{r+s}$.  Using \eqref{eq:shiftedstopped} twice, the terms
  $-Y^\circ_r + \kappa$ cancel, and
  \[
    E^R\bigl[ (\hat{Y}^{\circ,\tau_n}_t - \hat{Y}^{\circ,\tau_n}_s) Z^\circ_s \bigr]
      = E^P\Bigl[ \bigl( Y^\circ_{(r+t) \wedge \sigma}
          - Y^\circ_{(r+s) \wedge \sigma} \bigr) W \Bigr]
      = 0
  \]
  by \ref{it:L3}, which says exactly that
  $t \mapsto Y^\circ_{(r+t)\wedge\sigma} - Y^\circ_r$ is a
  $(\Filt^\circ_{r+t})$-martingale under $P$.  This is
  \eqref{eq:localrestart}.

  For a stopping time $\tau$ the argument is unchanged, with
  $\Filt^\circ_{\tau+s}$ in place of $\Filt^\circ_{r+s}$ and with
  $\sigma = \tau + \tau_n \circ \theta_\tau$, a stopping time by hypothesis.
\end{proof}

\begin{theorem}[Local uniqueness, Markov and strong Markov property]
  \label{thm:localuniq}
  Let $(\XX^\circ_r)$ be a shift system for $\XX^\circ$ with $\ZZ^\circ$
  determining for each $\XX^\circ_r$, let $\Sigma$ be a localizing system, and
  suppose that for every $r \in \T$
  \begin{equation}
    \label{eq:localonedim}
    P, Q \in \Msol_{\mathrm{loc}}(\XX^\circ_r), \ P\pi_0^{-1} = Q\pi_0^{-1}
    \quad \Longrightarrow \quad
    P\pi_t^{-1} = Q\pi_t^{-1} \ \text{ for all } t \in \T .
  \end{equation}
  Then Theorem~\ref{thm:absuniq} holds verbatim with $\Msol$ replaced by
  $\Msol_{\mathrm{loc}}$ throughout; and if $\Sigma$ is strongly shift covariant
  and the $Y^{\circ,\tau_n}$ are right continuous with uniformly integrable
  increments on bounded intervals, so does
  Theorem~\ref{thm:absstrongmarkov}.
\end{theorem}

\begin{proof}
  The proof of Theorem~\ref{thm:absuniq} uses exactly three things:
  Lemma~\ref{lem:restart}, applied with the four densities of
  Remark~\ref{rem:absuniqgain}(i); the hypothesis \eqref{eq:absonedim}; and
  \eqref{eq:pathsigma}.  Replace the first by Lemma~\ref{lem:localrestart} and the
  second by \eqref{eq:localonedim}; the third is a statement about $F$ and is
  untouched.  Every density used is nonnegative and $\Filt^\circ_r$-measurable
  --- for the induction step, $Z = \prod_{k \leq n} f_k(\pi_{t_k})$ with
  $t_k \leq t_n$ --- so Lemma~\ref{lem:localrestart} applies to each.

  For the strong Markov property, the proof of
  Theorem~\ref{thm:absstrongmarkov} consists of the stopping time version of
  Lemma~\ref{lem:restart}, which is the second assertion of
  Lemma~\ref{lem:localrestart}, together with Lemma~\ref{lem:mixture} for the
  mixture $\int P_x \mu_{F_0}(\dif x)$; the latter is
  Lemma~\ref{lem:localmix}\ref{it:lm_a}.  Optional sampling
  (Fact~\ref{fact:optsampl}) is applied to the $Y^{\circ,\tau_n}$, which are
  genuine martingales by \ref{it:L1}.
\end{proof}

\begin{remark}[What \JS{} do differently]
  \label{rem:jsdiff}
  Three points where \JS, Chapter~III go beyond the account above.

  \emph{(i) The shifted problem --- adopted.}
  \JS, III.2.39 carry a whole family of shifted triplets
  $(\rho_t B, \rho_t C, \rho_t \nu)$ with
  $(\rho_t B)_s(\theta_t \omega) = B_{t+s}(\omega) - B_t(\omega)$ and likewise for
  $C$ and $\nu$, rather than demanding that the shifted problem be the original
  one.  Definition~\ref{def:shiftstable} is that idea in the present notation, and
  Example~\ref{ex:shiftXA} is the corresponding computation for $\XX_A$.  The
  earlier version of this manuscript required $\XX^\circ_r = \XX^\circ$ and
  consequently had to assume that the clock be shift invariant and $g$ time
  independent; both hypotheses are now gone, and the conclusion is the Markov
  property in its time-inhomogeneous form.  The change costs nothing, because
  Lemmas~\ref{lem:restart} and \ref{lem:localrestart} use only that the
  \emph{shifted} process is a $P$-martingale.

  \emph{(ii) Local uniqueness.}
  \JS, Definition~III.2.37 isolate a notion strictly stronger than the uniqueness
  of \eqref{eq:localonedim}: \emph{local uniqueness} holds if for every strict
  stopping time $T$ any two solutions of the \emph{stopped} martingale problem
  agree on $\Filt_T$.  It implies uniqueness ($T \equiv \infty$) and is what
  absolute-continuity and limit-theorem arguments actually require.  Their
  Theorem~III.2.40 shows that uniqueness \emph{does} imply local uniqueness once
  the problem is of Markovian type in the sense of (i) --- precisely, once there
  is a kernel $P_{x,t}$ solving the shifted problem.  This is carried out in
  Section~\ref{ssec:localuniq}; it became available only after (i), which is why
  it is stated there and not here.

  \emph{(iii) Strict stopping times.}
  The canonical filtration $\Filt^\circ_t = \sigma(\pi_u : u \leq t)$ is not right
  continuous, and \JS, III.2.35 therefore distinguish \emph{strict} stopping
  times --- those with $\{T \leq t\} \in \Filt^\circ_t$ for the raw filtration ---
  from stopping times for the completed right continuous filtration, with
  $\Filt^\circ_T \subset \Filt_T$ and $\Filt_{T-} \subset \Filt^\circ_T$ on
  $\{T>0\}$.  The stopping times of Definition~\ref{def:localizing} are strict, by
  construction: they must be path functionals, and that is the same requirement
  seen from the other side.  The distinction costs nothing here but should be
  made explicit in the formalization, where Mathlib's \texttt{IsStoppingTime} is
  relative to whatever filtration it is given.
\end{remark}

\begin{remark}
  \label{rem:localsummary}
  Two things are worth extracting.

  First, the entire local theory rests on the single identity
  \eqref{eq:shiftedstopped}, which says that shifting a stopped test process
  produces a test process stopped at the \emph{shifted} time.  That is
  \ref{it:L2}, and it is the local form of shift stability.  Everything else is
  either the global argument verbatim or the linearity that \ref{it:L1} restores.

  Second, this route localizes the \emph{test processes} and leaves the operator
  alone.  \EK, Section~4.6 --- and hence Remark~\ref{rem:uniquelocal} --- instead
  replaces $A$ by $A^{(m)} = \{(f\one_{K_m}, g\one_{K_m})\}$.  In the abstract
  setting the present route is the right one, because $\XX^\circ$ is primitive
  here and $A$ is not; it is also the one that \KA{} takes.

  What is \emph{not} obtained is a local version of Theorem~\ref{thm:absreg}:
  Definition~\ref{def:localizing} presupposes stopping times, which presuppose
  the path regularity that Theorem~\ref{thm:absreg} exists to produce.  See
  Remark~\ref{rem:cadlaglocal}.
\end{remark}

\subsection{Local uniqueness}
\label{ssec:localuniq}

Theorem~\ref{thm:localuniq} concludes that two solutions with the same initial law
have the same one-dimensional --- hence all --- distributions.  \JS,
Definition~III.2.37 isolate a strictly stronger property, and it is the one that
absolute-continuity and limit arguments actually use.

\begin{definition}[Local uniqueness; (T2b) + (T4)]
  \label{def:localuniq}
  Write $\XX^{\circ,T} = \{ Y^{\circ,T} : Y^\circ \in \XX^\circ \}$ for the family
  stopped at a stopping time $T$.  \emph{Local uniqueness} holds for $\XX^\circ$
  if for every strict stopping time $T$ (Remark~\ref{rem:jsdiff}(iii)) and any two
  $P, P' \in \Msol(\XX^{\circ,T})$ with $P\pi_0^{-1} = P'\pi_0^{-1}$ one has
  $P = P'$ on $\Filt^\circ_T$.
\end{definition}

Taking $T \equiv \infty$ recovers uniqueness, so local uniqueness is at least as
strong; \JS{} note that it is strictly stronger in general.  What makes it
accessible is that the shift systems of Section~\ref{ssec:absuniq} are exactly
the structure \JS{} presuppose in their Theorem~III.2.40.  Two further pieces of
structure are needed, and both are properties of the path space rather than of
the martingale problem.

\begin{definition}[Stopping and concatenation on $F$; (T2b) + (T4)]
  \label{def:pasting}
  Let $T$ be a stopping time on $F$.
  \begin{enumerate}[label=(P\arabic*)]
  \item \label{it:P1}
    \emph{$T$ is strict}: $F$ is stable under the stopping operator
    $a_T : F \to F$, $\pi_t \circ a_T = \pi_{t \wedge T}$, and
    $\Filt^\circ_T = a_T^{-1}(\mathcal{S})$.  Equivalently, a set lies in
    $\Filt^\circ_T$ exactly when membership depends on the path only through
    $a_T$.  This is \JS, Lemma~III.2.43 for the canonical space, and it is what
    ``strict'' buys.
  \item \label{it:P2}
    \emph{$F$ admits concatenation at $T$}: there is a measurable
    $\gamma : \{(\alpha,\beta) \in F \times F : \beta_0 = \alpha_{T(\alpha)}\} \to F$
    with
    \[
      \gamma(\alpha,\beta)_t =
        \begin{cases} \alpha_t, & t < T(\alpha), \\
                      \beta_{t - T(\alpha)}, & t \geq T(\alpha), \end{cases}
      \qquad
      T(\gamma(\alpha,\beta)) = T(\alpha), \qquad
      \theta_{T} \gamma(\alpha,\beta) = \beta .
    \]
  \item \label{it:P3}
    \emph{The shift system is full}: in addition to
    Definition~\ref{def:shiftstable}, for every $Y^\circ \in \XX^\circ$ and every
    $r$ there are $\hat{Y}^\circ \in \XX^\circ_r$ and an $\Filt^\circ_r$-measurable
    $\kappa$ with $\hat{Y}^\circ_u \circ \theta_r = Y^\circ_{r+u} - Y^\circ_r + \kappa$.
  \end{enumerate}
  For $F = \DE$ all three hold, with $a_T$ the usual stopping of a path and
  $\gamma$ the usual concatenation.
\end{definition}

\begin{lemma}[Pasting; (T2b) + (T4)]
  \label{lem:pasting}
  Assume \ref{it:P1}--\ref{it:P3}, let $\ZZ^\circ$ be determining for each
  $\XX^\circ_r$, and let $(P_{x,r})$ be a measurable family with
  $P_{x,r} \in \Msol(\XX^\circ_r,\delta_x)$.  Let $T$ be a strict stopping time and
  $P \in \Msol(\XX^{\circ,T})$.  Let $Q$ be the law of $\gamma(\alpha,\beta)$
  where $\alpha \sim P$ and, given $\alpha$, $\beta \sim P_{\alpha_{T(\alpha)},\,T(\alpha)}$.
  Then, subject to integrability,
  \[
    Q \in \Msol(\XX^\circ)
    \qquad\text{and}\qquad
    Q = P \ \text{ on } \Filt^\circ_T .
  \]
\end{lemma}

\begin{proof}
  \emph{Agreement on $\Filt^\circ_T$.}  Let $A \in \Filt^\circ_T$.  By \ref{it:P1},
  $\one_A = \one_A \circ a_T$, and by \ref{it:P2}, $a_T \gamma(\alpha,\beta) = a_T\alpha$.
  Hence $\one_A(\gamma(\alpha,\beta)) = \one_A(\alpha)$ and $Q(A) = P(A)$.

  \emph{The martingale property.}  Fix $Y^\circ \in \XX^\circ$, $s \leq t$, and a
  bounded $\Filt^\circ_s$-measurable $W$; we show $E^Q[(Y^\circ_t - Y^\circ_s)W] = 0$,
  which by Definition~\ref{def:canonical} suffices.  Split
  \[
    Y^\circ_t - Y^\circ_s
      = \underbrace{\bigl( Y^{\circ,T}_t - Y^{\circ,T}_s \bigr)}_{(\mathrm{I})}
      + \underbrace{\bigl( Y^\circ_t - Y^\circ_{t \wedge T} \bigr)
                    - \bigl( Y^\circ_s - Y^\circ_{s \wedge T} \bigr)}_{(\mathrm{II})} .
  \]

  \emph{The term $(\mathrm{I})$.}  It vanishes on $\{T \leq s\}$, so only
  $\{T > s\}$ contributes.  There the concatenated path agrees with $\alpha$ on
  $\T_{\leq s}$, so $W(\gamma(\alpha,\beta)) = W(\alpha)$, and $(\mathrm{I})$ depends
  on the path only through $a_T$, hence equals its value at $\alpha$.  Therefore
  $E^Q[(\mathrm{I})W] = E^P[(Y^{\circ,T}_t - Y^{\circ,T}_s)W] = 0$, because
  $Y^{\circ,T}$ is a $P$-martingale by hypothesis.

  \emph{The term $(\mathrm{II})$.}  By \ref{it:P3} choose
  $\hat{Y}^\circ \in \XX^\circ_T$ and $\kappa$ with
  $\hat{Y}^\circ_u \circ \theta_T = Y^\circ_{T+u} - Y^\circ_T + \kappa$; taking
  $u = 0$ gives $\hat{Y}^\circ_0 \circ \theta_T = \kappa$, so that
  \begin{equation}
    \label{eq:pastingII}
    (\mathrm{II}) \;=\;
      \bigl( \hat{Y}^\circ_{(t-T)^+} - \hat{Y}^\circ_{(s-T)^+} \bigr) \circ \theta_T ,
    \qquad (u-T)^+ \coloneqq u - T \ \text{on } \{T \leq u\},\ 0 \text{ otherwise},
  \end{equation}
  the constant $\kappa$ cancelling on $\{T \leq s\}$ and against
  $\hat{Y}^\circ_0$ on $\{T > s\}$.  Under $Q$, conditionally on $\Filt^\circ_T$,
  the shifted path $\theta_T$ has law $P_{\pi_T,\,T} \in \Msol(\XX^\circ_T)$, so
  $\hat{Y}^\circ$ is a martingale under it.  Now split $W$:
  \begin{itemize}
  \item On $\{T > s\}$ the variable $W\one_{\{T>s\}}$ is $\Filt^\circ_T$-measurable,
    since $\{T>s\} \in \Filt^\circ_s$ and on that event the path up to $s$ is
    determined by $a_T$.  There $(s-T)^+ = 0$, and
    $E[\hat{Y}^\circ_{(t-T)^+} - \hat{Y}^\circ_0 \mid \Filt^\circ_T] = 0$.
  \item On $\{T \leq s\}$ the variable $W\one_{\{T \leq s\}}$ is measurable with
    respect to $\Filt^\circ_T \vee \theta_T^{-1}(\Filt^\circ_{s-T})$, because the
    path up to $s$ is the concatenation of $a_T$ with the shifted path run for
    time $s-T$.  Conditionally on $\Filt^\circ_T$, the martingale property of
    $\hat{Y}^\circ$ tested against an $\Filt^\circ_{s-T}$-measurable variable gives
    $E[(\hat{Y}^\circ_{t-T} - \hat{Y}^\circ_{s-T}) \, W \mid \Filt^\circ_T] = 0$.
  \end{itemize}
  In both cases the contribution vanishes, so $E^Q[(\mathrm{II})W] = 0$.
\end{proof}

\begin{theorem}[Uniqueness implies local uniqueness; \JS, Theorem III.2.40]
  \label{thm:localuniqueness}
  Assume \ref{it:P1}--\ref{it:P3} and the hypotheses of Lemma~\ref{lem:pasting},
  and suppose that $\Msol(\XX^\circ,\mu)$ contains at most one element for every
  $\mu \in \Prob(E)$.  Then local uniqueness holds in the sense of
  Definition~\ref{def:localuniq}.
\end{theorem}

\begin{proof}
  Let $T$ be a strict stopping time and $P \in \Msol(\XX^{\circ,T})$ with
  $P\pi_0^{-1} = \mu$.  Lemma~\ref{lem:pasting} produces
  $Q \in \Msol(\XX^\circ)$ with $Q = P$ on $\Filt^\circ_T$; in particular
  $Q\pi_0^{-1} = \mu$, since $\pi_0$ is $\Filt^\circ_T$-measurable.  By
  hypothesis $Q$ is the unique element of $\Msol(\XX^\circ,\mu)$, hence depends on
  $\mu$ alone.  Therefore $P|_{\Filt^\circ_T} = Q|_{\Filt^\circ_T}$ depends on
  $\mu$ alone, and two solutions of the stopped problem with the same initial law
  agree on $\Filt^\circ_T$.
\end{proof}

\begin{remark}
  \label{rem:localuniqcomment}
  Three comments.

  \emph{Where the new structure sits.}  Everything before this subsection needs
  only the shift $\theta_r$; local uniqueness needs its partial inverse, namely
  concatenation \ref{it:P2}.  That is a genuinely new hypothesis on $F$ and not a
  consequence of anything above: $\theta_r$ forgets the past, and no amount of
  forgetting reconstructs a pasting operation.  \ref{it:P3} is the mirror image of
  Definition~\ref{def:shiftstable} and is free in every example.

  \emph{Why strict stopping times.}  \ref{it:P1} is used exactly once, to know
  that $\Filt^\circ_T$-measurable functions see only $a_T$.  For the raw
  filtration $\Filt^\circ_t = \sigma(\pi_u : u \leq t)$ this is \JS,
  Lemma~III.2.43; for a completed or right continuous filtration it is false, and
  the statement of local uniqueness would have to be adjusted.  This is the
  practical reason for \JS's insistence on the distinction.

  \emph{What it is for.}  \JS{} introduce local uniqueness for questions of
  absolute continuity and singularity between solutions, and for limit theorems;
  neither is treated here.  Theorem~\ref{thm:localuniqueness} is included because
  it is the natural completion of Section~\ref{ssec:absuniq} and because, after
  Definition~\ref{def:shiftstable}, it costs one lemma.
\end{remark}

\subsection{The Markovian case}
\label{ssec:uniqmarkov}

\begin{theorem}[\EK, Theorem 4.4.2; (T2b) + (T4) + shift invariant clock, for homogeneity]
  \label{thm:uniqueness}
  Assume (E2) and let $A \subset \Bdd(E) \times \Bdd(E)$.  Suppose that
  for each $\mu \in \Prob(E)$ any two solutions $X$, $Y$ of the martingale
  problem for $(A,\mu)$ have the same one-dimensional distributions, that is,
  for each $t \in \T$,
  \begin{equation}
    \label{eq:onedim}
    P\{X(t) \in \Gamma\} \;=\; P\{Y(t) \in \Gamma\},
    \qquad \Gamma \in \Bor(E) .
  \end{equation}
  Then the following hold.
  \begin{enumerate}[label=(\alph*)]
  \item \label{it:uniq_a}
    Any solution of the martingale problem for $A$ with respect to a filtration
    $(\Gilt_t)$ is a Markov process with respect to $(\Gilt_t)$, and any two
    solutions of the martingale problem for $(A,\mu)$ have the same
    finite-dimensional distributions.  In particular, \eqref{eq:onedim} implies
    uniqueness.
  \item \label{it:uniq_b}
    If moreover $A \subset \Cb(E) \times \Bdd(E)$, $X$ is a solution of the
    martingale problem for $A$ with respect to $(\Gilt_t)$, and $X$ has sample
    paths in $\DE$, then for every a.s.\ finite $(\Gilt_t)$-stopping time $\tau$,
    \begin{equation}
      \label{eq:strongmarkov}
      E\bigl[ f(X(\tau + t)) \bigm| \Gilt_\tau \bigr]
        \;=\; E\bigl[ f(X(\tau + t)) \bigm| X(\tau) \bigr] ,
      \qquad f \in \Bdd(E),\ t \in \T .
    \end{equation}
  \item \label{it:uniq_c}
    If, in addition to the conditions of \ref{it:uniq_b}, for each $x \in E$
    there exists a solution $P_x \in \Prob(\DE)$ of the $\DE$ martingale problem
    for $(A,\delta_x)$ such that $x \mapsto P_x(B)$ is Borel measurable for each
    $B \in \Bor(\DE)$, then, setting
    $T(t) f(x) = \int f(\omega(t)) P_x(\dif\omega)$,
    \begin{equation}
      \label{eq:strongmarkov2}
      E\bigl[ f(X(\tau + t)) \bigm| \Gilt_\tau \bigr] \;=\; T(t) f(X(\tau))
    \end{equation}
    for all $f \in \Bdd(E)$, $t \in \T$ and a.s.\ finite $(\Gilt_t)$-stopping
    times $\tau$; that is, $X$ is strong Markov.
  \end{enumerate}
\end{theorem}

\begin{proof}
  We check the hypotheses of Theorems~\ref{thm:absuniq} and
  \ref{thm:absstrongmarkov} for $\XX^\circ = \XX_A$, i.e.\ for
  \[
    \XX^\circ = \Bigl\{ \bigl( f(\pi_t)
      - \int_{\langle 0,t \rangle_\iota} g(\pi_s)\, q(\dif s)
      \bigr)_{t \in \T} \;:\; (f,g) \in A \Bigr\}
  \]
  on the path space $F = E^{\T}$ (respectively $F = \DE$ in
  \ref{it:uniq_b}--\ref{it:uniq_c}).

  \emph{Determining set.}  Proposition~\ref{prop:fddchar} says precisely that the
  products $\prod_{k} h_k(\pi_{t_k})$, over arbitrary finite subsets
  $\{t_k\} \subset \T_{\leq s}$, form a determining set for $\XX_A$ in the sense
  of Definition~\ref{def:canonical}: \eqref{eq:fdd} for all such products is
  equivalent to the martingale property.  Note that this is stronger than what
  Definition~\ref{def:canonical} asks --- it delivers the martingale property for
  ${}^{*}\Filt^X$, not merely for the canonical filtration.

  \emph{Shift system.}  Example~\ref{ex:shiftXA} exhibits one, with
  $\theta_r\omega = \omega(r+\cdot)$, the pulled-back clock $q_r$ and
  $\kappa = f(\pi_r)$.  Since $A \subset \Cb(E) \times \Bdd(E)$ does not depend on
  time and $q$ is assumed shift invariant, $q_r = q$ and $\XX^\circ_r = \XX^\circ$
  for every $r$; the shift system is the constant one, and the Markov process
  obtained is homogeneous.  Dropping either assumption leaves
  Theorem~\ref{thm:absuniq} intact and yields a time-inhomogeneous Markov process
  instead --- this is the only point at which \EK, Theorem~4.4.2 is genuinely less
  general than Theorem~\ref{thm:absuniq}.

  \emph{Integrability.}  $f$ and $g$ are bounded, so every $Y^\circ \in \XX^\circ$
  satisfies
  $|Y^\circ_t| \leq \lVert f \rVert + \lVert g \rVert \, q(\T_{\leq t})$, which is
  finite by Definition~\ref{def:clock}; all the integrability provisos are
  automatic, and the increments of $Y^\circ$ on $\T_{\leq T}$ are bounded, hence
  uniformly integrable.

  \emph{Conclusion.}  Hypothesis \eqref{eq:onedim} is \eqref{eq:absonedim} by
  Remark~\ref{rem:filtirrelevant}, so
  Theorem~\ref{thm:absuniq}\ref{it:absuniq_a} gives the Markov property and
  \ref{it:absuniq_b} the equality of the finite-dimensional distributions, which
  is \ref{it:uniq_a}.  For \ref{it:uniq_b} and \ref{it:uniq_c}, $A \subset
  \Cb(E)\times\Bdd(E)$ and paths in $\DE$ make each $Y^\circ(X)$ a right
  continuous martingale with bounded increments, so
  Theorem~\ref{thm:absstrongmarkov} applies; its second assertion is
  \eqref{eq:strongmarkov2} with $T(t)f(x) = \int f(\omega(t)) P_x(\dif\omega)$.
\end{proof}

\begin{corollary}[\EK, Corollary 4.4.3; (T3)]
  \label{cor:DEuniqueness}
  Assume (E3) and let $A \subset \Bdd(E) \times \Bdd(E)$.  Suppose that for
  each $\mu \in \Prob(E)$ any two solutions $X$, $Y$ of the martingale problem
  for $(A,\mu)$ with sample paths in $\DE$ (respectively $\CE$) satisfy
  \eqref{eq:onedim}.  Then for each $\mu \in \Prob(E)$ any two such solutions
  have the same distribution on $\DE$ (respectively $\CE$).
\end{corollary}

\begin{proof}
  The processes $\tilde X$, $\tilde Y$ constructed in the proof of
  Theorem~\ref{thm:uniqueness} again have sample paths in $\DE$ (resp.\ $\CE$)
  whenever $X$ and $Y$ do, so the induction goes through and $X$, $Y$ have the
  same finite-dimensional distributions.  By Fact~\ref{fact:fdd} these determine
  the law on $\DE$ (resp.\ $\CE$).
\end{proof}

\begin{remark}
  \label{rem:uniquelocal}
  For the \emph{local} martingale problem, Theorem~\ref{thm:uniqueness} holds in
  the following localized form.  Suppose $A \subset \Meas(E) \times \Meas(E)$ and
  let $(K_m)$ be an increasing sequence of compact sets with
  $\tau_m = \inf\{t : X(t) \notin K_m \text{ or } X(t-) \notin K_m\}$ and
  $\tau_m \to \infty$ a.s.\ for every solution $X$.  If for each $\mu$ and each
  $m$ any two solutions of the local martingale problem for $(A,\mu)$ agree in
  the one-dimensional distributions of the \emph{stopped} processes
  $X(\cdot \wedge \tau_m)$, then the conclusions \ref{it:uniq_a}--\ref{it:uniq_c}
  hold.  The proof is the above one applied to the stopped martingale problem for
  $A^{(m)} = \{(f \one_{K_m}, g \one_{K_m}) : (f,g) \in A\}$, followed by
  $m \to \infty$.  This is the content of \EK, Section~4.6.  It is worth stating
  separately in the formalization, since the reduction ``local $\rightsquigarrow$
  global via stopping'' is reused in Section~\ref{sec:convergence}.

  This route modifies the operator.  Section~\ref{ssec:localmp} gives the
  alternative, which leaves $A$ alone and stops the test process instead; it is
  the one that fits the abstract setting, and the one \KA{} uses.
\end{remark}

% =====================================================================
\section{Uniqueness via duality}
\label{sec:duality}
% =====================================================================

Duality reduces the computation of $E[f(X(t),y)]$ for a process $X$ to the
computation of $E[f(x,Y(t))]$ for a second, typically simpler, process $Y$.
Combined with Theorem~\ref{thm:uniqueness} it is one of the main practical
routes to uniqueness.

Here, unlike in Sections~\ref{sec:cadlag} and \ref{sec:uniqueness}, there is
nothing to strip away on the probabilistic side: Theorem~\ref{thm:duality} below
already assumes only that $X$ and $Y$ are \emph{measurable} processes, uses no
path regularity, no Polish structure and no boundedness, and would be stated
verbatim in the \CPS{} setting.  What can be generalized is the \emph{clock}.
Section~\ref{ssec:antidiag} therefore isolates the analytic core as a statement
about a function of two time variables, valid for an arbitrary measure
$q$, and reads off exactly which $q$ make the duality relation true.  The answer
turns out to be sharp, and it is not the answer one gets in
Sections~\ref{sec:cadlag} and \ref{sec:uniqueness}.

\subsection{The anti-diagonal lemma}
\label{ssec:antidiag}

Duality is the assertion that a certain function of two time variables is
constant on anti-diagonals.  The following identity is that assertion stripped of
all analysis: it is an exact telescoping, with no hypothesis beyond the two
increment representations, and every version of the duality lemma is obtained
from it by choosing a chain.

\begin{lemma}[Chain identity; (T0) + (T4) + (E0) + clock]
  \label{lem:chain}
  Let $(\T,+,\leq)$ be an ordered commutative monoid, let $q$ be a measure on
  $\T$, and let $\Phi, \gamma_1, \gamma_2 : \T \times \T \to \R$ satisfy, for all
  $s \leq s'$ and $t \leq t'$,
  \begin{equation}
    \label{eq:incrementrep}
    \Phi(s',t) - \Phi(s,t) = \int_{[s,s')} \gamma_1(r,t) \, q(\dif r),
    \qquad
    \Phi(s,t') - \Phi(s,t) = \int_{[t,t')} \gamma_2(s,r) \, q(\dif r),
  \end{equation}
  all integrals being assumed to exist.  Then for every $t \in \T$ and every
  chain $0 = s_0 \leq s_1 \leq \dots \leq s_m = t$ in $\T$ such that $t - s_k$ is
  defined,
  \begin{equation}
    \label{eq:chain}
    \Phi(t,0) - \Phi(0,t)
      = \sum_{k=0}^{m-1} \Bigl[
          \int_{[s_k,s_{k+1})} \! \gamma_1(r,\, t - s_{k+1}) \, q(\dif r)
          \;-\;
          \int_{[t-s_{k+1},\, t-s_k)} \! \gamma_2(s_k,\, r) \, q(\dif r)
        \Bigr] .
  \end{equation}
\end{lemma}

\begin{proof}
  Telescope along the anti-diagonal:
  $\Phi(t,0) - \Phi(0,t)
    = \sum_{k} \bigl[ \Phi(s_{k+1}, t-s_{k+1}) - \Phi(s_k, t-s_k) \bigr]$,
  and split each summand through the corner $(s_k, t-s_{k+1})$:
  \[
    \bigl[ \Phi(s_{k+1},t-s_{k+1}) - \Phi(s_k,t-s_{k+1}) \bigr]
    - \bigl[ \Phi(s_k,t-s_k) - \Phi(s_k,t-s_{k+1}) \bigr] .
  \]
  Now apply the first representation in \eqref{eq:incrementrep} to the first
  bracket and the second to the second.
\end{proof}

The duality relation is the case in which the two brackets in \eqref{eq:chain}
cancel.  Under the balance condition $\gamma_1 = \gamma_2 = \gamma$ the $k$-th
summand is
\begin{equation}
  \label{eq:cancel}
  \int_{[s_k,s_{k+1})} \gamma(r,\, t-s_{k+1}) \, q(\dif r)
  \;-\;
  \int_{[t-s_{k+1},\,t-s_k)} \gamma(s_k,\, r) \, q(\dif r) ,
\end{equation}
in which the two integrals run over the block $[s_k,s_{k+1})$ and over its
reflection $[t-s_{k+1},t-s_k)$.  So the duality relation holds as soon as $q$
\emph{gives a block and its reflection the same mass, at every scale} --- and
that is precisely translation invariance.

\begin{proposition}[Which clocks admit duality; (T2a) + (T4), $\iota = \mathrm{p}$]
  \label{prop:haar}
  Assume \eqref{eq:incrementrep} with $\gamma_1 = \gamma_2 = \gamma$.
  \begin{enumerate}[label=(\alph*)]
  \item \label{it:haar_disc}
    If $\T = \N_0$ and $q$ is counting measure, then $\Phi(t,0) = \Phi(0,t)$ for
    every $t$, for every $\gamma$ whatsoever.
  \item \label{it:haar_cont}
    If $\T = \Rp$ and $q$ is Lebesgue measure, then $\Phi(t,0) = \Phi(0,t)$ for
    almost every $t$, under the integrability hypothesis \eqref{eq:calcint};
    this is Lemma~\ref{lem:calculus} below.
  \item \label{it:haar_fail}
    For a general $q$ the conclusion is false.
  \end{enumerate}
\end{proposition}

\begin{proof}
  \ref{it:haar_disc}  Take the chain $s_k = k$.  Then $[s_k,s_{k+1}) = \{k\}$
  and $[t-s_{k+1},t-s_k) = \{t-k-1\}$, so \eqref{eq:cancel} reads
  $\gamma(k,t-k-1) - \gamma(k,t-k-1) = 0$.  Every summand vanishes identically;
  no regularity of $\gamma$ is used.

  \ref{it:haar_cont}  Refine the chain; see Lemma~\ref{lem:calculus}.

  \ref{it:haar_fail}  Let $\T = \Rp$ and $q = \delta_a$ for some $a > 0$.  Then
  \eqref{eq:incrementrep} gives, for $s,t > a$,
  $\Phi(s,t) = \Phi(0,t) + \gamma(a,t) = \Phi(s,0) + \gamma(s,a)$, and hence
  \[
    \Phi(t,0) - \Phi(0,t) \;=\; \gamma(a,0) - \gamma(0,a) ,
    \qquad t > a ,
  \]
  which does not vanish for a general $\gamma$.  In \eqref{eq:cancel} the first
  integral is carried by the block containing $a$ in its \emph{first} coordinate
  and the second by the block containing $a$ in its \emph{second} coordinate;
  these are different points of the anti-diagonal, and no chain brings them
  together.
\end{proof}

\begin{remark}[Two conventions, and they disagree]
  \label{rem:conventionclash}
  Proposition~\ref{prop:haar}\ref{it:haar_disc} uses the half-open convention
  $[s,s')$ of \eqref{eq:incrementrep}, i.e.\ a compensator
  $C_t = \int_{[0,t)} g \dif q$, which is the \emph{predictable} one.  With
  $C_t = \int_{(0,t]} g \dif q$ instead, the $k$-th summand of \eqref{eq:cancel}
  becomes $\gamma(k+1,t-k-1) - \gamma(k,t-k)$, and the sum telescopes to
  $\gamma(t,0) - \gamma(0,t) \neq 0$: \emph{discrete-time duality fails}.

  This is the opposite of what Section~\ref{sec:cadlag} wants.  There the
  half-open convention must be $(0,t]$, since an atom of $q$ at $t$ must lie
  \emph{inside} the compensator up to $t$ for $E[C_s - C_t \mid \Filt_t] \to 0$
  to hold as $s \downarrow t$, and it is that limit which makes the c\`adl\`ag
  process a \emph{modification} rather than merely a regularization
  (Remark~\ref{rem:absreggain}(ii)).

  So Theorem~\ref{thm:absreg} wants $\iota = \mathrm{o}$ and
  Proposition~\ref{prop:haar} wants $\iota = \mathrm{p}$, and no single convention
  serves both.  The conflict is invisible whenever the clock is atomless --- in
  particular for $\T = \Rp$ with $q$ Lebesgue, the only case treated in \EK{} and
  \CPS{} --- which is why neither source has to choose.  A treatment that admits
  atoms does have to, and this manuscript resolves it by carrying both
  (Definition~\ref{def:clock} and Remark~\ref{rem:convention}) rather than by
  choosing.

  That this costs almost nothing is itself worth recording, and it is not an
  accident.  Sections~\ref{sec:MP} and \ref{sec:uniqueness} use the compensator
  only through the additivity \eqref{eq:clockadd}, which both conventions
  satisfy; so does Lemma~\ref{lem:chain}, which is why the analytic core of
  Section~\ref{sec:duality} is convention-agnostic and only its \emph{conclusion},
  Proposition~\ref{prop:haar}, is not.  The two conventions part company at
  exactly two theorems out of the whole manuscript, and the bundle table of
  Section~\ref{ssec:bundletable} marks them.
\end{remark}

\begin{remark}[Scope of Section~\ref{sec:duality}]
  \label{rem:dualscope}
  Proposition~\ref{prop:haar} answers, for duality, the question that
  Definition~\ref{def:shiftstable} answers for the Markov property.  Section
  \ref{sec:uniqueness} needs $q$ \emph{shift invariant}.  Section~\ref{sec:duality}
  needs more, and it is worth being precise about how much, because translation
  invariance alone is \emph{not} enough.

  What \eqref{eq:cancel} requires is that the reflection $u \mapsto t-u$ carry
  the block $[s_k,s_{k+1})$ onto the block $[t-s_{k+1},t-s_k)$ with the same
  $q$-mass.  Both blocks are differences of \emph{down-sets},
  $\T_{\leq s'} \setminus \T_{\leq s}$; but the reflection of a down-set is an
  \emph{up-set}, so the reflected first block is
  $\{u \geq t-s_{k+1}\} \setminus \{u \geq t-s_k\}$, which coincides with the
  second block only when the order is \emph{linear}.  Hence
  Section~\ref{sec:duality} needs a linearly ordered $\T$ \emph{and} a
  translation invariant $q$ on it; up to normalization that leaves Lebesgue
  measure on $\Rp$ (or $[0,T]$) and counting measure on $h\N_0$, the two cases of
  Proposition~\ref{prop:haar}\ref{it:haar_disc}--\ref{it:haar_cont}.

  In particular Lebesgue measure on $\T = \Rp^d$, $d \geq 2$, is translation
  invariant and yet does \emph{not} admit duality.  For $d = 2$, $t = (1,1)$ and
  the chain $0 \to (1,0) \to (1,1)$, the first block $[0,(1,0))$ is Lebesgue
  null while its partner $\T_{\leq(1,1)} \setminus \T_{\leq(0,1)}$ has mass $1$,
  so \eqref{eq:cancel} does not cancel.  The duality section is thus genuinely
  narrower than the rest of the manuscript --- narrower, in particular, than
  Sections~\ref{sec:MP} and \ref{sec:uniqueness}, which survive on a partially
  ordered index --- and \eqref{eq:chain} makes visible why.
\end{remark}

\subsection{The continuous-time theory}
\label{ssec:dualcont}

\begin{lemma}[\EK, Lemma 4.4.10; (T3)]
  \label{lem:calculus}
  Let $f : [0,\infty) \times [0,\infty) \to \R$ be absolutely continuous in $s$
  for each fixed $t$ and absolutely continuous in $t$ for each fixed $s$, and
  set $(f_1,f_2) = \nabla f$.  Suppose that
  \begin{equation}
    \label{eq:calcint}
    \int_0^T \! \int_0^T |f_i(s,t)| \dif s \dif t < \infty,
    \qquad i = 1,2, \quad T > 0 .
  \end{equation}
  Then for almost every $t \geq 0$,
  \begin{equation}
    \label{eq:calcconc}
    f(t,0) - f(0,t) \;=\; \int_0^t \bigl( f_1(s,t-s) - f_2(s,t-s) \bigr) \dif s .
  \end{equation}
\end{lemma}

\begin{proof}
  By \eqref{eq:calcint} and Fubini,
  \begin{align*}
    \int_0^T \!\! \int_0^t &\bigl( f_1(t-s,s) - f_2(s,t-s) \bigr) \dif s \dif t \\
    &= \int_0^T \!\! \int_0^t \! f_1(t-s,s) \dif s \dif t
       - \int_0^T \!\! \int_0^t \! f_2(s,t-s) \dif s \dif t \\
    &= \int_0^T \!\! \int_s^T \! f_1(t-s,s) \dif t \dif s
       - \int_0^T \!\! \int_s^T \! f_2(s,t-s) \dif t \dif s \\
    &= \int_0^T \bigl( f(T-s,s) - f(0,s) \bigr) \dif s
       - \int_0^T \bigl( f(s,T-s) - f(s,0) \bigr) \dif s \\
    &= \int_0^T \bigl( f(s,0) - f(0,s) \bigr) \dif s ,
  \end{align*}
  where the third equality uses absolute continuity in the first resp.\ second
  variable, and the fourth is the substitution $s \mapsto T-s$ in the first two
  integrals.  Differentiating with respect to $T$ gives \eqref{eq:calcconc} for
  a.e.\ $T$.
\end{proof}

\begin{theorem}[\EK, Theorem 4.4.11; (T3)]
  \label{thm:duality}
  Let $X$ and $Y$ be independent measurable processes with values in $E_1$ and
  $E_2$ respectively.  Let $f,g,h \in \Meas(E_1 \times E_2)$,
  $\alpha \in \Meas(E_1)$ and $\beta \in \Meas(E_2)$.  Suppose that for each
  $T > 0$ there exist an integrable random variable $\Gamma_T$ and a constant
  $C_T$ such that
  \begin{equation}
    \label{eq:dual1}
    \begin{aligned}
      \sup_{r,s,t \leq T} \bigl( |\alpha(X(r))| + 1 \bigr) |f(X(s),Y(t))| &\leq \Gamma_T , \\
      \sup_{r,s,t \leq T} \bigl( |\beta(Y(r))| + 1 \bigr) |f(X(s),Y(t))| &\leq \Gamma_T , \\
      \sup_{r,s,t \leq T} \bigl( |\alpha(X(r))| + 1 \bigr) |g(X(s),Y(t))| &\leq \Gamma_T , \\
      \sup_{r,s,t \leq T} \bigl( |\beta(Y(r))| + 1 \bigr) |h(X(s),Y(t))| &\leq \Gamma_T ,
    \end{aligned}
  \end{equation}
  and
  \begin{equation}
    \label{eq:dual2}
    \int_0^T |\alpha(X(u))| \dif u + \int_0^T |\beta(Y(u))| \dif u \;\leq\; C_T .
  \end{equation}
  Suppose that
  \begin{equation}
    \label{eq:dualmg1}
    f(X(t),y) - \int_0^t g(X(s),y) \dif s
  \end{equation}
  is a $({}^{*}\Filt^X_t)$-martingale for each $y \in E_2$, and
  \begin{equation}
    \label{eq:dualmg2}
    f(x,Y(t)) - \int_0^t h(x,Y(s)) \dif s
  \end{equation}
  is a $({}^{*}\Filt^Y_t)$-martingale for each $x \in E_1$ (the integrals in
  \eqref{eq:dualmg1} and \eqref{eq:dualmg2} being assumed to exist).  Then for
  almost every $t \geq 0$,
  \begin{multline}
    \label{eq:dualconc}
    E\Bigl[ f(X(t),Y(0)) \exp\Bigl\{ \int_0^t \alpha(X(u)) \dif u \Bigr\} \Bigr]
    - E\Bigl[ f(X(0),Y(t)) \exp\Bigl\{ \int_0^t \beta(Y(u)) \dif u \Bigr\} \Bigr] \\
    = E\Bigl[ \int_0^t \bigl\{ g(X(s),Y(t-s)) - h(X(s),Y(t-s))
        + \bigl( \alpha(X(s)) - \beta(Y(t-s)) \bigr) f(X(s),Y(t-s)) \bigr\} \\
      \times \exp\Bigl\{ \int_0^s \alpha(X(u)) \dif u
        + \int_0^{t-s} \beta(Y(u)) \dif u \Bigr\} \dif s \Bigr].
  \end{multline}
\end{theorem}

\begin{proof}
  Set
  \begin{equation}
    \label{eq:F}
    F(s,t) = E\Bigl[ f(X(s),Y(t))
      \exp\Bigl\{ \int_0^s \alpha(X(u)) \dif u \Bigr\}
      \exp\Bigl\{ \int_0^t \beta(Y(u)) \dif u \Bigr\} \Bigr] .
  \end{equation}
  \emph{Step 1: freezing $Y$.}
  Write $e_\alpha(s) = \exp\{\int_0^s \alpha(X(u))\dif u\}$ and
  $e_\beta(t) = \exp\{\int_0^t \beta(Y(u))\dif u\}$, and fix $t$.  Both $e_\beta(t)$
  and $Y(t)$ are $\sigma(Y)$-measurable, and $\sigma(Y)$ is independent of
  $\sigma(X)$; so every expectation below may be computed by conditioning on
  $\sigma(Y)$, freezing $Y(t) = y$ and $e_\beta(t)$, and using that
  $M^y_r = f(X(r),y) - \int_0^r g(X(v),y)\dif v$ is a
  $({}^{*}\Filt^X_r)$-martingale under the conditioned measure --- which it is,
  because conditioning on an independent $\sigma$-field leaves the law of $X$ and
  the filtration ${}^{*}\Filt^X$ unchanged.  We suppress the conditioning from the
  notation.

  \emph{Step 2: the increment identity.}  For $h > 0$ split
  \begin{multline}
    \label{eq:dualsplit}
    f(X(s{+}h),y) e_\alpha(s{+}h) - f(X(s),y) e_\alpha(s) \\
      = \underbrace{\bigl[ f(X(s{+}h),y) - f(X(s),y) \bigr] e_\alpha(s)}_{(A)}
      + \underbrace{f(X(s{+}h),y) \bigl[ e_\alpha(s{+}h) - e_\alpha(s) \bigr]}_{(B)} .
  \end{multline}
  For $(A)$, write $f(X(s{+}h),y) - f(X(s),y) = M^y_{s+h} - M^y_s
  + \int_s^{s+h} g(X(v),y)\dif v$.  Since $e_\alpha(s)$ is
  ${}^{*}\Filt^X_s$-measurable and bounded by $e^{C_T}$, the martingale term
  contributes $0$, so
  \begin{align*}
    E[(A)] &= \int_s^{s+h} E\bigl[ g(X(r),y) e_\alpha(s) \bigr] \dif r \\
      &= \int_s^{s+h} E\bigl[ g(X(r),y) e_\alpha(r) \bigr] \dif r \\
      &\qquad + \int_s^{s+h} E\bigl[ g(X(r),y) (e_\alpha(s) - e_\alpha(r)) \bigr] \dif r .
  \end{align*}
  For $(B)$, the chain rule gives
  $e_\alpha(s{+}h) - e_\alpha(s) = \int_s^{s+h} \alpha(X(r)) e_\alpha(r) \dif r$.
  For each $r \in [s,s+h]$ write
  \[
    f(X(s{+}h),y) \;=\; f(X(r),y) + \bigl( M^y_{s+h} - M^y_r \bigr)
      + \int_r^{s+h} g(X(v),y)\dif v .
  \]
  The martingale term again contributes $0$, since $\alpha(X(r))e_\alpha(r)$ is
  ${}^{*}\Filt^X_r$-measurable, and one obtains
  \begin{align*}
    E[(B)] &= \int_s^{s+h} E\bigl[ f(X(r),y)\alpha(X(r)) e_\alpha(r) \bigr] \dif r \\
      &\qquad + \int_s^{s+h} \! \int_r^{s+h}
          E\bigl[ g(X(v),y) \alpha(X(r)) e_\alpha(r) \bigr] \dif v \dif r .
  \end{align*}
  Multiplying by $e_\beta(t)$, integrating over the law of $Y$ and using Fubini ---
  legitimate because all integrands are dominated by $\Gamma_T e^{C_T}$ by
  \eqref{eq:dual1}--\eqref{eq:dual2} --- we obtain
  \begin{multline}
    \label{eq:Fincrement}
    F(s+h,t) - F(s,t) \;=\;
      \underbrace{\int_s^{s+h} \!\! E\bigl[ f\alpha(X(r),Y(t)) e_\alpha(r) e_\beta(t) \bigr] \dif r}_{T_1} \\
      + \underbrace{\int_s^{s+h} \!\! \int_r^{s+h} \!\!
          E\bigl[ g(X(v),Y(t)) \alpha(X(r)) e_\alpha(r) e_\beta(t) \bigr] \dif v \dif r}_{T_2} \\
      + \underbrace{\int_s^{s+h} \!\! E\bigl[ g(X(r),Y(t)) e_\alpha(r) e_\beta(t) \bigr] \dif r}_{T_3} \\
      + \underbrace{\int_s^{s+h} \!\!
          E\bigl[ g(X(r),Y(t)) \bigl( e_\alpha(s) - e_\alpha(r) \bigr) e_\beta(t) \bigr] \dif r}_{T_4} ,
  \end{multline}
  where $f\alpha$ abbreviates $(x,y) \mapsto f(x,y)\alpha(x)$.

  \emph{Step 3: $T_2$ and $T_4$ are $O(h^2)$.}  Let $t, s+h \leq T$.  The third
  line of \eqref{eq:dual1} gives
  $|\alpha(X(r))| \, |g(X(v),Y(t))| \leq \Gamma_T$, and \eqref{eq:dual2} gives
  $e_\alpha(r) e_\beta(t) \leq e^{C_T}$.  Hence
  \[
    |T_2| \;\leq\; e^{C_T} E[\Gamma_T] \int_s^{s+h} (s+h-r) \dif r
          \;=\; \tfrac{1}{2} h^2 e^{C_T} E[\Gamma_T] .
  \]
  For $T_4$ use $e_\alpha(s) - e_\alpha(r) = -\int_s^r \alpha(X(u)) e_\alpha(u)\dif u$
  and the same two bounds:
  \[
    |T_4| \;\leq\; \int_s^{s+h} \!\! \int_s^r
      E\bigl[ |g(X(r),Y(t))| \, |\alpha(X(u))| \, e_\alpha(u) e_\beta(t) \bigr] \dif u \dif r
      \;\leq\; \tfrac{1}{2} h^2 e^{C_T} E[\Gamma_T] .
  \]

  \emph{Step 4: passing to the limit.}  Since
  $T_1 + T_3 = \int_s^{s+h} E[(f\alpha + g)(X(r),Y(t)) e_\alpha(r) e_\beta(t)]\dif r$,
  summing \eqref{eq:Fincrement} over a partition
  $0 = s_0 < \dots < s_m = s$ gives
  \[
    \Bigl| F(s,t) - F(0,t)
      - \int_0^s E\bigl[ (f\alpha + g)(X(r),Y(t)) e_\alpha(r) e_\beta(t) \bigr] \dif r \Bigr|
    \;\leq\; \tfrac{1}{2} e^{C_T} E[\Gamma_T] \sum_{i=1}^m (s_i - s_{i-1})^2
  \]
  --- the main terms telescoping \emph{exactly}, with no limit needed --- and the
  right-hand side is at most
  $\tfrac{1}{2} e^{C_T} E[\Gamma_T] \, s \max_i (s_i - s_{i-1}) \to 0$.  This is
  \begin{equation}
    \label{eq:Fpartial1}
    F(s,t) - F(0,t)
      = \int_0^s E\bigl[ \bigl( f(X(r),Y(t)) \alpha(X(r)) + g(X(r),Y(t)) \bigr)
          e_\alpha(r) e_\beta(t) \bigr] \dif r .
  \end{equation}
  Symmetrically,  Symmetrically,
  \begin{equation}
    \label{eq:Fpartial2}
    F(s,t) - F(s,0)
      = \int_0^t E\bigl[ \bigl( f(X(s),Y(r)) \beta(Y(r)) + h(X(s),Y(r)) \bigr)
          e_\alpha(s) e_\beta(r) \bigr] \dif r .
  \end{equation}
  Thus $F$ satisfies the hypotheses of Lemma~\ref{lem:calculus}, with $F_1$ and
  $F_2$ the integrands in \eqref{eq:Fpartial1} and \eqref{eq:Fpartial2}; the
  integrability \eqref{eq:calcint} follows from \eqref{eq:dual1} and
  \eqref{eq:dual2}.  Now \eqref{eq:calcconc} is exactly \eqref{eq:dualconc}.
\end{proof}

\begin{remark}[\EK, Remark 4.4.12]
  \label{rem:dualext}
  Equation \eqref{eq:dualconc} can frequently be extended to more general
  $\alpha$ and $\beta$ by approximating them by bounded functions.
\end{remark}

\begin{corollary}[Duality relation; \EK, Corollary 4.4.13; (T3)]
  \label{cor:dualrel}
  If, in addition to the conditions of Theorem~\ref{thm:duality},
  \begin{equation}
    \label{eq:dualbalance}
    g(x,y) + \alpha(x) f(x,y) \;=\; h(x,y) + \beta(y) f(x,y) ,
    \qquad (x,y) \in E_1 \times E_2 ,
  \end{equation}
  then for all $t \geq 0$
  \begin{equation}
    \label{eq:dualrel}
    E\Bigl[ f(X(t),Y(0)) \exp\Bigl\{ \int_0^t \alpha(X(u)) \dif u \Bigr\} \Bigr]
    = E\Bigl[ f(X(0),Y(t)) \exp\Bigl\{ \int_0^t \beta(Y(u)) \dif u \Bigr\} \Bigr].
  \end{equation}
\end{corollary}

\begin{proof}
  Under \eqref{eq:dualbalance} the right-hand side of \eqref{eq:dualconc}
  vanishes, so \eqref{eq:dualrel} holds for almost every $t$.  It extends to all
  $t$ because $t \mapsto F(t,0)$ and $t \mapsto F(0,t)$ are continuous, which
  follows from \eqref{eq:Fpartial1} and \eqref{eq:Fpartial2}.
\end{proof}

\begin{remark}
  \label{rem:dualischain}
  Corollary~\ref{cor:dualrel} is Proposition~\ref{prop:haar}\ref{it:haar_cont}
  applied to $\Phi = F$ of \eqref{eq:F}, with
  $\gamma_1$ and $\gamma_2$ the integrands of \eqref{eq:Fpartial1} and
  \eqref{eq:Fpartial2}; the balance condition \eqref{eq:dualbalance} is exactly
  $\gamma_1 = \gamma_2$.  The whole of Theorem~\ref{thm:duality} is thus
  Lemma~\ref{lem:chain} plus the computation \eqref{eq:Fincrement} that produces
  the increment representations \eqref{eq:incrementrep} from the martingale
  hypotheses.  Only the latter uses probability; and it uses no path regularity,
  no topology on $E_1$, $E_2$ and no boundedness, so it is already stated at the
  level of generality of \CPS.  What it does use, and what
  Remark~\ref{rem:dualscope} isolates, is the translation invariance of the
  clock.
\end{remark}

The discrete counterpart, which \EK{} do not state, now costs nothing.

\begin{corollary}[Duality for Markov chains; $\T = \N_0$, (E0), $\iota = \mathrm{p}$]
  \label{cor:dualdiscrete}
  Let $P$ and $Q$ be transition kernels on measurable spaces $E_1$ and $E_2$, and
  let $f \in \Bdd(E_1 \times E_2)$ satisfy the duality relation
  \begin{equation}
    \label{eq:dualdiscrete}
    \int f(x',y) \, P(x,\dif x') \;=\; \int f(x,y') \, Q(y,\dif y') ,
    \qquad (x,y) \in E_1 \times E_2 .
  \end{equation}
  Let $X$ and $Y$ be independent Markov chains with kernels $P$ and $Q$ and
  $X(0) = x$, $Y(0) = y$.  Then
  \begin{equation}
    \label{eq:dualdiscreteconc}
    E\bigl[ f(X(n),y) \bigr] \;=\; E\bigl[ f(x,Y(n)) \bigr] ,
    \qquad n \in \N_0 .
  \end{equation}
\end{corollary}

\begin{proof}
  Put $\Phi(n,m) = E[f(X(n),Y(m))]$.  Writing $P$ for the action of the kernel on
  the first variable and $Q$ for its action on the second, the Markov property
  and independence give
  $\Phi(n+1,m) - \Phi(n,m) = E[((P-I)f)(X(n),Y(m))]$ and
  $\Phi(n,m+1) - \Phi(n,m) = E[((Q-I)f)(X(n),Y(m))]$, which is
  \eqref{eq:incrementrep} for $q = $ counting measure with
  $\gamma_1(n,m) = E[((P-I)f)(X(n),Y(m))]$ and $\gamma_2$ its analogue.  By
  \eqref{eq:dualdiscrete}, $(P-I)f = (Q-I)f$, so $\gamma_1 = \gamma_2$, and
  Proposition~\ref{prop:haar}\ref{it:haar_disc} applies.
\end{proof}

\begin{remark}
  \label{rem:dualdiscretemp}
  In the language of Section~\ref{sec:MP}, \eqref{eq:dualdiscrete} says that
  $(f(\cdot,y),\, g(\cdot,y)) \in A$ and $(f(x,\cdot),\, g(x,\cdot)) \in B$ hold
  with the \emph{same} $g = (P-I)f = (Q-I)f$, where $A$ and $B$ denote the
  one-step generators; the associated test process is then the Doob
  decomposition
  \[
    f(X(n),y) - \sum_{k < n} g(X(k),y) ,
  \]
  i.e.\ the compensator in the convention $[0,t)$ of
  Remark~\ref{rem:conventionclash}.  So Corollary~\ref{cor:dualdiscrete} is the
  martingale-problem duality theorem for $\T = \N_0$, and the convention it
  forces is the predictable one.
\end{remark}

\begin{corollary}[Stopped version; \EK, Corollary 4.4.14; (T3)]
  \label{cor:dualstopped}
  Let $(\Filt_t)$ and $(\Gilt_t)$ be independent filtrations, let $X$ be
  $(\Filt_t)$-progressive with $(\Filt_t)$-stopping time $\tau$, and let $Y$ be
  $(\Gilt_t)$-progressive with $(\Gilt_t)$-stopping time $\sigma$.  Suppose
  \eqref{eq:dual1} and \eqref{eq:dual2} hold with $X$, $Y$ replaced by
  $X(\cdot \wedge \tau)$, $Y(\cdot \wedge \sigma)$, and that
  $f(X(t \wedge \tau),y) - \int_0^{t \wedge \tau} g(X(s),y) \dif s$ is an
  $(\Filt_t)$-martingale for each $y$, and
  $f(x,Y(t \wedge \sigma)) - \int_0^{t \wedge \sigma} h(x,Y(s)) \dif s$ is a
  $(\Gilt_t)$-martingale for each $x$.  Then the analogue of
  \eqref{eq:dualconc} holds with $X$, $Y$ replaced by the stopped processes and
  the integrand multiplied by the indicators $\one_{\{s \leq \tau\}}$ resp.\
  $\one_{\{t-s \leq \sigma\}}$.
\end{corollary}

\begin{proof}
  Rewrite the martingale in the hypothesis as
  $f(X(t \wedge \tau),y) - \int_0^t \one_{\{s \leq \tau\}} g(X(s),y) \dif s$ and
  repeat the proof of Theorem~\ref{thm:duality}.
\end{proof}

\begin{remark}
  \label{rem:dualstopped}
  Corollary~\ref{cor:dualstopped} is the practically relevant form: the
  estimates \eqref{eq:dual1} are often unavailable globally but easy on the
  exit times $\tau$ from compact sets.  One then lets $\tau, \sigma \to \infty$
  along an exhausting sequence; the integrand on the right of the stopped version
  of \eqref{eq:dualconc} tends to zero, and the only difficulty is to justify the
  interchange of limits and expectations (\EK, Remark 4.4.16).  This is also the
  route by which duality is applied to \emph{local} martingale problems.
\end{remark}

We can now close the circle with Section~\ref{sec:uniqueness}.

\begin{corollary}[Uniqueness via duality; (T3)]
  \label{cor:uniqviadual}
  Assume (E2), let $A \subset \Bdd(E) \times \Bdd(E)$, let $(E_2,\mathcal{E}_2)$ be
  a measurable space, and let $B \subset \Meas(E_2) \times \Meas(E_2)$.  Suppose
  there is $f \in \Meas(E \times E_2)$ such that:
  \begin{enumerate}[label=(\roman*)]
  \item $\{ f(\cdot,y) : y \in E_2 \} \subset \Bdd(E)$ is separating for
    $\Prob(E)$;
  \item $(f(\cdot,y), g(\cdot,y)) \in A$ for every $y \in E_2$, and
    $(f(x,\cdot), h(x,\cdot)) \in B$ for every $x \in E$, with $g = h$;
  \item for every $y \in E_2$ there is a solution $Y^y$ of the martingale problem
    for $(B, \delta_y)$, and for every solution $X$ of the martingale problem for
    $A$ with $X$ independent of $Y^y$ the hypotheses
    \eqref{eq:dual1}--\eqref{eq:dual2} hold with $\alpha = \beta = 0$.
  \end{enumerate}
  Then for every $\mu \in \Prob(E)$ the martingale problem for $(A,\mu)$ has at
  most one solution, every solution is Markov, and under the additional
  hypotheses of Theorem~\ref{thm:uniqueness}\ref{it:uniq_b}--\ref{it:uniq_c} it
  is strong Markov.
\end{corollary}

\begin{proof}
  Let $X$ solve the martingale problem for $(A,\delta_x)$ and let $Y^y$ be as in
  (iii), independent of $X$.  By (ii), \eqref{eq:dualbalance} holds with
  $\alpha = \beta = 0$, so Corollary~\ref{cor:dualrel} gives
  \begin{equation}
    \label{eq:dualxy}
    E\bigl[ f(X(t), y) \bigr] \;=\; E\bigl[ f(x, Y^y(t)) \bigr] ,
    \qquad t \geq 0 .
  \end{equation}
  The right-hand side does not depend on $X$.  Hence $\int f(\cdot,y) \dif \nu_t$
  is the same for the one-dimensional distributions $\nu_t = P X(t)^{-1}$ of any
  two solutions, and by (i) these coincide.  For general $\mu$, integrate
  \eqref{eq:dualxy} over $\mu(\dif x)$, which is legitimate because the
  one-dimensional distributions of a solution for $(A,\mu)$ are mixtures of those
  for $(A,\delta_x)$ by Theorem~\ref{thm:uniqueness}\ref{it:uniq_a}.  Now apply
  Theorem~\ref{thm:uniqueness}.
\end{proof}

% =====================================================================
\section{Existence}
\label{sec:convergence}
% =====================================================================

Sections~\ref{sec:uniqueness} and \ref{sec:duality} are a uniqueness theory, and
by themselves they are vacuous: nothing so far exhibits a single solution of a
single martingale problem.  This section collects the four routes to existence,
in increasing order of cost.

\begin{enumerate}[label=(\alph*)]
\item \emph{From a transition semigroup} (Section~\ref{ssec:exsemigroup}).
  Immediate, and correspondingly weak: it presupposes the object one usually
  wants to construct.
\item \emph{Explicit construction} (Section~\ref{ssec:jump}).  Jump processes are
  built by hand out of a Markov chain and exponential holding times, and the
  martingale property is verified by a two-line computation with the exponential
  distribution.  This is the only place in the manuscript where a solution is
  written down.
\item \emph{From stochastic differential equations}
  (Section~\ref{ssec:exsde}).  The most important route in practice and the most
  expensive here, since it presupposes stochastic integration; we state the chain
  of results and prove none of it.
\item \emph{By convergence} (Sections~\ref{ssec:EKconv}--\ref{ssec:CPSconv}).
  Weak limits of solutions of approximating problems.  Conditional, like~(a), but
  on much less.
\end{enumerate}

Routes (b) and (c) are the two that produce a solution out of nothing, and it is
worth noticing before we start that they are the same argument at two levels of
difficulty: both are Picard--Lindel\"of.  In (b) it is a Gronwall estimate on the
first-jump equation \eqref{eq:firstjump}; in (c) it is the Picard iteration for a
Lipschitz stochastic differential equation.  What differs is the space on which
the contraction acts.

\subsection{From a transition semigroup}
\label{ssec:exsemigroup}

\begin{proposition}[\EK, Proposition 4.1.7; (T3) + (E0)]
  \label{prop:exsemigroup}
  Let $X$ be a Markov process with measurable transition semigroup
  $\{T(t)\}$ on $\Bdd(E)$ and full generator $\hat{A}$ in the sense of
  Definition~\ref{def:fullgenerator}.  Then $X$ solves the martingale problem for
  $\hat{A}$.
\end{proposition}

This is Remark~\ref{rem:fullgenerator}, restated as an existence statement, and it
is recorded only for completeness: it turns a Markov process into a solution, and
producing the Markov process is the whole difficulty.  The remaining three routes
do not presuppose one.

\subsection{Jump processes}
\label{ssec:jump}

Throughout this subsection $\T = \Rp$ with $q$ Lebesgue measure, so that the two
conventions of Definition~\ref{def:clock} agree and $\iota$ plays no r\^ole.  The
state space needs only (E0) for the construction and the martingale property;
(E2) enters only when one wants to say that the paths lie in $\DE$.

\begin{setting}[Jump data; (T3) + (E0)]
  \label{set:jumpdata}
  Let $\lambda : E \to [0,\infty)$ be measurable (the \emph{rate}) and let $\mu$ be
  a transition kernel on $E$ (the \emph{jump distribution}).  Define
  \begin{equation}
    \label{eq:jumpgen}
    A f(x) \;=\; \lambda(x) \int_E \bigl( f(y) - f(x) \bigr) \, \mu(x,\dif y) ,
    \qquad f \in \Bdd(E) ,
  \end{equation}
  and $A = \{(f, Af) : f \in \Bdd(E)\}$.  If
  $\bar{\lambda} \coloneqq \sup_x \lambda(x) < \infty$ then $A$ maps $\Bdd(E)$ into
  itself with $\lVert Af \rVert \leq 2 \bar{\lambda} \lVert f \rVert$; in general
  $Af$ need not be bounded.
\end{setting}

\begin{definition}[The construction; (T3) + (E0)]
  \label{def:jumpconstruction}
  Let $\nu \in \Prob(E)$.  On some probability space let $(Y_n)_{n \geq 0}$ be a
  Markov chain with initial distribution $\nu$ and transition kernel $\mu$, and
  let $\varepsilon_1, \varepsilon_2, \dots$ be i.i.d.\ exponential random
  variables of parameter $1$, independent of $(Y_n)$.  Put
  \begin{equation}
    \label{eq:jumptimes}
    \tau_0 = 0 , \qquad
    \tau_{n+1} = \tau_n + \frac{\varepsilon_{n+1}}{\lambda(Y_n)} ,
    \qquad \zeta = \lim_n \tau_n ,
  \end{equation}
  with the convention $1/0 = \infty$, and
  \begin{equation}
    \label{eq:jumpprocess}
    X_t = Y_n \quad \text{for } \tau_n \leq t < \tau_{n+1} , \qquad t < \zeta .
  \end{equation}
  We call $\zeta$ the \emph{explosion time} and $(\tau_n)$ the \emph{jump times}.
\end{definition}

Three observations before the theorem.  The paths of $X$ are piecewise constant
and right continuous by construction, hence lie in $\DE$ as soon as $E$ carries a
topology, i.e.\ under (E2).  If $\lambda(Y_n) = 0$ then $\tau_{n+1} = \infty$ and
$X$ is absorbed at $Y_n$.  And each $\tau_n$ is a functional of the path --- the
time of its $n$-th jump --- which is what will make it a legitimate localizing
sequence in the sense of \ref{it:L1}.

\begin{lemma}[No explosion for bounded rates; (T3) + (E0)]
  \label{lem:noexplosion}
  If $\bar{\lambda} < \infty$ then $\zeta = \infty$ almost surely.
\end{lemma}

\begin{proof}
  $\tau_{n+1} - \tau_n \geq \varepsilon_{n+1}/\bar{\lambda}$, and
  $\sum_n \varepsilon_{n+1}/\bar{\lambda} = \infty$ a.s.\ by the strong law of
  large numbers, the $\varepsilon_n$ being i.i.d.\ with mean $1$.
\end{proof}

\begin{theorem}[Jump processes solve their martingale problem; (T3) + (E0)]
  \label{thm:jumpMP}
  Let $\bar{\lambda} < \infty$ and let $X$ be as in
  Definition~\ref{def:jumpconstruction}.  Then $X$ solves the martingale problem
  for $(A,\nu)$ of \eqref{eq:jumpgen}.
\end{theorem}

\begin{proof}
  Fix $f \in \Bdd(E)$ and write $\mu f(x) = \int f(y)\mu(x,\dif y)$, so that
  $Af = \lambda \cdot (\mu f - f)$.  Let $\Gilt_t = \sigma(X_u : u \leq t)$ and
  \[
    M_t \;=\; f(X_t) - f(X_0) - \int_0^t Af(X_r) \dif r .
  \]

  \emph{Step 1: the paths make both terms explicit.}  Since $X$ is constant on
  $[\tau_n, \tau_{n+1})$,
  \begin{align*}
    \int_0^t Af(X_r)\dif r
      &= \sum_{n \geq 0} Af(Y_n) \bigl( (\tau_{n+1} \wedge t) - (\tau_n \wedge t) \bigr) , \\
    f(X_t) - f(X_0)
      &= \sum_{n \geq 1} \bigl( f(Y_n) - f(Y_{n-1}) \bigr) \one_{\{\tau_n \leq t\}} ,
  \end{align*}
  both sums being finite a.s.\ by Lemma~\ref{lem:noexplosion}.  Hence
  \begin{equation}
    \label{eq:jumpbrackets}
    M_t \;=\; \sum_{n \geq 0} D_n , \qquad
    D_n = \bigl( f(Y_{n+1}) - f(Y_n) \bigr) \one_{\{\tau_{n+1} \leq t\}}
        - Af(Y_n) \bigl( (\tau_{n+1} \wedge t) - (\tau_n \wedge t) \bigr) .
  \end{equation}

  \emph{Step 2: each bracket is a martingale difference.}  Let
  $\mathcal{H}_n = \sigma(Y_0,\dots,Y_n,\tau_1,\dots,\tau_n)$.  Conditionally on
  $\mathcal{H}_n$, and on the event $\{\tau_n \leq t\}$, put $x = Y_n$,
  $\lambda = \lambda(x)$ and $a = t - \tau_n \geq 0$.  By construction
  $\tau_{n+1} - \tau_n$ is exponential with parameter $\lambda$ and independent of
  $Y_{n+1} \sim \mu(x,\cdot)$, so
  \begin{align*}
    E\bigl[ (f(Y_{n+1}) - f(Y_n)) \one_{\{\tau_{n+1} \leq t\}} \bigm| \mathcal{H}_n \bigr]
      &= \bigl( \mu f(x) - f(x) \bigr) \bigl( 1 - e^{-\lambda a} \bigr) , \\
    E\bigl[ (\tau_{n+1} \wedge t) - \tau_n \bigm| \mathcal{H}_n \bigr]
      &= E\bigl[ \mathrm{Exp}(\lambda) \wedge a \bigr]
       = \frac{1 - e^{-\lambda a}}{\lambda} .
  \end{align*}
  Multiplying the second identity by
  $Af(x) = \lambda(\mu f(x) - f(x))$ gives exactly the first.  Hence
  $E[D_n \mid \mathcal{H}_n] = 0$ on $\{\tau_n \leq t\}$, and $D_n = 0$ on
  $\{\tau_n > t\}$.  \emph{This cancellation --- the factor $\lambda$ in
  \eqref{eq:jumpgen} against the mean $1/\lambda$ of the holding time --- is the
  whole content of the theorem.}

  \emph{Step 3: from $0$ to $s$.}  For $s \leq t$ the same computation applies with
  $\tau_n$ replaced by $s$ on the interval straddling $s$: by the memorylessness
  of the exponential distribution, conditionally on $\Gilt_s$ the residual holding
  time $\tau_{N_s+1} - s$ is again exponential with parameter $\lambda(X_s)$ and
  the subsequent chain starts afresh at $X_s$, where $N_s$ is the number of jumps
  in $[0,s]$.  Summing \eqref{eq:jumpbrackets} from that interval onwards and
  applying Step~2 to each term gives $E[M_t - M_s \mid \Gilt_s] = 0$.  As
  $|M_t| \leq 2\lVert f\rVert + 2\bar{\lambda}\lVert f\rVert t$ is bounded on
  bounded intervals, $M$ is a $(\Gilt_t)$-martingale, and since
  ${}^{*}\Filt^X_t = \Filt^X_t \subset \Gilt_t$ for the right continuous $X$, the
  same holds for the filtration of Definition~\ref{def:markovMP}.
\end{proof}

The construction gives existence.  Uniqueness comes from
Section~\ref{sec:uniqueness}, and the input it needs is a Gronwall estimate.

\begin{proposition}[Well-posedness for bounded rates; (T3) + (E2)]
  \label{prop:jumpwellposed}
  Let $\bar{\lambda} < \infty$.  Then the martingale problem for $(A,\nu)$ is well
  posed for every $\nu \in \Prob(E)$, every solution is a Markov process, and if
  $E$ satisfies (E2) it is strong Markov.
\end{proposition}

\begin{proof}
  Existence is Theorem~\ref{thm:jumpMP}.  For uniqueness it suffices, by
  Theorem~\ref{thm:absuniq}, to show that the one-dimensional distributions are
  determined by $\nu$; and since $\Bdd(E)$ is separating it suffices to determine
  $v(t,x) = E_x[f(X_t)]$ for $f \in \Bdd(E)$.  Conditioning on the first jump ---
  which under any solution has the law prescribed by \eqref{eq:jumpgen}, by
  Proposition~\ref{prop:fddchar} --- gives the \emph{first-jump equation}
  \begin{equation}
    \label{eq:firstjump}
    v(t,x) \;=\; e^{-\lambda(x)t} f(x)
      + \int_0^t \lambda(x) e^{-\lambda(x)s}
        \bigl( \mu\, v(t-s,\cdot) \bigr)(x) \dif s .
  \end{equation}
  If $v_1, v_2$ are two bounded measurable solutions and $w = v_1 - v_2$, then
  \[
    |w(t,x)| \;\leq\; \int_0^t \lambda(x) e^{-\lambda(x)s}
      \lVert w(t-s,\cdot) \rVert \dif s
    \;\leq\; \bar{\lambda} \int_0^t \lVert w(r,\cdot) \rVert \dif r ,
  \]
  so $\rho(t) = \lVert w(t,\cdot)\rVert$ satisfies
  $\rho(t) \leq \bar\lambda\int_0^t\rho$, and Gronwall gives $\rho \equiv 0$.
  Hence the one-dimensional distributions are unique, and
  Theorem~\ref{thm:absuniq} applies.  The strong Markov property is
  Theorem~\ref{thm:absstrongmarkov}, whose hypotheses hold because the paths are
  piecewise constant and $M$ has increments bounded on bounded intervals.
\end{proof}

\begin{remark}[Explosion, and why localization is not academic]
  \label{rem:jumpexplosion}
  Drop the assumption $\bar\lambda < \infty$.  Then $\zeta$ may be finite with
  positive probability --- take $E = \N$, $\mu(n,\cdot) = \delta_{n+1}$ and
  $\lambda(n) = 2^n$, so that $E[\zeta] = \sum_n 2^{-n} < \infty$ --- and
  \eqref{eq:jumpprocess} defines $X$ only on $[0,\zeta)$.  There is then no
  solution of the \emph{global} martingale problem for $A$ on $\Rp$, and the
  honest statement is the local one:
  \begin{itemize}
  \item Step~2 of the proof of Theorem~\ref{thm:jumpMP} is a finite computation
    and is unaffected, so $M^{\tau_n}$ is a martingale for every $n$;
  \item $\tau_n$ is a functional of the path, the same for every solution, and
    $\tau_n \uparrow \zeta$;
  \item hence $(\tau_n)$ is a localizing system in the sense of
    Definition~\ref{def:localizing}, and \ref{it:L1} holds by construction rather
    than by assumption --- exactly as Lemma~\ref{lem:L1auto} predicts, since the
    stopped test processes are bounded.
  \end{itemize}
  So $X$ solves the local martingale problem for $A$ on $[0,\zeta)$, and
  Section~\ref{ssec:localmp} applies verbatim.  This is the example that
  Definition~\ref{def:localizing} was written for; up to this point it had none.
\end{remark}

\begin{remark}[Discrete time, and the two conventions]
  \label{rem:jumpdiscrete}
  Deleting the holding times from Definition~\ref{def:jumpconstruction} leaves the
  chain $(Y_n)$ itself, which is the case $\T = \N_0$ with $q$ counting measure.
  There $A = \mu - I$ under the predictable convention $\iota = \mathrm{p}$, so
  that $M_n = f(Y_n) - f(Y_0) - \sum_{k<n}(\mu f - f)(Y_k)$ is the Doob
  decomposition; under $\iota = \mathrm{o}$ one would obtain the condition
  $\mu f - f = \mu(\mu f - f)$ instead, which is not what one wants.  This is the
  concrete form of Remark~\ref{rem:convention}, and it is also the case in which
  Corollary~\ref{cor:dualdiscrete} lives.
\end{remark}

\begin{remark}[What this example is for]
  \label{rem:jumppoint}
  Besides being the only explicit solution in the manuscript, the construction
  exercises three of its abstractions and is the reason they were made.  The state
  space needs only (E0): $\lambda$ measurable and $\mu$ a kernel, no topology, no
  metric, no Polish structure --- so jump processes on a standard Borel $E$, for
  instance on a space of measures or of distributions, are covered by
  Theorem~\ref{thm:jumpMP} as it stands.  The localizing system of
  Remark~\ref{rem:jumpexplosion} is a genuine instance of \ref{it:L1}--\ref{it:L3}.
  And the discrete case of Remark~\ref{rem:jumpdiscrete} is a genuine instance of a
  clock with atoms.
\end{remark}

\subsection{From stochastic differential equations}
\label{ssec:exsde}

For $E = \R^d$ and continuous paths there is a route that is more important than
all the others put together, and we state it without proofs.  Let
$\sigma : \R^d \to \R^{d \times m}$ and $b : \R^d \to \R^d$ be measurable, put
$a = \sigma\sigma^{\top}$, and let
\begin{equation}
  \label{eq:sdegen}
  A f(x) \;=\; \tfrac{1}{2} a^{ij}(x) \partial_{ij} f(x) + b^i(x) \partial_i f(x) ,
  \qquad f \in C^\infty_c(\R^d) .
\end{equation}
A \emph{weak solution} of the equation $(\sigma,b)$ is a pair $(X,B)$ on some
filtered probability space with $B$ an $m$-dimensional Brownian motion and
$X_t = X_0 + \int_0^t \sigma(X_s)\dif B_s + \int_0^t b(X_s)\dif s$.

\begin{fact}[Picard--Lindel\"of for SDEs]
  \label{fact:picard}
  If $\sigma$ and $b$ are Lipschitz, then for every $x$ the equation
  $(\sigma,b)$ has a strong solution with $X_0 = x$, and pathwise uniqueness
  holds.  The proof is Picard iteration in $L^2$, with the It\^o isometry in place
  of the triangle inequality.
\end{fact}

\begin{fact}[Yamada--Watanabe]
  \label{fact:yamadawatanabe}
  Pathwise uniqueness together with weak existence implies uniqueness in law.
\end{fact}

\begin{fact}[Stroock--Varadhan; \KA, Theorem 32.7]
  \label{fact:strookvaradhan}
  For progressive $\sigma, b$ and $P \in \Prob(C_{\Rp,\R^d})$ the following are
  equivalent: $P$ is the law of a weak solution of $(\sigma,b)$; and $P$ solves
  the local martingale problem for $A$ of \eqref{eq:sdegen}.
\end{fact}

\begin{corollary}[Well-posedness for Lipschitz coefficients; (T3), $E = \R^d$]
  \label{cor:sdewellposed}
  If $\sigma$ and $b$ are Lipschitz, the local martingale problem for $A$ of
  \eqref{eq:sdegen} is well posed for every initial distribution, and its solution
  is a strong Markov process.
\end{corollary}

\begin{proof}
  Facts~\ref{fact:picard} and \ref{fact:yamadawatanabe} give weak existence and
  uniqueness in law; Fact~\ref{fact:strookvaradhan} transfers both to the local
  martingale problem.  The Markov properties are
  Theorems~\ref{thm:localuniq} and \ref{thm:absstrongmarkov}, or directly
  \KA, Theorem~32.11.
\end{proof}

\begin{remark}[What a proof would cost]
  \label{rem:sdecost}
  Section~\ref{ssec:ch2} contains nothing about stochastic integration, and this
  is the only place in the manuscript that reaches outside \EK, Chapters~1--3.  A
  self-contained treatment of Corollary~\ref{cor:sdewellposed} would need the
  It\^o integral against a continuous local martingale, quadratic variation and
  the Kunita--Watanabe inequality, It\^o's formula, and L\'evy's characterization
  of Brownian motion --- the last being what turns a solution of the martingale
  problem back into a solution of the equation in Fact~\ref{fact:strookvaradhan}.
  That is a chapter, not a section, and it is deliberately not attempted here.

  For the formalization the situation is less bad than it looks: the
  \texttt{brownian-motion} development of Remark~\ref{rem:bmrepo} is building
  stochastic integration in Lean, so the dependency may become available rather
  than having to be created.  Until then, Section~\ref{ssec:jump} is the route
  that can actually be carried out.
\end{remark}

\subsection{By convergence: the Ethier--Kurtz version}
\label{ssec:EKconv}

The remaining route obtains solutions as weak limits of solutions of approximating
martingale problems.  We first state the classical result and then its
generalization.


\begin{lemma}[\EK, Lemma 4.5.1; (T3)]
  \label{lem:EKconv}
  Assume (E3), let $A \subset \Cb(E) \times \Cb(E)$ and
  $A_n \subset \Bdd(E) \times \Bdd(E)$ for $n \in \N$.  Suppose that for each
  $(f,g) \in A$ there exist $(f_n,g_n) \in A_n$ with
  \begin{equation}
    \label{eq:approxA}
    \lim_{n \to \infty} \lVert f_n - f \rVert = 0,
    \qquad
    \lim_{n \to \infty} \lVert g_n - g \rVert = 0 .
  \end{equation}
  If for each $n$ the process $X_n$ is a solution of the martingale problem for
  $A_n$ with sample paths in $\DE$, and if $X_n \wto X$ in $\DE$, then $X$ is a
  solution of the martingale problem for $A$.
\end{lemma}

\begin{proof}
  Let $D(X) = \{u \geq 0 : P\{X(u) = X(u-)\} = 1\}$; by
  Fact~\ref{fact:Dcountable} its complement is countable, so $D(X)$ is dense.
  Let $0 \leq t_1 \leq \dots \leq t_k \leq t < s$ with
  $t_i, t, s \in D(X)$, and $h_i \in \Cb(E)$.  For $u \in D(X)$ the evaluation
  map $\omega \mapsto \omega(u)$ is continuous at $\omega$ for $P X^{-1}$-almost
  every $\omega \in \DE$, hence the functional
  \[
    \omega \;\longmapsto\;
      \Bigl( f(\omega(s)) - f(\omega(t)) - \int_t^s g(\omega(u)) \dif u \Bigr)
      \prod_{i=1}^k h_i(\omega(t_i))
  \]
  is bounded and $P X^{-1}$-a.s.\ continuous on $\DE$.  By the continuous
  mapping theorem (Fact~\ref{fact:cmt}) and \eqref{eq:approxA},
  \begin{align}
    \label{eq:convlimit}
    E\Bigl[ \Bigl( f(X(s)) - f(X(t)) - \int_t^s g(X(u)) \dif u \Bigr)
      \prod_{i=1}^k h_i(X(t_i)) \Bigr] \\
    \notag
    = \lim_{n \to \infty}
      E\Bigl[ \Bigl( f_n(X_n(s)) - f_n(X_n(t)) - \int_t^s g_n(X_n(u)) \dif u \Bigr)
      \prod_{i=1}^k h_i(X_n(t_i)) \Bigr] \;=\; 0 ,
  \end{align}
  the last equality because $X_n$ solves the martingale problem for $A_n$.  By
  right continuity of $X$ and density of $D(X)$, \eqref{eq:convlimit} extends to
  all $0 \leq t_1 \leq \dots \leq t_k \leq t < s$.  By
  Proposition~\ref{prop:fddchar}, $X$ solves the martingale problem for $A$.
\end{proof}

\begin{remark}[\EK, Remark 4.5.2]
  \label{rem:EKrelcompact}
  Lemma~\ref{lem:EKconv} presupposes convergence.  Suppose $(E,r)$ is Polish and
  $\dom(A)$ contains an algebra that separates points and vanishes nowhere; by
  Fact~\ref{fact:stoneweierstrass} such an algebra is separating, and it is dense
  in $\Cb(E)$ in the topology of uniform convergence on compact sets.  Then
  $\{X_n\}$ is relatively compact provided \eqref{eq:approxA} holds for each
  $(f,g) \in A$ and the compact containment condition \eqref{eq:ccn} holds.  This
  is Facts~\ref{fact:relcompact}, \ref{fact:relcompact2},
  \ref{fact:fddconv} and \ref{fact:prohorov}.  Combining with
  Lemma~\ref{lem:EKconv}: every limit point of $\{X_n\}$ solves the martingale
  problem for $A$; in particular the martingale problem for $A$ has a solution
  with sample paths in $\DE$.  If in addition uniqueness holds (e.g.\ by
  Theorem~\ref{thm:uniqueness} or Corollary~\ref{cor:uniqviadual}), then
  $X_n \wto X$ for the unique solution $X$.
\end{remark}

\subsection{By convergence: the abstract theorem}
\label{ssec:absconv}

Theorem~\ref{thm:CPSconv} below is an instance of one
statement, which uses neither an operator, nor c\`adl\`ag paths, nor the
Skorokhod topology, nor any structure on the time index beyond a dense subset ---
and none at all if that subset is all of $\T$.  It is \CPS, Theorem~3.14 and
Corollary~3.17, and it is the cheapest genuine theorem in this manuscript; we
give it with its proof, since it is what the formalization should target
(step~\ref{it:F5a} of Section~\ref{sec:formalization}).  Lemma~\ref{lem:EKconv}
is in turn a special case of Theorem~\ref{thm:CPSconv}, so all of
Section~\ref{sec:convergence} rests on Theorem~\ref{thm:absconv}.

\begin{definition}[$P$-continuity at $X$; (T0) + (E3)]
  \label{def:Pcont}
  A Borel map $\psi : F \to \R$ is \emph{$P$-continuous at $X$} if there is
  $C \in \Bor(F)$ with $P\{X \in C\} = 1$ such that $\psi(\alpha_m) \to \psi(\alpha)$
  whenever $\alpha_m \to \alpha$ in $F$ with $\alpha \in C$.
\end{definition}

\begin{theorem}[Abstract convergence; \CPS, Theorem 3.14 and Corollary 3.17; (T0) + (E3)]
  \label{thm:absconv}
  In Setting~\ref{set:abstract}, let $\XX$ be canonical for $X$, let
  $D \subset \T$, and for each $n$ let $B^n = (\Omega^n,\Filt^n,\mathbb{F}^n,P^n)$
  be a filtered probability space carrying a process $X^n$ with paths in $F$.
  Assume:
  \begin{enumerate}[label=(C\arabic*)]
  \item \label{it:C1} $X^n \wto X$ weakly on $F$;
  \item \label{it:C2} $\XX$ admits a determining set $\ZZ^\circ$
    (Definition~\ref{def:canonical});
  \item \label{it:C3} for every $Y \in \XX$ there is a canonical version
    $(Y^\circ_t)$ such that for all $t \in D$, $s \in D \cap \T_{\leq t}$ and
    $Z^\circ_s \in \ZZ^\circ_s$:
    \begin{enumerate}[label=(\alph*)]
    \item \label{it:C3a} $Y^\circ_t$ and $Y^\circ_t Z^\circ_s$ are $P$-continuous
      at $X$;
    \item \label{it:C3b} $\{ Y^\circ_r(X^n) : r \in D \cap \T_{\leq t},\ n \in \N \}$
      is uniformly integrable;
    \item \label{it:C3c}
      $\displaystyle \lim_{n \to \infty}
        E^{P^n}\bigl[ \bigl( Y^\circ_t(X^n) - Y^\circ_s(X^n) \bigr) Z^\circ_s(X^n) \bigr]
        = 0$.
    \end{enumerate}
  \end{enumerate}
  Then $E^P[Y_t \mid \Filt_s] = Y_s$ for all $s \leq t$ in $D$ and all
  $Y \in \XX$.  If moreover $D = \T$, then $P \in \Msol(\XX)$.  The same
  conclusion holds for a proper $D \subsetneq \T$ provided $\T$ satisfies (T2b)
  with $D$ its countable dense subset and every $Y \in \XX$ is right continuous.
\end{theorem}

\begin{proof}
  Fix $Y \in \XX$ with canonical version $Y^\circ$, and $t \in D$.

  \emph{Step 0: integrability.}  For $r \in D \cap \T_{\leq t}$ the map
  $|Y^\circ_r|$ is $P$-continuous at $X$ by \ref{it:C3a}, so
  $|Y^\circ_r(X^n)| \wto |Y^\circ_r(X)|$ by \ref{it:C1} and
  Fact~\ref{fact:cmt}; the family is uniformly integrable by \ref{it:C3b}, so
  Fact~\ref{fact:ui} gives
  $E^{P^n}|Y^\circ_r(X^n)| \to E^P|Y^\circ_r(X)| < \infty$.  In particular
  $Y_r = Y^\circ_r(X) \in L^1(P)$.

  \emph{Step 1: the martingale identity on $D$.}
  Let $s \in D \cap \T_{\leq t}$ and $Z^\circ_s \in \ZZ^\circ_s$, and put
  $\psi = (Y^\circ_t - Y^\circ_s) Z^\circ_s : F \to \R$.  By \ref{it:C3a} and
  Definition~\ref{def:Pcont} --- which is stable under differences and products,
  the exceptional sets being intersected --- $\psi$ is $P$-continuous at $X$.
  Hence $\psi(X^n) \wto \psi(X)$ by \ref{it:C1} and Fact~\ref{fact:cmt}.  The
  family $\{\psi(X^n)\}_n$ is uniformly integrable, being dominated by
  $\lVert Z^\circ_s \rVert ( |Y^\circ_t(X^n)| + |Y^\circ_s(X^n)| )$ and
  $Z^\circ_s$ bounded by Definition~\ref{def:canonical}(i).  So
  Fact~\ref{fact:ui} applies and, with \ref{it:C3c},
  \[
    E^P\bigl[ (Y_t - Y_s) Z^\circ_s(X) \bigr]
      \;=\; E^P\bigl[ \psi(X) \bigr]
      \;=\; \lim_{n \to \infty} E^{P^n}\bigl[ \psi(X^n) \bigr]
      \;=\; 0 ,
  \]
  the first equality because $\XX$ is canonical.  As $Z^\circ_s \in \ZZ^\circ_s$
  was arbitrary and $Y_s, Y_t \in L^1(P)$ by Step~0, part~(ii) of
  Definition~\ref{def:canonical} yields $E^P[Y_t \mid \Filt_s] = Y_s$.  If
  $D = \T$ this is already the assertion.

  \emph{Step 2: uniform integrability under $P$.}
  Let $\varphi_N(x) = |x| - |x| \wedge N$.  For $r \in D \cap \T_{\leq t}$ the map
  $\varphi_N \circ Y^\circ_r$ is $P$-continuous at $X$, and
  $\{\varphi_N(Y^\circ_r(X^n))\}_n$ is uniformly integrable because it is
  dominated by $|Y^\circ_r(X^n)|$.  So Fact~\ref{fact:ui} applies twice and gives
  \[
    E^P\bigl[ \varphi_N(Y_r) \bigr]
      \;=\; \lim_{n \to \infty} E^{P^n}\bigl[ \varphi_N(Y^\circ_r(X^n)) \bigr]
      \;\leq\; \sup_{n \in \N} E^{P^n}\bigl[ \varphi_N(Y^\circ_r(X^n)) \bigr] .
  \]
  Taking the supremum over $r \in D \cap \T_{\leq t}$ and letting $N \to \infty$,
  the right-hand side tends to $0$ by \ref{it:C3b} and the second
  characterization of uniform integrability in Fact~\ref{fact:ui}.  Hence
  $\{ Y_r : r \in D \cap \T_{\leq t} \}$ is uniformly integrable under $P$.

  \emph{Step 3: from $D$ to $\T$.}
  Assume (T2b) and $Y$ right continuous, and let $s < t$ in $\T$ be arbitrary.
  Since $D$ is dense and $\T$ has no isolated points from the right, there are
  $s_m \downarrow s$ and $t_m \downarrow t$ in $D$ with $s_m < t_m$.  Let
  $G \in \Filt_s$.  Then $G \in \Filt_{s_m}$ because $s \leq s_m$, so Step~1 gives
  $E^P[Y_{t_m} \one_G] = E^P[Y_{s_m} \one_G]$.  By right continuity
  $Y_{t_m} \to Y_t$ and $Y_{s_m} \to Y_s$ $P$-a.s., and both families are
  uniformly integrable by Step~2, so Fact~\ref{fact:ui} applied to
  $Y_{t_m}\one_G$ and $Y_{s_m}\one_G$ yields
  $E^P[Y_t \one_G] = E^P[Y_s \one_G]$.  As $G \in \Filt_s$ was arbitrary,
  $E^P[Y_t \mid \Filt_s] = Y_s$, and $P \in \Msol(\XX)$.
\end{proof}

\begin{remark}[Where the index enters]
  \label{rem:absconvindex}
  Steps~0--2 use no structure on $\T$ whatsoever beyond the preorder: $D$ is any
  subset, and $\T_{\leq t}$ is a down-set.  The only place where more is needed is
  Step~3, and it is needed only to pass from $D$ to its complement.  With
  $D = \T$ --- in particular for $\T = \N_0$, where every point is isolated from
  the right and no approximation is available or necessary --- Step~3 is empty and
  the theorem holds under (T0).  This is why step~\ref{it:F5a} of
  Section~\ref{sec:formalization} can be done immediately after \ref{it:F2},
  before any of the Skorokhod theory.
\end{remark}

\begin{remark}[What has to be checked, and what does not]
  \label{rem:absconvcheck}
  Of the three hypotheses, \ref{it:C1} is the input, \ref{it:C2} is a property of
  the path space and is supplied once and for all by
  Example~\ref{ex:determining} or by Proposition~\ref{prop:fddchar}, and
  \ref{it:C3} is the work.  Within \ref{it:C3}, \ref{it:C3a} is where the topology
  of $F$ enters --- and it is the only place; for $F = \DE$ it is the reason for
  the ``countable complement'' hypothesis of Theorem~\ref{thm:CPSconv}, since
  $\alpha \mapsto \alpha(t)$ is $J_1$-continuous exactly at paths not jumping at
  $t$.  Condition \ref{it:C3c} is where the approximating martingales are used,
  and \ref{it:C3b} is what makes the two limits interchangeable.
\end{remark}

\subsection{By convergence: the general version}
\label{ssec:CPSconv}

The following theorem is \CPS, Theorem~4.1; it is the specialization of the
abstract convergence theorem \CPS, Corollary~3.17, to the Markovian martingale
problem, and it generalizes Lemma~\ref{lem:EKconv} in three independent
directions:
\begin{enumerate}[label=(\arabic*)]
\item the approximating objects are \emph{not} required to solve a martingale
  problem for an operator; they are arbitrary pairs $(\xi^n,\varphi^n)$ of
  progressively measurable processes such that
  $\xi^n - \int_0^\cdot \varphi^n_s \dif s$ is a martingale;
\item the approximation \eqref{eq:approxA} in the uniform norm is replaced by
  the much weaker \emph{asymptotic} conditions \eqref{eq:cps2} and
  \eqref{eq:cps3}, which are conditions on the laws only;
\item uniform boundedness is replaced by uniform integrability
  \eqref{eq:cps1}.
\end{enumerate}

\begin{theorem}[\CPS, Theorem 4.1; (T3)]
  \label{thm:CPSconv}
  Assume (E3).  Let $B = (\Omega,\Filt,\mathbb{F},P)$ and
  $B^n = (\Omega^n,\Filt^n,\mathbb{F}^n,P^n)$ be filtered probability spaces
  carrying $E$-valued c\`adl\`ag adapted processes $X$ and $X^n$, and assume that
  $\mathbb{F}$ is generated by $X$.  Let $A \subset \Cb(E) \times \Cb(E)$ and let
  $\XX = \XX_A(X)$ be as in \eqref{eq:XXA}.  For each $n$, let $\XX^n$ be a set
  of pairs $(\xi,\varphi)$ of real-valued progressively measurable processes on
  $B^n$ such that $\sup_{s \leq T} E^{P^n}[|\xi_s| + |\varphi_s|] < \infty$ for
  all $T>0$ and such that $\xi - \int_0^{\cdot} \varphi_s \dif s$ is a martingale
  on $B^n$.

  Suppose that $X^n \wto X$ weakly on $\DE$ with the Skorokhod $J_1$ topology,
  and that there is a set $\Gamma \subset \Rp$ with countable complement such
  that for each $(f,g) \in A$ and $T > 0$ there is a sequence
  $(\xi^n, \varphi^n) \in \XX^n$ with
  \begin{align}
    \label{eq:cps1}
    \sup_{n \in \N} \sup_{s \leq T} E^{P^n}\bigl[ |\xi^n_s| + |\varphi^n_s| \bigr]
      &< \infty , \\
    \label{eq:cps2}
    \lim_{n \to \infty}
      E^{P^n}\Bigl[ \bigl( \xi^n_t - f(X^n_t) \bigr)
        \prod_{i=1}^k h_i(X^n_{t_i}) \Bigr] &= 0 , \\
    \label{eq:cps3}
    \lim_{n \to \infty}
      E^{P^n}\Bigl[ \int_s^t \bigl( \varphi^n_u - g(X^n_u) \bigr) \dif u
        \prod_{i=1}^k h_i(X^n_{t_i}) \Bigr] &= 0 ,
  \end{align}
  for all $k \in \N$, $t_1,\dots,t_k \in \Gamma \cap [0,t]$,
  $t \in \Gamma \cap [0,T]$, $h_1,\dots,h_k \in \Cb(E)$.
  Then $P \in \Msol(\XX)$, i.e.\ $X$ solves the martingale problem for $A$.
\end{theorem}

\begin{proof}
  We verify \ref{it:C1}--\ref{it:C3} of Theorem~\ref{thm:absconv}.

  \emph{The set $D$, and \ref{it:C2}.}
  Set $D = \{ t \in \Gamma : P\{X_t \neq X_{t-}\} = 0 \}$.  By
  Fact~\ref{fact:Dcountable} and the hypothesis on $\Gamma$, $D^c$ is countable,
  so $D$ is dense.  Let $\ZZ^\circ$ be the determining set of
  Example~\ref{ex:determining} built from $D$; this is \ref{it:C2}.
  Hypothesis \ref{it:C1} is the assumed weak convergence $X^n \wto X$ on $\DE$.

  \emph{\ref{it:C3b}.}
  Let $Y = f(X) - \int_0^\cdot g(X_s)\dif s \in \XX$ with canonical version
  $Y^\circ_t(\omega) = f(\omega(t)) - \int_0^t g(\omega(s))\dif s$.  Since $f,g$
  are bounded, $Y^\circ_r$ is bounded on $r \in [0,t]$, so
  $\{Y^\circ_r(X^n) : r \in D \cap [0,t],\, n \in \N\}$ is uniformly integrable.

  \emph{\ref{it:C3a}.}
  The evaluation $\omega \mapsto \omega(t)$ is $J_1$-continuous at every $\omega$
  with $\omega(t) = \omega(t-)$, so the definition of $D$ makes $Y^\circ_t$ and
  $Y^\circ_t Z^\circ_s$ $P$-continuous at $X$ in the sense of
  Definition~\ref{def:Pcont}, for $s,t \in D$.  This is the only point at which
  the Skorokhod topology is used, and it is what the hypothesis on $\Gamma$ is
  for.

  \emph{\ref{it:C3c}.}
  It remains to verify, for $s,t \in D$ with $s \leq t$ and
  $Z^\circ_s = \prod_i h_i(\cdot(t_i))$,
  \begin{equation}
    \label{eq:cpskey}
    \lim_{n \to \infty}
      E^{P^n}\bigl[ \bigl( Y^\circ_t(X^n) - Y^\circ_s(X^n) \bigr) Z^\circ_s(X^n) \bigr]
      \;=\; 0 .
  \end{equation}
  Put $Y^n = \xi^n - \int_0^{\cdot} \varphi^n_u \dif u$, a martingale on $B^n$.
  Because $Z^\circ_s(X^n)$ is $\Filt^n_s$-measurable,
  $E^{P^n}[(Y^n_t - Y^n_s) Z^\circ_s(X^n)] = 0$, so
  \[
    E^{P^n}\bigl[ ( Y^\circ_t(X^n) - Y^\circ_s(X^n) ) Z^\circ_s(X^n) \bigr]
    = E^{P^n}\bigl[ ( Y^\circ_t(X^n) - Y^n_t + Y^n_s - Y^\circ_s(X^n) )
      Z^\circ_s(X^n) \bigr] ,
  \]
  and the right-hand side tends to $0$ by \eqref{eq:cps2} and \eqref{eq:cps3},
  since
  $Y^\circ_t(X^n) - Y^n_t = (f(X^n_t) - \xi^n_t)
    + \int_0^t (\varphi^n_u - g(X^n_u)) \dif u$.
  Thus \eqref{eq:cpskey} holds, which is \ref{it:C3c}.

  Theorem~\ref{thm:absconv} now gives $E^P[Y_t \mid \Filt_s] = Y_s$ for
  $s \leq t$ in $D$, and its Step~3 --- which uses exactly the uniform
  integrability \eqref{eq:cps1} and the right continuity of $Y$ --- extends this
  to arbitrary $s < t$.
\end{proof}

\begin{remark}
  \label{rem:convlocal}
  Theorem~\ref{thm:CPSconv} concludes that $P$ solves the \emph{global}
  martingale problem.  To obtain a solution of a \emph{local} martingale problem
  in the limit, apply the theorem to a localized family of test processes, i.e.\
  to $\XX^{(m)} = \{ Y^{\tau_m} : Y \in \XX \}$ for a suitable exhausting
  sequence of stopping times $\tau_m$, and let $m \to \infty$.  This is how
  \CPS{} obtain their semimartingale results (their Section 4.3), and it is the
  same reduction as in Remark~\ref{rem:uniquelocal}.
\end{remark}

\begin{remark}[Path spaces other than $\DE$]
  \label{rem:CPSabstract}
  Theorem~\ref{thm:absconv} is stated for an arbitrary Polish path space
  $F \subset E^{\T}$, and Theorem~\ref{thm:CPSconv} is the instance $F = \DE$.
  Nothing in the abstract theorem refers to c\`adl\`ag paths, to a Markov
  structure, or to the Skorokhod topology; $F$ may for instance be
  $L^p_{\mathrm{loc}}(\Rp,\R^d)$, which is the state space \CPS{} use for
  solutions of Volterra equations (their Example 3.13), or a product of such
  spaces.  What changes from one $F$ to another is only which maps are
  $P$-continuous at $X$, i.e.\ hypothesis \ref{it:C3a} --- see
  Remark~\ref{rem:absconvcheck}.  \CPS{} also give a version in which the
  uniform integrability of \ref{it:C3b} is dropped (their Corollary 3.21) and one
  with control variables and weak-strong convergence (their Theorem 3.14); the
  latter is what their Section~5.3 needs for processes with fixed times of
  discontinuity.
\end{remark}

% =====================================================================
\section{Notes for the formalization}
\label{sec:formalization}
% =====================================================================

\subsection{What is available and what has to be built}

Section~\ref{sec:notation} lists the prerequisites; the following is a stocktaking
of which of them Mathlib already supplies.

\emph{Available.}  Filtrations, adapted and progressively measurable processes,
stopping times and $\Filt_\tau$ (Definition~\ref{def:process},
Fact~\ref{fact:stoppingtimes}); martingales and the optional stopping theorem
(Definition~\ref{def:martingale}, Fact~\ref{fact:optsampl}); Doob's inequalities
(Fact~\ref{fact:doob}); Polish spaces; weak convergence of measures with the
portmanteau theorem, the L\'evy--Prokhorov metric and Prohorov's theorem
(Facts~\ref{fact:portmanteau}, \ref{fact:PSpolish}, \ref{fact:prohorov});
uniform integrability (Fact~\ref{fact:ui}); monotone class theorems
(Fact~\ref{fact:monotoneclass}).

\emph{To be built.}
\begin{itemize}
\item the Skorokhod space $\DE$ with the $J_1$ metric, and the facts that it is
  Polish, that its Borel field is generated by the coordinates, and the
  characterization of its compact sets (Definition~\ref{def:skorokhod},
  Facts~\ref{fact:DEpolish}, \ref{fact:DEcompact}, \ref{fact:fdd});
\item regularization of submartingales along a dense set
  (Facts~\ref{fact:cadlagext}, \ref{fact:submgreg}) --- but see
  Remark~\ref{rem:bmrepo}: this exists already, outside Mathlib;
\item local martingales in the \EK{} sense
  (Definition~\ref{def:localmartingale}, Fact~\ref{fact:stoppedlocalmg});
\item separating and convergence determining classes and their Stone--Weierstrass
  theory (Definition~\ref{def:separating},
  Facts~\ref{fact:convdet}, \ref{fact:stoneweierstrass}, \ref{fact:fdd},
  \ref{fact:sepcond});
\item bounded pointwise convergence and bp-closures
  (Definition~\ref{def:bp}, Fact~\ref{fact:bp});
\item the continuous mapping theorem in the form of Fact~\ref{fact:cmt} and the
  Skorokhod representation (Fact~\ref{fact:PSpolish});
\item convergence of finite-dimensional distributions in $\DE$ and the relative
  compactness criteria (Facts~\ref{fact:fddconv}, \ref{fact:relcompact},
  \ref{fact:relcompact2});
\item dissipative operators and the full generator of a measurable contraction
  semigroup (Definitions~\ref{def:dissipative}, \ref{def:fullgenerator},
  Fact~\ref{fact:fullgenerator}); these are needed only for
  Proposition~\ref{prop:dissipative} and Remark~\ref{rem:fullgenerator}, which
  are themselves optional context.
\end{itemize}

\begin{remark}[Prior art: the \texttt{brownian-motion} project]
  \label{rem:bmrepo}
  The c\`adl\`ag half of \ref{it:F4} is already formalized, not in Mathlib but in
  the Lean development accompanying \BM, in the repository
  \texttt{RemyDegenne/}\allowbreak\texttt{brownian-motion} on GitHub.  Its file
  \texttt{Cadlag.lean} defines \texttt{IsRightContinuous} for a preorder and
  \texttt{IsCadlag} as a structure, and develops jump sets, the discreteness of
  large jumps, local boundedness and the usual closure properties, with no
  remaining gaps.  Its \texttt{CadlagModification.lean} then constructs the
  c\`adl\`ag modification; the blueprint chapter proves that a submartingale has a
  c\`adl\`ag modification if and only if $t \mapsto E[X_t]$ is right continuous,
  and that a martingale always has one.  These are precisely
  Facts~\ref{fact:cadlagext} and \ref{fact:submgreg}, and they are proved for
  \emph{quasimartingales}, which is more general than what
  Theorem~\ref{thm:absreg} needs.

  Three consequences for the present plan.  First, \ref{it:F4} loses its most
  expensive prerequisite.  Second, the index there is a preorder with an order
  topology, which agrees with the bundle (T2b) of Definition~\ref{def:bundles} ---
  the two developments make the same choice independently.  Third, what is
  \emph{not} there is the Skorokhod space: no $J_1$ metric, no
  $\Bor(\DE) = \sigma(\pi_t)$, no compactness criterion.  So \ref{it:F5b} is
  untouched, and Fact~\ref{fact:DEpolish} and Fact~\ref{fact:DEcompact} remain to
  be built.
\end{remark}

This dictates a natural order of work.  Each item names the bundle of
Definition~\ref{def:bundles} it targets; the point of doing so is that
\ref{it:F1}, \ref{it:F5a} and half of \ref{it:F2} live under
\texttt{[Preorder \(\iota\)]}, which is how Mathlib already indexes
\texttt{Filtration} and \texttt{Martingale}, and can therefore be stated without
any bespoke index theory.

\begin{enumerate}[label=\textbf{(F\arabic*)}]
\item \label{it:F1}
  \emph{Targets (T0) + (E0) + clock.}
  Definitions only: Definition~\ref{def:absMP}, Definition~\ref{def:markovMP},
  the filtration \eqref{eq:starfilt}, Definition~\ref{def:cc} (which needs
  (T2b)), and Proposition~\ref{prop:fddchar}.  The clock is a
  \texttt{MeasureTheory.Measure} on $\T$ together with the hypothesis that every
  down-set is measurable and of finite mass; the interval algebra
  \eqref{eq:clockinterval} is then definitional.  These need nothing beyond what Mathlib has and
  should be done first; note that Proposition~\ref{prop:fddchar} is what turns
  every subsequent statement into a statement about finite-dimensional
  distributions.
\item \label{it:F2}
  \emph{Targets (T0) + (T4) + (E0); (T2a) for the uniqueness half, (E1) for
  Lemma~\ref{lem:disint}.}
  Lemma~\ref{lem:mixture}, Lemma~\ref{lem:restart} and
  Theorem~\ref{thm:absuniq}, then Theorem~\ref{thm:uniqueness} and
  Corollary~\ref{cor:DEuniqueness} as corollaries.  Modulo \eqref{eq:pathsigma}
  this is pure measure theory, and Theorem~\ref{thm:absuniq} needs no path
  regularity at all, so it is the best entry point after \ref{it:F1}.  Almost all
  of the work sits in Lemma~\ref{lem:restart}, which is four lines; the two
  theorems that follow are bookkeeping.  Theorem~\ref{thm:absstrongmarkov}
  additionally needs Fact~\ref{fact:optsampl} and c\`adl\`ag paths and should be
  deferred until after \ref{it:F4}, and so should the local theory of
  Section~\ref{ssec:localmp} --- Definition~\ref{def:localizing} is about
  stopping times.  Note that
  Lemma~\ref{lem:localmix}\ref{it:lm_conv} and Lemma~\ref{lem:L1auto} are the
  only genuinely new proofs there; Lemma~\ref{lem:localrestart} is
  Lemma~\ref{lem:restart} plus the one identity \eqref{eq:shiftedstopped}, and
  Theorem~\ref{thm:localuniq} is bookkeeping.
\item \label{it:F3}
  \emph{Targets (T0) + (T4) + (E0) for Lemma~\ref{lem:chain}, (T2a) + a
  reflective clock with $\iota = \mathrm{p}$ for Proposition~\ref{prop:haar},
  (T3) + (E2) for the rest.}
  Lemma~\ref{lem:chain} --- an exact telescoping identity, and the cheapest item
  in this list --- then Proposition~\ref{prop:haar},
  Corollary~\ref{cor:dualdiscrete}, and finally Lemma~\ref{lem:calculus},
  Theorem~\ref{thm:duality} and their corollaries.  Self-contained: Fubini,
  absolute continuity, dominated convergence.  No Skorokhod space is needed.
  Note that Corollary~\ref{cor:dualdiscrete} needs none of the analysis and can
  be formalized directly after Lemma~\ref{lem:chain}.  Together with \ref{it:F2} this already yields
  Corollary~\ref{cor:uniqviadual}, a complete uniqueness theory.
\item \label{it:F4}
  \emph{Targets (T2b) + (E2), $\iota = \mathrm{o}$.}
  Fact~\ref{fact:cadlagext} and Fact~\ref{fact:submgreg}, then
  Theorem~\ref{thm:absreg} and, as a one-line corollary,
  Theorem~\ref{thm:cadlag}.  This requires the notion of a c\`adl\`ag path but
  \emph{not} the Skorokhod topology.  Formalize the abstract theorem, not the
  Markovian one: Definition~\ref{def:regclass} is three hypotheses on a pair of
  processes, whereas $\dom(A)$ drags in the operator, its domain and the
  filtration \eqref{eq:starfilt}, none of which the proof uses.
\item \label{it:F5a}
  \emph{Targets (T0) + (E3), with $D = \T$.}
  Theorem~\ref{thm:absconv}, the abstract convergence theorem.  It needs an
  abstract Polish path space, a determining set, $P$-continuity at $X$, uniform
  integrability and Vitali's theorem --- and \emph{no} Skorokhod topology, no
  c\`adl\`ag paths and no Markov structure.  It can therefore be done
  immediately after \ref{it:F2}, and it is the cheapest genuine theorem in this
  list rather than, as the ordering suggests, part of the most expensive one.
\item \label{it:F0}
  \emph{Targets (T3) + (E0).}  Out of order, and deliberately: the jump processes
  of Section~\ref{ssec:jump}.  Theorem~\ref{thm:jumpMP} needs a Markov chain,
  i.i.d.\ exponential variables and the computation of Step~2, and nothing else;
  Proposition~\ref{prop:jumpwellposed} adds a Gronwall estimate.  It is by a wide
  margin the cheapest theorem with real content in this manuscript, it is the only
  one that exhibits a solution, and it supplies the missing example for
  Definition~\ref{def:localizing}.  It could be done before \ref{it:F1}, as a way
  of testing the definitions against something concrete.
\item \label{it:F5b}
  \emph{Targets (T3) + (E3).}
  The Skorokhod space and its topology, then Lemma~\ref{lem:EKconv},
  Remark~\ref{rem:EKrelcompact} and Theorem~\ref{thm:CPSconv} as instances of
  \ref{it:F5a}.  This is by far the largest item.
\end{enumerate}

\subsection{Design decisions}

\begin{enumerate}[label=(\alph*)]
\item \emph{Operators as relations.}  $A$ is a subset of
  $\Meas(E) \times \Meas(E)$, not a function.  In Lean this is
  \texttt{A : Set ((E \textrightarrow{} R) \texttimes{} (E \textrightarrow{} R))} together with measurability side
  conditions, or a structure bundling the pair with its measurability.  Keeping
  $A$ multivalued costs nothing and is needed for
  Remark~\ref{rem:fddconsequences}(b).
\item \emph{Two notions of solution.}  Definition~\ref{def:markovMP} defines
  ``solution with respect to $(\Gilt_t)$'' as the primitive and ``solution''
  as the special case $\Gilt_t = {}^{*}\Filt^X_t$.  Doing it the other way round
  duplicates every statement.
\item \emph{Compact containment as a predicate.}  Definition~\ref{def:cc} occurs
  in Theorem~\ref{thm:cadlag}, in Remark~\ref{rem:EKrelcompact} and implicitly in
  Remarks~\ref{rem:uniquelocal} and \ref{rem:convlocal}.  It should be one
  definition with the two variants \eqref{eq:cc} (over $D$, single process) and
  \eqref{eq:ccn} (over $[0,T]$, family) related by a lemma for right continuous
  processes.
\item \emph{Local martingale problems.}  Following
  Section~\ref{ssec:recommendation}, define the local martingale problem at the
  level of Definition~\ref{def:absMP} and derive the localized statements
  (Remarks~\ref{rem:uniquelocal}, \ref{rem:convlocal}) by a single reusable lemma
  of the form: \emph{if $X$ solves the local martingale problem for $A$ with
  localizing sequence $(\tau_m)$, then $X(\cdot \wedge \tau_m)$ solves the
  (global) martingale problem for the stopped operator $A^{(m)}$.}
\item \emph{Separating classes.}  Facts~\ref{fact:sepcond} and \ref{fact:fdd}
  should be developed independently of martingale problems; they are of general
  interest and are the only place where Polishness of $E$ is really used.
\item \emph{No semigroup theory.}  The manuscript rests on Chapters~2 and~3 of
  \EK{} alone, as envisaged in the project plan; of Chapter~1 only the notions of
  a dissipative operator and of the full generator are borrowed, and only for the
  optional Proposition~\ref{prop:dissipative} and Remark~\ref{rem:fullgenerator}.
  See Remark~\ref{rem:noch1}.
\end{enumerate}

% =====================================================================
\begin{thebibliography}{9}

\bibitem{BM25}
R.~Degenne, D.~Ledvinka, E.~Marion and P.~Pfaffelhuber.
\newblock Formalization of Brownian motion in Lean.
\newblock Preprint, arXiv:2511.20118, 2025.

\bibitem{CPS23}
D.~Criens, P.~Pfaffelhuber and T.~Schmidt.
\newblock The martingale problem method revisited.
\newblock \emph{Electron.\ J.\ Probab.} \textbf{28} (2023), article no.~19, 1--46.

\bibitem{EK86}
S.~N.~Ethier and T.~G.~Kurtz.
\newblock \emph{Markov Processes: Characterization and Convergence}.
\newblock Wiley, New York, 1986.

\bibitem{JS03}
J.~Jacod and A.~N.~Shiryaev.
\newblock \emph{Limit Theorems for Stochastic Processes}.
\newblock 2nd edition, Springer, Berlin, 2003.

\bibitem{Kal21}
O.~Kallenberg.
\newblock \emph{Foundations of Modern Probability}.
\newblock 3rd edition, Springer, Cham, 2021.

\bibitem{SV79}
D.~W.~Stroock and S.~R.~S.~Varadhan.
\newblock \emph{Multidimensional Diffusion Processes}.
\newblock Springer, Berlin, 1979.

\end{thebibliography}

\end{document}
